% DER_VII-a: Grothendieck's text with the commentary set in, then the overview, inside a group that applies the chunk's own definitions
\dgrchapter{VII}{Chapter VII}{Catégories de chemins (2)}
\dgrpart{0}{VII-a}{Part 1 of 3, manuscript pages 1--57}{VII-a\enspace\textperiodcentered\enspace Chapter VII, part 1 of 3\enspace\textperiodcentered\enspace manuscript pages 1--57}{Chapter VII \textperiodcentered\ \textup{VII-a}}
\dgrchunkbegin
% DER_VII-a: the definitions of the chunk's preamble, applied inside its group
\let\printedpage\relax\newcommand{\printedpage}[1]{\if\relax\detokenize{#1}\relax\else\marginpar{\raggedright\footnotesize p.~#1}\fi}
\let\msp\relax\newcommand{\msp}[1]{\marginpar{\raggedright\footnotesize [#1]}}
\let\Fltr\relax\newcommand{\Fltr}{\operatorname{Fl}\operatorname{tr}}
\let\Ob\relax\newcommand{\Ob}{\operatorname{Ob}}
\let\bDelta\relax\newcommand{\bDelta}{\boldsymbol{\Delta}}
\let\bD\relax\newcommand{\bD}{\boldsymbol{\Delta}}
\let\iddots\relax\newcommand{\iddots}{\mathinner{\mkern1mu\raise1pt\hbox{.}\mkern2mu\raise4pt\hbox{.}\mkern2mu\raise7pt\hbox{.}\mkern1mu}}
\let\Ch\relax\newcommand{\Ch}{\underline{\mathrm{Ch}}}
\let\uHom\relax\newcommand{\uHom}{\underline{\mathrm{Hom}}}
\let\uFl\relax\newcommand{\uFl}{\underline{\mathrm{Fl}}}
\let\opp\relax\newcommand{\opp}{\mathrm{opp}}
\let\Hom\relax\newcommand{\Hom}{\mathrm{Hom}}
\let\uCh\relax\newcommand{\uCh}{\underline{\mathrm{Ch}}}
\let\Ty\relax\newcommand{\Ty}{\mathrm{Ty}}
\let\uTy\relax\newcommand{\uTy}{\underline{\mathrm{Ty}}}
\let\Cat\relax\newcommand{\Cat}{\mathrm{Cat}}
\let\circled\relax\newcommand{\circled}[1]{\text{\textcircled{$\scriptstyle #1$}}}
\let\ul\relax\newcommand{\ul}[1]{\underline{\mathrm{#1}}}
\let\red\relax\newcommand{\red}{\text{\upshape réd}}
\let\Tr\relax\newcommand{\Tr}{\mathrm{Tr}}
\let\Cr\relax\newcommand{\Cr}{\mathrm{Cr}}
\let\lng\relax\newcommand{\lng}{\operatorname{long}}
\let\card\relax\newcommand{\card}{\operatorname{card}}
\let\uTyp\relax\newcommand{\uTyp}{\underline{\mathrm{Typ}}}
\let\ovalbox\relax\newcommand{\ovalbox}[1]{\tikz[baseline=(X.base)]\node[draw,rounded corners=6pt,inner sep=2pt](X){#1};}

% ---- from Der-ChVII-456_p001-005.tex
% ===================== page 1 =====================
\printedpage{}
% head: (none; title page, no page number printed)

\begin{center}
{\Large\textsf{LES DÉRIVATEURS}}

\bigskip

\textbf{ALEXANDRE GROTHENDIECK}

\bigskip\bigskip

\textbf{Chapitre VII}

\medskip

\textbf{Catégories de chemins (2)}
\end{center}

\newpage

% ===================== page 2 =====================
\printedpage{}
% head: (none; no page number printed)

Ce texte a été déchiffré et transcrit en \LaTeXe\ par M. Künzer. Il a été édité par M. Künzer, J. Malgoire, et G. Maltsiniotis. La transcription est aussi fidèle que possible au manuscrit. Pour les quelques corrections évidentes, ou rares commentaires des éditeurs, ainsi que pour la numérotation originale des pages du manuscrit, les caractères de machine à écrire \texttt{[typewriter]} entre crochets sont utilisés. Un point d'interrogation entre crochets signifie que l'on n'est pas sûr du mot qui précède.

Cette édition est provisoire. Les remarques, commentaires et corrections sont bienvenus. Envoyer un message à :

\textbf{maltsin@math.jussieu.fr}

\hfill G. Maltsiniotis\hspace*{3em}

\newpage

% ===================== page 3 =====================
\printedpage{1}
% head: A. Grothendieck : Dérivateurs, Ch. VII, Édition des Universités de Montpellier II et Paris VII
% foot: Version 22 décembre 2001

\msp{page 1}
Comme le formalisme ‘simplifié’ des chemins que j'ai voulu dégager à la va-vite s'est révélé d'une lourdeur insupportable, de nature à le rendre quasiment inutilisable, je vais essayer de dégager un formalisme souple, collant à l'intuition, au prix seulement d'un effort conceptuel supplémentaire.

\section*{1\quad Intervalles et leur variantes} % [sic] leur variantes -> leurs variantes, plural agreement

On appelle \textit{catégorie intervalle standard} une catégorie définie par un ensemble ordonné, dont l'ensemble des élements est un intervalle d'entiers $[0,n]$ ($n\in\mathbf{N}$, on n'exclut pas le cas $n=0$, i.e. de la catégorie ponctuelle $\{0\}$) : % [sic] élements -> éléments, missing accent
\[
\Ob I=[0,n]\subseteq\mathbf{N},
\]
et satisfaisant l'hypothèse suivante :

\begin{enumerate}
\item[a)] pour toute paire de deux sommets consécutifs $i-1$, $i$ de $I$ ($1\leq i\leq n$), il y a une flèche qui les relie, donc soit $i\longrightarrow i+1$, soit $i\longleftarrow i+1$ \texttt{[plutôt} $i-1\longrightarrow i$\texttt{,} $i-1\longleftarrow i$\texttt{]} (cette flèche est notée $\underline{u}_i$,
\[
(i-1)\xrightarrow{\underline{u}_i}i\qquad\quad\text{ou}\qquad\quad(i-1)\xleftarrow{\underline{u}_i}i\ \ ).
\]
\item[b)] La catégorie $I$ est engendrée par l'ensemble de ses flèches $\underline{u}_i$ (appelées \textit{flèches de transition}).
\end{enumerate}

La flèche de transition est dite \textit{directe} si elle va de $i-1$ à $i$, \textit{rétrograde} dans le sens inverse. Si on affecte l'intervalle $[i-1,i]$ du signe $+$ ou $-$ suivant que $\underline{u}_i$ est directe ou rétrograde, on voit donc que les
\msp{page 2}
intervalles standard de longueur donnée $n$ correspondent biunivoquement aux systèmes de signes $+$ ou $-$ sur les $n$ intervalles de transition $[i-1,i]$ ($1\leq i\leq n$), il y en a donc $2^n$. Le sommet $0$ est l'\textit{origine} ou la \textit{source} de l'intervalle, le sommet $n$ l'\textit{extrémité} ou le \textit{but} de l'intervalle.

\dgrblock{dgrb278}{A combinatorial interval carries two orders}
\begin{dcomment}{exposition}{p1}{A combinatorial interval carries two orders}
Grothendieck begins again because a path should be easy to concatenate and easy to change without losing its endpoints. The extra conceptual effort announced here consists in making its shape an object of study. Before forming a category of all paths, he isolates the finite categories which can serve as their domains. The commentary on Chapter~VI, manuscript p.~69 (``Concatenation and the two iterated units'') discusses the preceding path construction; the present chapter reorganizes its combinatorics.

A standard interval is the free category on an oriented path graph. Write its transition signs as
\[
 \tau=(\varepsilon _1,\ldots,\varepsilon _n),\qquad
 \varepsilon_i\in\{+,-\}.
\]
There is one object at each integer from $0$ to $n$. A positive sign supplies the generating arrow $i-1\to i$ and a negative sign supplies $i\to i-1$. The free category is thin: between two fixed vertices there is at most one directed path, since the underlying undirected graph has no cycle. Thus the description as a category associated with a partially ordered set introduces no additional relations. Conversely, the two conditions in the text reconstruct precisely this free category.

For later use we record the complete hom-set calculation. If $i<j$, then
\[
\begin{split}
 I_\tau(i,j)\ne\varnothing
   &\quad\Longleftrightarrow\quad
   \varepsilon_{i+1}=\cdots=\varepsilon_j=+,\\
 I_\tau(j,i)\ne\varnothing
   &\quad\Longleftrightarrow\quad
   \varepsilon_{i+1}=\cdots=\varepsilon_j=-.
\end{split}
\]
Every nonempty hom-set contains exactly one arrow. These formulas prove both the classification by sign words and the count $2^n$. They also explain why a functor out of an interval will be specified freely by its values on the transition arrows.

The numerical order and the categorical partial order have different jobs. The first tells us in which order the vertices are visited; the second tells us which composites exist. For the shape $0\to1\leftarrow2$, the designated source of the interval is $0$ and its designated target is $2$, although neither statement says that $0$ is initial or $2$ terminal. A path of this shape travels along the first arrow and against the second. Confusing these two orders would already make the later distinction between an opposite path and an opposite category disappear.

Length counts transitions. The interval of length zero has one vertex and no transition, and the interval of length $n$ has $n+1$ vertices. We keep this convention throughout; it will identify the isolated equality $n=\operatorname{card}I$ on ms.~13 as a printed slip rather than a competing convention.
\end{dcomment}


Un intervalle standard est \textit{orienté,} du fait qu'on à la relation d'ordre \textit{total} naturelle sur l'ensemble de ses sommets, savoir celle des entiers dans $[0,n]$. Quand on parlera de la relation d'ordre sur $I$, qu'un sommet $i$ en \textit{‘précède’} ou \textit{‘suit’} un autre $j$ ($i\leq j$ ou $i\geq j$ respectivement), c'est de cette relation d'ordre qu'on voudra parler. Celle associée à la structure de catégorie ne sera jamais regardée comme une relation d'ordre. Il y a également une relation d'ordre total naturelle sur l'ensemble des flèches de transition $\underline{u}_i$ ($1\leq i\leq n$), donnée par la relation d'ordre des entiers $i$, $1\leq i\leq n$, qui les indexent. On dira qu'une flèche $i$ \texttt{[plutôt} $\underline{u}_i$\texttt{]} de transition \textit{précède} ou \textit{suit} une autre \texttt{[}$\underline{u}_j$\texttt{]}, ou qu'elle est \textit{successeur immédiat} ou \textit{prédécesseur immédiat} d'une autre (si $i=j+1$ ou $i=j-1$). % [sic] qu'on à -> qu'on a, verb; [sic] un sommet i en ‘précède’ -> un sommet i ‘précède’, superfluous "en"

% Reader layout: continuous text; edition coordinate retained below.
% ===================== page 4 =====================
\printedpage{2}
% head: A. Grothendieck : Dérivateurs, Ch. VII, Édition des Universités de Montpellier II et Paris VII
% foot: Version 22 décembre 2001

Si on a un ensemble d'entiers
\[
E\subseteq\mathbf{Z},
\]
on considérera dans $E$ la relation d'équivalence engendrée par la relation entre deux éléments $(i,j)$ de $E$ : $j$ est successeur immédiat de $i$ \textit{dans} $\mathbf{Z}$ (i.e. $j=i+1$). Deux éléments distinct $i$, $j$ de $E$ sont congrus pour cette relation si et seulement si l'intervalle d'entiers qu'ils déterminent est $\subseteq E$, donc % [sic] distinct -> distincts, plural agreement
\[
\text{si } i<j,\ i \text{ et } j \text{ sont \textit{‘connectés’} dans } E \text{ si et seulement si } [i,j]\subseteq E.
\]

\msp{page 3}
Quand on parlera des ‘composantes connexes’ de $E$, il s'agira toujours des classes suivant cette relation d'équivalence.

Nous identifierons l'ensemble des flèches de transition $\Fltr(I)$ de $I$ à l'intervalle d'entiers $[1,n]$, donc à une partie de $\mathbf{Z}$. Si on a une partie $E\subseteq\Fltr(I)$, on y a donc la notion de flèches de transition ‘connectées’ dans $E$, et la notion de connexité pour $E$. Si $i<j$, $\underline{u}_i$ et $\underline{u}_j$ sont connectées si et seulement si toutes les flèches de transition intermédiaires entre $\underline{u}_i$ et $\underline{u}_j$ sont dans $E$. Et $E$ est connexe si et seulement s'il est soit vide, soit formée d'une suite de flèches de transition dont chacune est successeur immédiat de celle qui la précède, % [sic] soit formée -> soit formé, agreement with "il"
\[
E=\{\underline{u}_i,\underline{u}_{i+1}\ldots,\underline{u}_{i+l-1}\}, % [sic] \underline{u}_{i+1}\ldots -> \underline{u}_{i+1},\ldots, missing comma
\]
où $l$ est le cardinal de $E$. Une telle $E$ est appelée une \textit{chaîne} de flèches de transition, et $l$ est sa longueur. Une chaîne est dite \textit{composable} si le composé $\underline{u}_{i+l-1}\circ\cdots\circ\underline{u}_{i+1}\circ\underline{u}_i$ \texttt{[ou} $\underline{u}_i\circ\underline{u}_{i+1}\circ\ldots\circ\underline{u}_{i+l-1}$\texttt{]} a un sens, i.e. si elles sont toutes de même sens, soit directe, soit rétrograde.

On désigne par $\Fltr^+(I)$, $\Fltr^-(I)$ l'ensemble des flèches de transition directes, ou rétrogrades, respectivement. Leurs composantes connexes ne sont autres que les chaînes composables maximales. Ainsi, les chaînes composables maximales forment une partition de l'ensemble des flèches de transition, en intervalles d'entiers mutuellement disjoints. Les chaînes forment elles-mêmes un ensemble totalement ordonné
\[
C_1,\ C_2,\ \ldots\ C_r\qquad\quad(1\leq r\leq n).
\]
L'entier $r$ s'appelle la \textit{‘longueur utile’} de l'intervalle
\msp{page 4}
$I$, c'est donc le cardinal de l'ensemble de ses chaînes composables maximales.

Deux chaînes composables maximales \texttt{[consécutives]} sont toujours de signe opposé, si l'une est directe l'autre est rétrograde, et inversement. La première chaîne composable maximale $C_1$ (disons simplement \textit{‘chaîne structurale’}) a le même signe que $\underline{u}_1$; \texttt{[on]} dit que la \textit{direction initiale} de $I$ est directe, resp. rétrograde, s'il en est ainsi de $\underline{u}_1$, ou ce qui revient au même, de $C_1$. De même on définit la \textit{direction finale} de $I$ comme celle de $\underline{u}_n$, ou encore celle de $C_r$. Direction initiale et finale sont égales si la longueur utile $r$ est
\printedpage{3}
% head: A. Grothendieck : Dérivateurs, Ch. VII, Édition des Universités de Montpellier II et Paris VII
% foot: Version 22 décembre 2001
impaire, inégales si elle est paire. (\textbf{NB} Ce n'est que pour $n=0$, i.e. quand $I=\{0\}$ est la catégorie ponctuelle, que direction initiale et direction finale ne sont pas définies.)

Les chaînes composables correspondent aux sous-catégories pleines de $I$ qui sont isomorphes à une catégorie $\bDelta_i$ (où $i$ est la longueur de la chaîne). Notons que $\bDelta_i$ n'a pas d'autre automorphisme que l'identité, donc les chaînes maximales de longueur donnée $i\geq 1$ correspondent exactement aux foncteurs
\[
\bDelta_i\lhook\joinrel\longrightarrow\ I
\]
qui sont injectifs sur les sommets. Un tel foncteur est pleinement fidèle, et un monomorphisme. Ce foncteur définit une application strictement croissante ou strictement décroissante sur
\msp{page 5}
les sommets, suivant que la chaîne est directe ou rétrograde.

\dgrblock{dgrb279}{Runs, simplices, and the meaning of maximality}
\begin{dcomment}{exposition}{p4}{Runs, simplices, and the meaning of maximality}
Grothendieck next separates the places where a path can be composed from the places where its direction turns. A maximal constant-sign run of the word $\tau$ is his structural chain. If the run lengths are $l_1,\ldots,l_r$, then
\[
 n=l_1+\cdots+l_r,
\]
and the sign of the $a$th run is $(-1)^{a-1}\varepsilon_1$, when signs are written as $\pm1$. This proves the parity rule for initial and final directions. The useful length $r$ records the number of changes of compositional regime, rather than the number of arrows which happen to have been written down.

The paragraph identifying chains with simplex subcategories needs a small but substantive separation. A chain was defined on ms.~3 to be a contiguous set of transitions. An object-injective functor $\Delta_k\to I$ may skip vertices. For example,
\[
 \Delta_1\longrightarrow\Delta_2,
 \qquad 0\longmapsto0,\quad1\longmapsto2,
\]
is fully faithful, but its image is not the domain of one transition of $\Delta_2$; nor does it describe its maximal structural chain, whose length is two. The words ``chaînes maximales'' and the unrestricted injectivity condition both occur in the edition. They therefore remain in the transcription, with their correction supplied here.

The exact statement is useful in its own right. For $k>0$, an injective functor from $\Delta_k$ picks $k+1$ vertices in a single structural interval, in their categorical order. Indeed its first and last image vertices are connected by an arrow, so every vertex between them in the numerical order belongs to the same constant-sign run. All intermediate image vertices lie between these two in that run's categorical order. Conversely, any such selection determines an injective functor. It is fully faithful because the target is thin and its order on those vertices is exactly the simplex order. Its object function is strictly increasing in the external orientation for a positive run and strictly decreasing for a negative run.

To recover a chain in the earlier sense, impose numerical contiguity of the image. To recover a maximal structural chain, require the image to be the whole structural interval. These distinctions do not disturb the intended decomposition into runs. They identify which extra condition the printed correspondence omitted, while preserving its correct full-faithfulness assertion.
\end{dcomment}


Notons que tout simplexe type $\bDelta_i$, $i\geq 0$, peut être vu comme un intervalle, de longueur $i$, de longueur utile $1$ (si $i\geq 1$, $0$ si $i=0$). Son origine est $0$, son extrémité est $i$. Mais appelons une catégorie quelconque une \textit{catégorie intervalle} si elle est isomorphe à un intervalle standard. Notons ici

\textbf{Proposition 1.} \textit{Soit $I$ un intervalle standard de longueur $n\geq 1$.}

\begin{enumerate}
\item[\textit{a)}] \textit{Soit $i\in I$. Pour que $i\in\{0,n\}$, i.e. pour que $i$ soit un des deux} ‘bouts’ \textit{(ou} ‘extrémités’\textit{) (source ou but) de $I$, il faut et il suffit que l'une des deux catégories $I/i$ ou $i\backslash I$ soit la catégorie ponctuelle et que l'autre soit isomorphe à un} $\bDelta_j$ (\footnote{où $j$ est la longueur de l'unique chaîne structurale qui ‘touche’ $i$ (i.e. telle que $i$ soit dans la sous-catégorie pleine définie par ladite chaîne).}). \textit{(Caractérisation intrinsèque des bouts d'un intervalle.) Pour que $i$ soit la source, il faut et il suffit que $i\backslash I$} \texttt{[phrase incomplète]}
\item[\textit{b)}] \textit{Tout foncteur $I\longrightarrow I$ qui fixe les deux extrémités de $I$ est l'identité.}
\end{enumerate}

\textbf{Corollaire.} \textit{Le groupe des automorphismes de $I$ a au plus deux éléments. S'il en a deux, l'automorphisme non trivial échange les deux extrémités. (On dit alors que l'intervalle est} symétrique.\textit{)}

\dgrblock{dgrb280}{The endpoint determines an orientation, not an endofunctor}
\begin{dcomment}{exposition}{p5:Le groupe des automorphismes}{The endpoint determines an orientation, not an endofunctor}
The purpose of Proposition~1 is to make orientation intrinsic enough to transport it through isomorphisms. Its first assertion is correct. At an endpoint, one of the two slices is a point and the other is the linearly ordered category running along the unique adjacent structural chain. An interior vertex at which the signs agree has both a nonidentity incoming and a nonidentity outgoing arrow, so neither slice is a point. At a turning vertex, one slice is a point, but the other has two nonempty branches meeting at their common source or common target. In the source case it has two incomparable maximal objects, and in the target case two incomparable minimal objects. It cannot be a simplex category. This gives the stated characterization of the two endpoints.

The second assertion is false with the word ``foncteur'' which the edition prints. On $\Delta_2$ the monotone map
\[
 0\longmapsto0,\qquad1\longmapsto0,\qquad2\longmapsto2
\]
defines a nonidentity endofunctor fixing both endpoints. The arrow $1\to2$ is sent to the composite $0\to2$; there is no requirement that a functor preserve a generating transition as a generating transition. This is exactly the flexibility which the later category of types is intended to accommodate.

What the subsequent argument proves is that an \emph{automorphism} fixing the endpoints is the identity. Here is a direct proof which also supplies the uniqueness statements used on ms.~7--9. A nonidentity arrow of an interval is indecomposable into two nonidentity arrows if and only if it is a transition: the displayed hom-set formula proves this, by factoring every longer directed segment at an intermediate vertex. Isomorphisms preserve indecomposability. The graph of indecomposable arrows, after forgetting directions, is the ordinary path with vertices $0,\ldots,n$. A graph automorphism fixing $0$ fixes its unique neighbour, then that neighbour's still unvisited neighbour, and so on. It fixes every vertex and hence every arrow. Every graph automorphism is consequently the identity or vertex reversal. The categorical automorphism group is a subgroup of this two-element group.

This is also the intrinsic reconstruction of the external total order: choose one endpoint, call it $s_0$, and successively take the unique neighbour not already visited. Grothendieck's slice argument implements the same reconstruction in categorical terms. In its inductive paragraph the slices must be taken in the remaining full tail $I'$ at $s_i$; the printed $i\backslash I$ and $I/i$ have dropped that qualification. Neither reconstruction is preserved by arbitrary functors, which may collapse or lengthen transitions.

The unfinished sentence about characterizing the \emph{source} of the interval should not be completed by guessing. The underlying category determines the pair of endpoints, but a symmetric interval has an automorphism interchanging them. Selecting one of them is precisely the additional orientation. After the definition of a reduced interval on ms.~12 we prove a sharp version of the stronger rigidity assertion: all endpoint-fixing endofunctors are identities exactly for reduced intervals.
\end{dcomment}


\msp{page 6}
NB Si $i$ n'est pas un bout, et si néanmoins l'une des deux catégories $i\backslash I$, $I/i$ est ponctuelle, disons $I/i$, alors au voisinage immédiat de $i$, $I$ a la forme
\[
\begin{tikzcd}[row sep=0pt]
\bullet & \bullet \arrow[l, "\underline{u}_i"'] \arrow[r, "\underline{u}_{i+1}"] & \bullet \\
i-1 & i & i+1
\end{tikzcd}
\]
($\underline{u}_i$ rétrograde, $\underline{u}_{i+1}$ directe - si c'était $i\backslash I$ la ponctuelle, alors $\underline{u}_i$ serait directe, $\underline{u}_{i+1}$ rétrograde - en tous les cas ces deux flèches sont de sens opposé, i.e. non composables). On trouve donc que $i\backslash I$ est une catégorie de la forme

% ---- from Der-ChVII-456_p006-010.tex
\printedpage{4}
% head: A. Grothendieck : Dérivateurs, Ch. VII, Édition des Universités de Montpellier II et Paris VII
% foot: Version 22 décembre 2001

\[
\begin{tikzcd}[column sep=small, row sep=small]
 & & & i \arrow[dl] \arrow[dr] & & & \\
 & & i-1 & & i+1 & & \\
 & \iddots \arrow[dl] & & & & \ddots \arrow[dr] & \\
 i-r & & & & & & i+r, % [sic] i+r -> i+s, s est la longueur de la chaîne contenant u_{i+1}
\end{tikzcd}
\]

où $r$ est la longueur de la chaîne structurale contenant $\underline{u}_i$, $s$ celle de la chaîne structurale contenant $\underline{u}_{i+1}$. Cette catégorie est une somme amalgamée $\bD_r \amalg \bD_s$ (recollées suivant leur origine), mais n'est pas isomorphe à une $\bD_j$, p.~ex. faute d'avoir un objet final.

Cette observation prouve a). Pour prouver b), il suffit de prouver que l'ordre total sur $I$ est déterminé de façon intrinsèque, par le choix de \textit{l'origine} $s_0$ de $I$ (parmi ces % [sic] ces -> ses, possessif
deux extrémités). Supposons que l'on ait déjà déterminé intrinsèquement les sommets $s_0$, $s_1$, $\ldots$, $s_i$ ($i \geq 0$). Comment déterminer $s_{i+1}$? (Quand $i < n$.) Soit $I'$ la sous-catégorie pleine de $I$
\msp{page 7}
formée des $j \in I$ qui ne sont pas parmi les $s_j$ \texttt{[plutôt} $s_k$\texttt{]} déjà construits, sauf le dernier $s_i$ (donc $s_i \in I'$, $\operatorname{card} I' \geq 2$). Alors $I'$ est isomorphe à un intervalle standard, avec $s_i$ correspondant à l'origine. Une des catégories $i\backslash I$, $I/i$ est ponctuelle, l'autre isomorphe à une catégorie $\bD_j$, et de plus, $i$ est objet initial ou final de $\bD_j$. Si c'est l'objet initial $0$, on prend pour $s_{i+1}$ le successeur immédiat de $i$ dans $\bD_j$. Si c'est l'objet final $j$, on prend pour $s_{i+1}$ le prédécesseur immédiat $j-1$.

Si on était parti du but $n$ de $I$, au lieu de partir de la source, on aurait trouvé la relation d'ordre total opposée. Affectant l'élément but du numéro d'ordre $0$, on trouve une renumération de $I$, $i$ devenant $n-i$, et les flèches de transition qui étaient directes pour la relation d'ordre initiale, deviennent rétrogrades dans la relation d'ordre opposée, et inversement. On voit ainsi :

\textbf{Proposition 2.} \textit{Si $I$ est un intervalle standard, il existe un unique couple $(I',\theta)$, où $I'$ est un intervalle standard, et $\theta$ un isomorphisme $I \xrightarrow{\sim} I'$ qui transforme le but de $I$ dans la source de $I'$ (et la source de $I$ en le but de $I'$).}

Cet intervalle $I'$ est appelé l'intervalle \textit{opposé} de $I$. Il faut le voir, intrinsèquement, comme la `même' categorie $I$, % [sic] categorie -> catégorie, accent manquant
mais avec une `orientation' opposée.

\msp{page 8}

Une façon plus claire de voir ceci est la suivante. Soit $J$ une catégorie intervalle, i.e. isomorphe à un intervalle standard. Alors il y a exactement (si $I \not\simeq e$) \textit{deux} couples $(I,\theta)$ d'un intervalle standard $I$ et d'un isomorphisme $\theta : I \xrightarrow{\sim} J$. Ces deux isomorphismes induisent sur $J$ deux ordres totaux, opposés l'un de l'autre. Chaque couple est déterminé de façon unique par ledit ordre total, qu'on appelera % [sic] appelera -> appellera, orthographe
une \textit{orientation} de l'intervalle $J$. Ainsi $J$ a deux orientations, et chacune détermine un \textit{isomorphisme orienté} (i.e. respectant
\printedpage{5}
% head: A. Grothendieck : Dérivateurs, Ch. VII, Édition des Universités de Montpellier II et Paris VII
% foot: Version 22 décembre 2001
les orientations) de $J$ avec un intervalle standard (\footnote[2]{\textbf{NB} On parlera de l'intervalle \textit{orienté opposé} $I^{\mathrm{opp}}$ d'un intervalle \textit{orienté} $I$, qu'il faut se garder de confondre avec la catégorie opposée $I^o$. Ici $I^{\mathrm{opp}}$ a même catégorie sous-jacente que $I$, mais avec une orientation opposée.}). Ces deux intervalles sont les mêmes si et seulement s'il existe un automorphisme de $J$ (nécessairement unique) qui échange ses deux orientations, ou ce qui revient au même, qui échange entre elles les deux extrémités. (Ainsi les catégories opposées l'une de l'autre
\[
\begin{tikzcd}[column sep=small, row sep=small]
 & \bullet \arrow[dl] \arrow[dr] & & & & \bullet & \\
 \bullet & & \bullet & & \bullet \arrow[ur] & & \bullet \arrow[ul]
\end{tikzcd}
\]
sont chacunes % [sic] chacunes -> chacune, orthographe
des intervalles symétriques.) Choisir une orientation revient aussi `a % [sic] `a -> à, accent imprimé comme guillemet
choisir l'origine parmi les deux extrémités.

Notons que la \textit{catégorie opposée} à un intervalle est un intervalle, et la catégorie opposée à un intervalle standard est l'intervalle standard qui s'en déduit en changeant tous les signes sur les intervalles de transition $[i-1,i]$ ($1 \leq i \leq n$). On se gardera de confondre la propriété $I \simeq I^o$, avec la propriété de symétrie.

\msp{page 9}

Ainsi les deux catégories ($\Psi_1$ et $\Phi_1$) ci-dessus sont symétriques, mais opposées l'une de l'autre sans être isomorphes. D'autre part, les $\bD_i$ ne sont pas symétriques, mais chacune isomorphe à la catégorie opposée.

Si $I$ est isomorphe à la catégorie opposée $I^o$, et si $\theta$ est un tel isomorphisme (nécessairement unique), alors il est immédiat que $\theta$ renverse l'orientation (sinon ce serait l'identité sur les sommets, et les flèches de transition seraient à la fois directes et indirectes), donc $\theta(0) = n$, $\theta(1) = n-1$, et on veut que $\underline{u}_1$ et $\underline{u}_n$ sont de même signe. Mais si $I$ est symétrique, $\underline{u}_1$ et $\underline{u}_n$ sont de signe opposé. Donc $I$ ne peut être à la fois symétrique et isomorphe à l'intervalle opposé \texttt{[plutôt à la catégorie opposée; exception :} $I = \{0\}$\texttt{]}.

Pour toute catégorie $\bD_i$, celle-ci étant un intervalle, il y a (si $i > 0$) \textit{deux orientations.} En tant qu'intervalles orientés, il y a $\bD_i^+$ (avec l'orientation ordinaire, donné par l'ordre ordinaire sur $\operatorname{Ob} \bD_i = [0,i]$), et $\bD_i^-$ avec l'orientation opposée.

\printedpage{6}
% head: A. Grothendieck : Dérivateurs, Ch. VII, Édition des Universités de Montpellier II et Paris VII
% foot: Version 22 décembre 2001

\msp{page 10}

\dgrblock{dgrb281}{Reversal and categorical duality in a sign word}
\begin{dcomment}{examples}{p9}{Reversal and categorical duality in a sign word}
Grothendieck's two uses of ``opposite'' describe operations which can be written down independently. This makes the examples of symmetry and self-duality entirely explicit. For a sign word $\tau=(\varepsilon_1,\ldots,\varepsilon_n)$, put
\[
\begin{aligned}
 R\tau&=(-\varepsilon_n,\ldots,-\varepsilon_1),\\
 D\tau&=(-\varepsilon_1,\ldots,-\varepsilon_n).
\end{aligned}
\]
The first is the standard model of the same category with its external orientation reversed; the standardizing isomorphism sends $i$ to $n-i$. The second is the categorical opposite with the original numerical orientation. The two involutions commute, and $RD\tau$ is the word read backwards without changing its signs.

For $n>0$, the interval is symmetric precisely when
\[
 \varepsilon_i=-\varepsilon_{n+1-i}\qquad(1\leq i\leq n).
\]
Necessity follows because a nonidentity automorphism must be reversal, by the preceding proof; sufficiency follows by checking each transition. This condition forces $n$ to be even. Likewise $I\cong I^o$ precisely when
\[
 \varepsilon_i=\varepsilon_{n+1-i}\qquad(1\leq i\leq n).
\]
The isomorphism must reverse the numerical order: an order-preserving bijection of the vertex line would be the identity and would have to reverse a nonidentity transition without changing its endpoints. These two sign conditions cannot hold together for a nonempty word.

For instance, $(+,-)$ and $(-,+)$ satisfy the first condition. They give the cospan and span, each symmetric, but the two categories are not isomorphic: one has a terminal object and two minimal objects, the other an initial object and two maximal objects. The constant word $(+,+,+)$ satisfies the second condition and gives the self-dual category $\Delta_3$ with no nonidentity automorphism. A less uniform self-dual example is $(+,-,+)$; its middle sink and source are exchanged under the anti-isomorphism. For a reduced alternating word, symmetry holds at positive even length and self-duality at odd length.

The point is separate: it has a unique orientation and only the identity automorphism, and both operations leave it unchanged. Keeping this case explicit avoids construing its coincident endpoints as two different choices.
\end{dcomment}


\section*{2\quad Sous-catégories pleines d'un intervalle}

Soit $I$ une catégorie intervalle orientée, de longueur $n \geq 1$. Nous identifierons $\operatorname{Ob} I$ à l'ensemble des entiers $[0,n]$. Pour une chaîne $\underline{u}_i$, $\underline{u}_{i+1}$, $\ldots$, $\underline{u}_{i+l-1}$ de flèches de transition dans $I$, on appelle \textit{domaine} de la chaîne, l'ensemble des sources et de buts de ces éléments. C'est donc un intervalle d'entiers $[i-1, i+l-1]$, de longueur $l$ égale à celle de la chaîne, dont l'origine est déterminé % [sic] déterminé -> déterminée, accord avec l'origine
par la première flèche $\underline{u}_i$ de la chaîne (c'est sa source, si $\underline{u}_i$ est directe, et son but si $\underline{u}_i$ est rétrograde), et dont l'extrémité est déterminée par la dernière $\underline{u}_{i+l-1}$. Les intervalles d'entiers (vus comme des sommets de $I$) correspondant aux chaînes structurales sont appelés les \textit{intervalles structuraux} ou \textit{intervalles critiques} de $I$. Contrairement aux chaînes structurales, ces intervalles ne forment pas une partition de $I$, puisque deux intervalles structuraux peuvent avoir une extrémité commune. Si les chaînes structurales sont
\[
C_1,\ C_2,\ \ldots\ ,\ C_r \qquad (r = \text{ longueur utile de } I)
\]
de longueurs
\[
l_1,\ l_2,\ \ldots\ ,\ l_r,
\]
alors les intervalles structuraux sont
\[
\underbrace{[0,l_1]}_{I_1},\ \underbrace{[l_1,l_1+l_2]}_{I_2},\ \ldots\ ,\ \underbrace{[l_1+\cdots+l_{r-1},l_1+\cdots+l_{r-1}+l_r]}_{I_r},
\]
une suite d'intervalles d'entiers dont deux consécutifs ont une extrémité commune, et tels que si $i < j$, chaque élément de $I$ % [sic] I -> I_i, comparaison des éléments de I_i et de I_j
soit $\leq$ chaque élément de $I_j$. On appelle \textit{sommet critique} de $I$ les extrémités des intervalles structuraux.

\textbf{Proposition 3.} \textit{Soit $s \in I$, soit $I' = I/s$ (sous-catégorie pleine de $I$ des $s'$ tels qu'il existe $s' \longrightarrow s$), $I'' = I/s$ \textup{\texttt{[plutôt} $I'' = s\backslash I$\texttt{]}}, $I'_0 = (I/s)\backslash\{s\}$, $I''_0 = (s\backslash I)\backslash\{s\}$} (\footnote[3]{i.e. qu'on ait : il n'y a pas de flèches entre un élément de $I'_0$ et un élément de $I''_0$.}). Pour que $s$ soit critique, il faut et il suffit que $I\backslash\{s\}$ soit catégorie somme de $I'_0$ et $I''_0$.

\msp{page 11}

\dgrblock{dgrb282}{Critical vertices and the correct deletion test}
\begin{dcomment}{exposition}{p10}{Critical vertices and the correct deletion test}
Grothendieck wants to retain exactly the vertices which a path cannot bypass by composing arrows. These are the endpoints of the structural intervals. If their positions are
\[
 0=c_0<c_1<\cdots<c_r=n,
\]
then $[c_{a-1},c_a]$ is one constant-sign run, and each interior $c_a$ is a change of sign. In the categorical partial order, the critical vertices are exactly the minimal and maximal vertices. This observation explains the two slices in Proposition~3, but it also exposes an error in its global formulation.

Let
\[
 A=(I/s)\setminus\{s\},\qquad
 B=(s\backslash I)\setminus\{s\}.
\]
These are respectively the strict predecessors and strict successors of $s$ in the categorical order. They do not in general exhaust the other vertices. In $0\to1\leftarrow2$, take $s=0$. It is critical; $A$ is empty and $B$ consists of $1$, whereas $I\setminus\{0\}$ still contains $2$. Thus the printed equality with all of $I\setminus\{s\}$ cannot be right. The edition prints the same reading.

The footnote gives the valid local test: there is no arrow between $A$ and $B$. If $s$ is critical, it is minimal or maximal and one of these categories is empty. If $s$ is not critical, its two adjacent transitions have the same sign. One adjacent vertex is a strict predecessor and the other a strict successor; their composite through $s$ is an arrow between $A$ and $B$. Equivalently, the full subcategory on the vertices comparable with $s$, after removing $s$, is the coproduct $A\amalg B$ exactly when $s$ is critical.

For the application to subintervals on the next page, the appropriate \emph{global} test uses the numerical order. Suppose $0<s<n$, and let $L$ and $H$ be the full subcategories on $[0,s-1]$ and $[s+1,n]$. Then
\[
 I\setminus\{s\}=L\amalg H
 \quad\Longleftrightarrow\quad s\text{ is critical}.
\]
If $s$ is a turning vertex, an arrow from one side to the other would traverse two transitions of opposite signs, which is impossible. If it is not a turning vertex, the two neighbouring vertices are connected by their composite through $s$, an arrow still present in the full subcategory obtained by deleting $s$. This corrected statement proves the separation property used in Proposition~5. At a numerical endpoint, deleting it leaves a connected interval; its criticality is already accounted for by its being an endpoint.
\end{dcomment}


\textbf{Proposition 4.} \textit{Soient $i,j$ dans $I$, avec $i < j$. Conditions équivalentes :}
\begin{enumerate}
\item[\textit{a)}] \textit{Il existe une flèche entre $i$ et $j$ (soit $i \longrightarrow j$, soit $j \longrightarrow i$).}
\item[\textit{b)}] \textit{$i$ et $j$ appartiennent à un même intervalle structural (nécessairement unique).}
\item[\textit{c)}] \textit{La sous-catégorie pleine de $I$ définie par l'ensemble d'objets $[i,j]$ est isomorphe à une $\bD_\alpha$ (on aura $\alpha = j-i$).}
\end{enumerate}

\printedpage{7}
% head: A. Grothendieck : Dérivateurs, Ch. VII, Édition des Universités de Montpellier II et Paris VII
% foot: Version 22 décembre 2001
On dira alors que $i$ et $j$ sont \textit{directement connectés} dans $I$. (Et on utilisera la même terminologie dans une catégorie quelconque : $i$ et $j$ sont directement connectés si et seulement si $\operatorname{Hom}_I(i,j)$ \textit{ou} $\operatorname{Hom}_I(j,i)$ est non vide.)

\textbf{Proposition 5.} \textit{Soit $J \subseteq I$ une partie de $I$, considérons la sous-catégorie pleine de $I$ qu'elle définit. Alors :}
\begin{enumerate}
\item[\textit{a)}] \textit{Les composantes connexes $J_\alpha$ de $J$ sont des intervalles et l'ordre induit sur un $J_\alpha$ par celui de $I$ est une orientation (dite \textup{orientation induite}) de $J_\alpha$.}
\item[\textit{b)}] \textit{Supposons $\operatorname{card} J > 1$. Soit $i$ le plus petit élément de $J$, $j$ le plus grand élément, donc $i < j$. Pour que $J$ lui-même soit un intervalle, il faut et il suffit que tout sommet critique dans $[i,j]$ soit élément de $J$.}
\end{enumerate}
\nopagebreak[4]
(\footnote[4]{On appelle \textit{sous-intervalle} de $I$ une sous-categorie % [sic] sous-categorie -> sous-catégorie, accent manquant
pleine connexe, i.e. qui est elle-même un intervalle. Notons d'ailleurs que \textit{toute} sous-catégorie $I'$ de $I$ qui est un intervalle, est nécessairement pleine, cf. plus bas \texttt{[paragraphe]} des foncteurs entre intervalles.}).

\textbf{NB} Si $s \in [i,j]$ est un sommet critique de $I$ qui n'est pas dans $J$, alors par la proposition 3 il disconnecterait $J$, i.e. $J$ serait somme directe de deux sous-catégories $J'$ et $J''$ non vides ($i \in J'$, $j \in J''$), donc $J$ ne serait pas un intervalle.

\msp{page 12}

\dgrblock{dgrb283}{Connected full subcategories and the vertices they must retain}
\begin{dcomment}{examples}{p11}{Connected full subcategories and the vertices they must retain}
Grothendieck is preparing a reduction which composes consecutive arrows without identifying a change of direction. The distinction between a full subcategory and an arbitrary subcategory matters here. A full subcategory retains every composite whose endpoints remain, even when an intermediate vertex has been removed. Thus numerical gaps do not necessarily disconnect it.

Let $J=\{j_0<\cdots<j_m\}$ be a nonempty subset of the vertices of $I$, with its full induced category. Two successive chosen vertices $j_a,j_{a+1}$ are joined by an arrow if and only if the signs between them are constant. By Proposition~4 this is equivalent to there being no critical vertex strictly between them. Whenever this condition fails, there can be no arrow from any chosen vertex at or before $j_a$ to any chosen vertex at or after $j_{a+1}$, in either direction: such an arrow would cross the same change of sign. The connected components are therefore the maximal consecutive blocks of the list $(j_0,\ldots,j_m)$ in which successive entries are directly connected.

Within one such block, make a generating transition between each successive pair, in the direction of its arrow in $I$. The free category on this oriented line is exactly the induced category. To verify fullness, an arrow between nonconsecutive retained vertices exists in $I$ precisely when the signs on the whole intervening segment agree. That is also precisely the condition that the successive arrows of the block compose. This proves Proposition~5(a), including the claim about the induced orientation. It proves (b) as well: the whole chosen list is connected exactly when every ambient critical vertex between its extreme vertices was retained.

For a concrete calculation take
\[
 I:\quad0\longrightarrow1\longrightarrow2
       \longleftarrow3\longleftarrow4\longrightarrow5.
\]
Its critical vertices are $0,2,4,5$. The full category on $\{0,2,4,5\}$ is the reduced interval
\[
 0\longrightarrow2\longleftarrow4\longrightarrow5.
\]
Deleting $1$ or $3$ has merely replaced a pair of arrows by its composite. The full category on $\{0,1,4,5\}$ has two components, $\{0,1\}$ and $\{4,5\}$, because the omitted turning vertex $2$ separates them. The least subinterval containing $1$ and $5$ is $\{1,2,4,5\}$: the endpoints must remain, the two intervening turns must remain, and the preceding calculation shows that nothing else is required. This is the corollary on ms.~12.

The last sentence of footnote~4 asserts more than fullness permits: an interval subcategory need not be full. Inside $\Delta_2$, keep all three objects and the arrows $0\to1$ and $0\to2$, but omit $1\to2$. This is a subcategory, because its two nonidentity arrows are not composable. Abstractly it is the span interval with vertex order $1,0,2$. It is not full in $\Delta_2$. Thus fullness must remain part of the definition of a subinterval; it cannot be inferred solely from the abstract shape of the subcategory.

Finally, for a full subinterval with at least two vertices, an interior retained vertex is a turn in the induced shape exactly when it was an ambient critical vertex. This proves Proposition~6. The reduced interval is consequently the smallest full connected part of $I$ which contains its endpoints, and its length is the useful length of $I$.
\end{dcomment}


\textbf{Corollaire.} \textit{L'ensemble des sommets critiques de $I$ est aussi l'ensemble des sommets du plus petit sous-intervalle de $I$ contenant les extrémités. Plus généralement, si $i < j$ sont deux sommets de $I$, le plus petit sous-intervalle de $I$ contenant $i$, $j$ est formé de $i$, $j$ et des sommets critiques $s$ tels que $i < s < j$.}

\textbf{Définition.} Un intervalle est dit \textit{réduit} (ou \textit{irrédondant}) si tous ses sommets sont critiques, ou ce qui revient au même, si ses chaînes structurales sont de longueur 1.

Un intervalle orienté \textit{réduit} est connu, à isomorphisme près, quand on connaît
\begin{enumerate}
\item[a)] sa longueur $\mathrm{n}$,
\item[b)] sa direction initiale $+$ (directe) ou $-$ (rétrograde).
\end{enumerate}

\textbf{NB} Cette direction est la même que la direction finale si $n$ est impair, elle est opposée si $n$ est pair.

Ainsi, suivant les deux cas de direction initiale, on trouve
\[
\begin{aligned}
0 \longrightarrow 1 \longleftarrow 2 \longrightarrow \cdots \qquad\qquad \cdots\ n & \qquad (\text{Avec } & n-1 \longrightarrow n & \quad \text{si } n \text{ impair,}\\
& & n-1 \longleftarrow n & \quad \text{si } n \text{ pair.)}\\
0 \longleftarrow 1 \longrightarrow 2 \longleftarrow \cdots \qquad\qquad \cdots\ n & \qquad (\text{Avec } & n-1 \longleftarrow n & \quad \text{si } n \text{ impair,}\\
& & n-1 \longrightarrow n & \quad \text{si } n \text{ pair.)}
\end{aligned}
\]

\printedpage{8}
% head: A. Grothendieck : Dérivateurs, Ch. VII, Édition des Universités de Montpellier II et Paris VII
% foot: Version 22 décembre 2001

\textbf{Proposition 6.} \textit{Un sous-intervalle $J$, ayant $i$ et $j$ ($i < j$) comme plus petit et plus grand sommet, est réduit si et seulement si tout $s \in J$ distinct de $i$ et de $j$ est critique.}

\textbf{Corollaire.} \textit{La sous-catégorie pleine de $I$ formée des sommets critiques est un sous-intervalle. On l'appelle l'intervalle réduit associé à $I$, noté $I_{\textup{réd}}$.}

% ---- from Der-ChVII-456_p011-015.tex
\printedpage{9}
% head: A. Grothendieck : Dérivateurs, Ch. VII, Édition des Universités de Montpellier II et Paris VII
% foot: Version 22 décembre 2001

\msp{page 13}

\dgrblock{dgrb284}{The critical-vertex inclusion has two adjoints}
\begin{dcomment}{development}{p12}{The critical-vertex inclusion has two adjoints}
The reduced interval retains every turn, but Grothendieck has not yet described how to return to it from an arbitrary vertex. There are two particularly useful choices, one following the categorical source of a run and one following its categorical target. They relate reduction to the endpoint-preserving transformations which a path calculus needs.

\textbf{Two canonical collapses.}\ours{}
Let $j:R=I_{\mathrm{red}}\hookrightarrow I$ be the full inclusion of the critical vertices of an oriented interval. There are retractions $r_s,r_t:I\to R$ with
\[
 r_t\dashv j\dashv r_s,\qquad
 r_sj=r_tj=\mathrm{id}_R.
\]
Both retractions fix the endpoints. At a noncritical vertex $x$, $r_s(x)$ is the categorical source of its structural interval and $r_t(x)$ is that interval's categorical target. At a critical vertex both are the identity. For the point we set $R=I$ and take all three functors to be the identity.

\emph{Proof.} A noncritical vertex lies strictly inside a unique structural interval. On a positive run $[a,b]$, the source-collapse sends every interior vertex to $a$, leaves $b$ at $b$, and is a monotone map of its categorical chain. On a negative run the source is $b$ and the same argument applies with the categorical order reversed. The target-collapse is obtained by exchanging source and target. The run-wise maps agree at shared critical vertices. Every arrow belongs to one run, so these formulas give functors on the whole interval.

To verify the adjunctions rather than merely exhibit homotopies, let $q$ be a critical vertex. If $x$ is itself critical, fullness of $j$ gives the required hom-set bijections. If $x$ is interior to a run, the only critical vertices comparable with $x$ are its source $a$ and target $b$, with $a\to x\to b$ in the categorical order. Indeed an arrow extending beyond either endpoint of that run would have to pass a change of sign. Consequently
\[
\begin{split}
 I(jq,x)\ne\varnothing&\ \Longleftrightarrow\ q=a
     \ \Longleftrightarrow\ R(q,r_sx)\ne\varnothing,\\
 I(x,jq)\ne\varnothing&\ \Longleftrightarrow\ q=b
     \ \Longleftrightarrow\ R(r_tx,q)\ne\varnothing.
\end{split}
\]
For the middle equivalences we use that a source endpoint is a minimal vertex of $I$ and a target endpoint a maximal vertex; hence the same holds in $R$. All these hom-sets are empty or singletons, so the bijections are unique and natural. They establish $j\dashv r_s$ and $r_t\dashv j$. The two retraction identities follow directly from the definition. The counit and unit are the unique arrows
\[
 j r_s(x)\longrightarrow x,\qquad x\longrightarrow j r_t(x),
\]
which are identities whenever $x$ is critical. This also verifies the triangle identities directly.\hfill$\square$

These natural transformations are categorical deformations relative to all critical vertices, in particular relative to both endpoints. Applying the nerve, a natural transformation $F\Rightarrow G$ supplies a simplicial homotopy through the functor $I\times\Delta_1\to I$ which it defines. Thus $Nj$ has homotopy inverses $Nr_s$ and $Nr_t$. Nothing in this argument asserts that the interval categories themselves are equivalent: a skeletal interval with an interior noncritical vertex has strictly more objects than its reduction.

The effect on diagrams gives a more informative comparison. For every category $X$, precomposition yields
\[
 r_s^*\dashv j^*\dashv r_t^*,\qquad
 j^*r_s^*=j^*r_t^*=\mathrm{id}_{\operatorname{Fun}(R,X)}.
\]
This reversal of adjunctions under precomposition is the standard calculation of \citet[Proposition~4.3.7]{Riehl2016}; it can also be read directly from the two displayed natural transformations by applying each diagram $I\to X$. The source extension inserts identity arrows before the final arrow of each directed run, and the target extension inserts them after its first arrow. Restriction to $R$ composes each run.

For $I=\Delta_2$, $R=\{0,2\}$ and $X=\Delta_1$, consider the two diagrams with vertex values
\[
 F=(0,0,1),\qquad G=(0,1,1).
\]
They have the same restriction $(0,1)$ to $R$. There is no natural transformation $G\to F$, since its middle component would have to be $1\to0$ in $\Delta_1$, whereas their restrictions have an identity map. Hence $j^*$ is not full and is not an equivalence. If $X$ is a groupoid, every component of both comparison transformations is invertible, and $j^*$ is an equivalence, with either section as an inverse. Reduction therefore loses factorization data in general, but loses no such data when all arrows of the target are already invertible.

The general precomposition-adjunction statement is known and cited above. The critical-vertex reduction itself also appears in \citet[Lemma~4.16]{ArakawaCnossen2026}. For the additional explicit adjoints and the sharp rigidity criterion below, we checked that construction, the local interval arguments here, and the adjunction calculus in \citet[\S4.3]{Riehl2016}; targeted searches for critical-vertex reductions and their adjoints did not locate these particular formulations. We supply the extra deductions as elementary developments of the present definitions, without a priority claim.
\end{dcomment}


\dgrblock{dgrb285}{The exact scope of endpoint rigidity}
\begin{dcomment}{development}{p12}{The exact scope of endpoint rigidity}
The false generality of Proposition~1(b) suggests asking precisely which part of an interval allows an endpoint-preserving functor to move a vertex. The preceding collapses give one implication, while a finite displacement argument gives the other.

\textbf{Endpoint rigidity criterion.}\ours{}
For a finite oriented interval $I$, the following are equivalent: $I$ is reduced; every endofunctor of $I$ fixing its two designated endpoints is the identity. We include the point as a reduced interval.

\emph{Proof.} Suppose $I$ is reduced, of length $n>0$, and identify its vertices with $[0,n]$. A nonidentity arrow of $I$ is then a single transition: two consecutive transitions have opposite signs and cannot compose. If $f:I\to I$ is any functor, the image of a transition is therefore either an identity or a transition. Its object function satisfies
\[
 -1\leq f(i)-f(i-1)\leq1\qquad(1\leq i\leq n).
\]
If $f(0)=0$ and $f(n)=n$, summation gives
\[
 n=f(n)-f(0)=\sum_{i=1}^n\bigl(f(i)-f(i-1)\bigr)\leq n.
\]
Equality forces every summand to be $1$. Thus $f(i)=i$ for all $i$, and thinness forces $f$ to be the identity on arrows as well. We have made no monotonicity assumption on $f$; it is a consequence of the equality. The point case is immediate.

Conversely, if $I$ is not reduced, a structural interval has an interior vertex. The endofunctor $j r_s$ of the preceding result fixes every critical vertex and hence both endpoints, but sends that interior vertex to the source of its run. It is not the identity.\hfill$\square$

The obstruction is therefore exactly a nontrivial factorization inside a constant-sign run. Automorphisms cannot change such a factorization because they preserve the graph of indecomposable arrows. General functors can replace it by an identity followed by a longer composite. This is why the automorphism proof on ms.~6--7 is sound while its printed extension to all endofunctors fails.
\end{dcomment}


\section*{3\quad Chemins. Catégorie de chemins d'un type donné $\tau$.}

Soient $I$ un intervalle orienté, et $X$ une catégorie. On appelle \textit{chemin de type $I$ dans $X$} tout foncteur $I \longrightarrow X$, et \textit{catégorie de chemins de type $I$ dans $X$} la catégorie
\[
\Ch_I(X) = \uHom(I,X).
\]
Comme $I$ est \textit{canoniquement} isomorphe à une catégorie intervalle standard $I_0$, on peut identifier les chemins de type $I$, au de type $I_0$. Le type $\tau$ d'un chemin le longueur $n$ $(= \operatorname{card} I)$ dans $X$ peut donc être décrit par une suite de $n$ signes $+$ ou $-$ % [sic] au de type -> à ceux de type, word(s) missing; [sic] le longueur -> de longueur, wrong preposition
\[
\tau_1,\ \tau_2,\ \dots\ ,\ \tau_n \qquad (\tau_i \in \{-1,+1\} \text{ pour } 1 \leq i \leq n).
\]
À cause du fait que $I$ est engendré pas ses flèches de transition, un foncteur $c : I \longrightarrow X$ est déterminé quand on connaît % [sic] pas ses flèches -> par ses flèches, misprint

\begin{enumerate}
\item[a)] Les $c(i)$ $(i \in I)$, i.e. une suite de $n+1$ objets de $X$ :
\[
a_0,\ a_1,\ \dots\ ,\ a_n\ (\text{\footnotemark[5]});
\]
\footnotetext[5]{$a_0$ est dit la \textit{source} du chemin, $a_n$ son but.}

\item[b)] Les $c(\underline{u}_i)$ (où les $\underline{u}_i$ sont les flèches de transition dans $I$), i.e. pour tout $i \in [1,n]$ une flèche $u_i$ dans $X$ entre $a_{i-1}$ et $a_i$, allant de $a_{i-1}$ dans $a_i$ si $\underline{u}_i$ est directe $(\tau_i = +1)$, et de $a_i$ dans $a_{i-1}$ si $\underline{u}_i$ est rétrograde $(\tau_i = -1)$,
\end{enumerate}

et à toute telle donnée a), b) correspond un foncteur. Ainsi on pourra identifier les chemins de type $I$ \msp{page 14} dans $X$ aux diagrammes du type précisé dans $X$. Et les morphismes entre chemins correspondent de même aux morphismes entre diagrammes $a$ et $a'$, une flèche $a \xrightarrow{\alpha} a'$ étant donc une suite de $n+1$ flèches
\[
\alpha_i : a_i \longrightarrow a'_i, \qquad i \in [0,n],
\]
satisfaisant la condition de commutativité pour les $n$ carrés
\[
\begin{tikzcd}
a_{i-1} \arrow[r, "u_i"] \arrow[d, "\alpha_{i-1}"'] & a_i \arrow[d, "\alpha_i"] \\
a'_{i-1} \arrow[r, "u'_i"] & a'_i
\end{tikzcd}
\qquad\qquad \text{(cas $\underline{u}_i$ directe)}
\]

\printedpage{10}
% head: A. Grothendieck : Dérivateurs, Ch. VII, Édition des Universités de Montpellier II et Paris VII
% foot: Version 22 décembre 2001
ou
\[
\begin{tikzcd}
a_{i-1} \arrow[d, "\alpha_{i-1}"'] & a_i \arrow[l, "u_i"'] \arrow[d, "\alpha_i"] \\
a'_{i-1} & a'_i. \arrow[l, "u'_i"']
\end{tikzcd}
\qquad\qquad \text{(cas $\underline{u}_i$ rétrograde)}
\]
Si on a une flèche entre deux sommets quelconque $i$, $j$ dans $I$, avec $i < j$ et $j$ non successeur immédiat de $i$, alors $i$, $j$ appartiennent à un même intervalle critique, et dans $I$ on aura % [sic] quelconque -> quelconques, agreement
\[
i \xrightarrow{\underline{u}_{i+1}} i+1 \xrightarrow{\underline{u}_{i+2}}\ \cdots\ j-1 \xrightarrow{\underline{u}_j} j \qquad \text{une chaîne de longueur } j-i \geq 2 % doubt resolved: PDF p. 12 (edition p. 10) clearly underlines u; the existing reading is retained.
\]
ou
\[
i \longleftarrow i+1 \longleftarrow\ \ \cdots\ j-1 \longleftarrow j,
\]
et la flèche correspondante entre $a_i$ et $a_j$ s'obtient en composant les flèches $u_i$ correspondantes.

\msp{page 15}

\dgrblock{dgrb286}{A fixed path type is a freely generated diagram}
\begin{dcomment}{exposition}{p14}{A fixed path type is a freely generated diagram}
Grothendieck now uses the work on intervals to make path data and path morphisms completely explicit. Fixing the type does not impose any compatibility among the arrows of the path beyond their sources and targets. This is the benefit of having chosen a free path category as the domain.

For example, let the type be $(+,+,-)$. A diagram consists of
\[
 a_0\xrightarrow{u_1}a_1\xrightarrow{u_2}a_2
       \xleftarrow{u_3}a_3.
\]
The only nonidentity arrow of the domain which is not a transition is $0\to2$, and its image must be $u_2u_1$. Given a second diagram $a'_0\to a'_1\to a'_2\leftarrow a'_3$, a map between them consists of four arrows $\alpha_i:a_i\to a'_i$ satisfying
\[
 \alpha_1u_1=u'_1\alpha_0,\qquad
 \alpha_2u_2=u'_2\alpha_1,\qquad
 \alpha_2u_3=u'_3\alpha_3.
\]
Naturality for the longer arrow follows by calculation:
\[
 \alpha_2(u_2u_1)
 =(u'_2\alpha_1)u_1
 =u'_2(u'_1\alpha_0)
 =(u'_2u'_1)\alpha_0.
\]
Every longer constant-sign run is handled by the same induction, with the index order reversed when its arrows are retrograde. Thus the displayed transition squares are both necessary and sufficient for a natural transformation. Identities and composition are taken vertex by vertex; composing commutative squares verifies that the result is again a morphism of paths.

This also gives a useful exact model of the fixed-type category. Let $\operatorname{Arr}(X)=\operatorname{Fun}(\Delta_1,X)$, with ordinary domain and codomain functors $d,c$. For a positive transition its ordered endpoint pair is $(d,c)$; for a negative one it is $(c,d)$. Then $\operatorname{Fun}(I_\tau,X)$ is the iterated strict fibre product of $n$ copies of $\operatorname{Arr}(X)$, matching the second ordered endpoint of one to the first ordered endpoint of the next. On objects this identifies the shared vertices. On arrows it identifies the corresponding vertical components of the naturality squares. The two descriptions are therefore isomorphic as categories, not merely bijective on objects. At length zero the category is $X$ itself.

The construction retains all factorizations and all commuting ladders between them. It has not inverted the retrograde arrows in $X$ and has not identified paths related by a ladder. Those operations would require additional choices or a localization. Here a retrograde transition simply says that the arrow of $X$ points against the chosen order of traversal. Keeping this distinction explicit will matter when the category of types starts to vary.
\end{dcomment}


\textbf{Exemples.}

\begin{enumerate}
\item[1.] $I = \bDelta_0 = e$, $\Ch_I(X) \simeq X$. Les chemins de longueur 0 s'identifient aux objets $a$ de $X$. L'unique élément $\mathrm{id}_e$ de $\mathrm{Fl}(I)$ donne $\mathrm{id}_a$ dans $X$.

\item[2.] $I = \bDelta_1^+$, $\Ch_I(X) \simeq \uFl(X)$. Les chemins de type $\bDelta_1^+$ sont ceux de longueur 1 et \textit{directs,} i.e. donnés par une flèche
\[
a_0 \longrightarrow a_1
\]
dont la source est aussi la source $a_0$ du chemin. Par contre, les chemins de type $\bDelta_1^-$ sont donnés par des diagrammes
\[
a_0 \longleftarrow a_1,
\]
i.e. par des flèches dont le but est considéré comme source du chemin. Bien sûr, les catégories $\Ch_{\bDelta_1^+}(X)$ et $\Ch_{\bDelta_1^-}(X)$ sont canoniquement isomorphes, mais par un isomorphisme qui \textit{échange source et but du chemin.} Nous nous garderons donc bien de confondre les deux chemins précédents.

\msp{page 16}

\item[3.] $\bDelta_2^+$, $\bDelta_2^-$, correspondant aux diagrammes
\[
a_0 \longrightarrow a_1 \longrightarrow a_2
\]
resp.
\[
a_0 \longleftarrow a_1 \longleftarrow a_2,
\]
on peut faire une remarque similaire que pour $\bDelta_1^+$, $\bDelta_1^-$.

\printedpage{11}
% head: A. Grothendieck : Dérivateurs, Ch. VII, Édition des Universités de Montpellier II et Paris VII
% foot: Version 22 décembre 2001

\item[4.] $\Phi_1$, $\Psi_1$, ce sont respectivement les intervalles
\[
\Phi_1 :\ 0 \xleftarrow{\underline{u}_0} 1 \xrightarrow{\underline{u}_1} 2 \qquad\qquad \text{ou} \qquad\qquad
\begin{tikzcd}[row sep=small, column sep=small]
 & 0 \\
1 \arrow[ur] \arrow[dr] & \\
 & 2
\end{tikzcd}
\qquad , % doubt resolved: PDF p. 13 (edition p. 11) clearly underlines u in both labels; the existing reading is retained.
\]
et
\[
\Psi_1 :\ 0 \xrightarrow{\underline{u}_0} 1 \xleftarrow{\underline{u}_1} 2 \qquad\qquad \text{ou} \qquad\qquad
\begin{tikzcd}[row sep=small, column sep=small]
 & 0 \arrow[dl] \\
1 & \\
 & 2 \arrow[ul]
\end{tikzcd}
\qquad .
\]
Ils correspondent respectivement aux diagrammes
\[
a_0 \xleftarrow{u_0} a_1 \xrightarrow{u_1} a_2
\]
et
\[
a_0 \longrightarrow a_1 \longleftarrow a_2.
\]
Les catégories $\Phi_1$, $\Psi_1$ sont opposées l'une de l'autre, et ne sont pas isomorphes. Aussi n'y a-t-il en général aucun isomorphisme entre $\Ch_{\Phi_1}(X)$ et $\Ch_{\Psi_1}(X)$.
\end{enumerate}

\dgrblock{dgrb287}{Zigzags before taking a localization}
\begin{dcomment}{context}{p16:sont opposées l'une de l'autre}{Zigzags before taking a localization}
Grothendieck's paths belong to the same categorical language as fractions, while retaining information which a category of fractions has already identified. In \citet[Chapter~I, \S1.1]{GabrielZisman1967}, localization adjoins inverses to a specified collection of arrows and imposes the relations needed for functoriality and inversion. A span $x\leftarrow z\to y$ can consequently represent a morphism from $x$ to $y$ when its left arrow is among those to be inverted. The fixed-type path category written here makes no such selection: the objects and every component of a natural transformation are ordinary categorical data.

The two-dimensional comparison is equally concrete. The reduced hammocks of \citet[Definition~2.1]{DwyerKan1980Calculating} consist of commutative ladders of zigzags, but their vertical arrows and their left-pointing arrows are required to lie in a specified subcategory of weak equivalences; identity columns are removed and neighbouring columns of the same direction are composed. Their Proposition~3.1 identifies connected components of the resulting mapping spaces with ordinary localized morphisms. Grothendieck's full functor category of one fixed interval is therefore an ingredient of a comparable calculus, not by itself a hammock mapping space. The later commentary on the varying category of types explains the precise relation; here the historical comparison tells us why it is useful to retain ladders rather than only strings of arrows.
\end{dcomment}


Soit $I$ un intervalle orienté, et considérons l'\textit{intervalle opposé $I^{\opp}$}, ayant même catégorie sous-jacente mais l'orientation opposé. Comme la définition de la catégorie des chemins du type $I$ ne fait pas intervenir l'orientation de $I$, on a donc un isomorphisme canonique % [sic] l'orientation opposé -> l'orientation opposée, agreement
\[
\Ch_I(X) \xrightarrow{\sim} \Ch_{I^{\opp}}(X).
\]

\msp{page 17}

\dgrblock{dgrb288}{A span and a cospan give different diagram categories}
\begin{dcomment}{examples}{p16}{A span and a cospan give different diagram categories}
Grothendieck emphasizes that the direction of a turn changes the fixed-type category. The difference can be seen without any localization or homological algebra. We give a finite example which also distinguishes equivalence of categories, not only strict isomorphism.

Let $X$ be the three-object poset
\[
             a\longleftarrow b\longrightarrow c.
\]
A span diagram in $X$ is a triple $(x_0,x_1,x_2)$ with $x_1\leq_X x_0$ and $x_1\leq_X x_2$. If its middle value is $b$, each endpoint can independently be $a,b,$ or $c$, giving nine diagrams. If the middle value is $a$, both endpoints must be $a$; if it is $c$, both must be $c$. Thus
\[
 \bigl|\operatorname{Ob}\operatorname{Fun}(\Phi_1,X)\bigr|
 =3^2+1^2+1^2=11.
\]
A cospan diagram has $x_0\leq_X x_1$ and $x_2\leq_X x_1$. For middle value $b$ there is one possibility. For middle value $a$, both endpoints independently lie in $\{b,a\}$, giving four; for middle value $c$, both independently lie in $\{b,c\}$, again giving four. Hence
\[
 \bigl|\operatorname{Ob}\operatorname{Fun}(\Psi_1,X)\bigr|
 =1^2+2^2+2^2=9.
\]
Morphisms in both functor categories are completely determined: there is one from $(x_0,x_1,x_2)$ to $(y_0,y_1,y_2)$ exactly when $x_i\leq_Xy_i$ for each $i$. The naturality squares then commute automatically because $X$ is thin. These categories are posets and have no distinct isomorphic objects. An equivalence between posets is an isomorphism, since essential surjectivity is actual surjectivity and full faithfulness reflects their order. Their different object counts therefore rule out an equivalence.

Orientation reversal is a separate operation. For the span, interchanging $x_0$ and $x_2$ gives another span and an automorphism of its diagram category. It leaves the middle vertex and both arrow directions of the underlying span shape intact, while exchanging the two endpoint evaluations. Passing from the span to the cospan reverses the categorical arrows at the turn. The preceding count detects exactly that change.
\end{dcomment}


\dgrblock{dgrb289}{When a change of orientation becomes a stable equivalence}
\begin{dcomment}{later}{p16}{When a change of orientation becomes a stable equivalence}
The span--cospan distinction remains real for the ordinary functor categories which Grothendieck has just defined. A later theorem explains why analogous diagram categories can nevertheless become equivalent after one passes to a stable homotopy theory. For a finite oriented tree $Q$, \citet[Theorem~9.19 and Corollary~9.23, arXiv v3]{GrothStovicek2016Trees} construct reflection equivalences and deduce strong stable equivalence between any two orientations of the same underlying tree. In particular, for a stable derivator $\mathbb D$,
\[
 \mathbb D^{\Phi_1}\simeq\mathbb D^{\Psi_1}.
\]
These are equivalences of shifted derivators, natural with respect to exact morphisms. They are constructed using stable operations such as fibres, cofibres, and biproducts; they are not induced by an isomorphism of the span and cospan categories. The finite example above therefore presents no obstruction. It concerns arbitrary ordinary targets, whereas the reflection theorem uses stability to transform the values and arrows of the diagram. This gives a precise later setting in which changing the signs of an interval preserves a diagram theory even though it does not preserve its indexing category.
\end{dcomment}


Néanmoins on se gardera de confondre un chemin de type $I$ avec le chemin opposé de type $I^{\opp}$, la source de l'un étant le but de l'autre et inversement. C'est dire qu'un chemin doit être vu comme \textit{plus} qu'un simple foncteur
\[
I \longrightarrow X,
\]
il faut de plus se donner l'\textit{orientation de $I$} comme faisant partie de la structure de chemin. Donc la définition formelle raisonnable, qui colle à l'intuition du chemin comme allant d'une \textit{source} bien définie vers un \textit{but} également bien défini, c'est qu'un chemin (défini en termes d'une catégorie intervalle $I$) est la donné d'un couple $(\omega,c)$, où $\omega$ est une orientation de $I$, et $c$ un foncteur $c : I \longrightarrow X$. Quand on regarde la catégorie des chemins d'un `type' déterminé, ce \textit{type} est $(I,\omega)$, et non $I$. C'est quand $\omega$ est fixé une bonne fois que cette catégorie s'identifie impunément à $\uHom(I,X)$ sans plus. % [sic] la donné -> la donnée, missing final e

Il faut voir cette catégorie de chemins comme munie de la structure essentielle
\[
\begin{tikzcd}
\Ch_{(I,\omega)}(X) \arrow[d, "{(\sigma,\beta)}"] \\
X \times X
\end{tikzcd}
\qquad
\begin{array}{l}
c \longmapsto (\sigma(c),\beta(c)) \\
\text{ou $\sigma(c)$ est la \textit{source,}} \\
\qquad \text{$\beta(c)$ le \textit{but} de $c$.}
\end{array} % [sic] ou sigma(c) -> où sigma(c), missing accent
\]

\printedpage{12}
% head: A. Grothendieck : Dérivateurs, Ch. VII, Édition des Universités de Montpellier II et Paris VII
% foot: Version 22 décembre 2001

\msp{page 18}

\dgrblock{dgrb290}{The endpoints are part of the fixed-type object}
\begin{dcomment}{exposition}{p17}{The endpoints are part of the fixed-type object}
Grothendieck's insistence on orientation is most transparent over the category of endpoint pairs. The category attached to an oriented type $(I,\omega)$ is not merely $\operatorname{Fun}(I,X)$; it comes with the functor
\[
 (\sigma,\beta):\operatorname{Fun}(I,X)\longrightarrow X\times X.
\]
On a natural transformation $\alpha:c\Rightarrow d$, its value is $(\alpha_0,\alpha_n)$. Thus a morphism of paths may move the endpoints, and the displayed functor records exactly how. Its strict fibre over $(x,y)$ consists of paths with endpoints $x,y$ and natural transformations which are the identity at both endpoints. This fibre is a concrete category of fixed-endpoint ladders.

Let $\omega'$ be the reverse orientation on the same underlying interval category. The identity on that underlying category identifies the two ordinary functor categories; under this identification the endpoint pair changes from $(c(0),c(n))$ to $(c(n),c(0))$. Consequently the comparison is a functor over the transposition $s(x,y)=(y,x)$. If $x\ne y$, it carries the fibre over $(x,y)$ to the fibre over $(y,x)$, not to itself. This explains why equality of the underlying functor categories does not erase the distinction between a path and its reverse.

The endpoint map $\alpha:\{0,1\}_{\mathrm{disc}}\to I$ packages the same structure as restriction. For positive length its two images are distinct. For length zero it is the fold onto the sole vertex, so the word ``inclusion'' on the next page needs that qualification; restriction is then the diagonal $X\to X\times X$. All the formulas still hold. The subsequent discussion of symmetric types compares this identity-on-underlying-diagram reversal with the additional comparison induced by a reflection of the shape. Their difference belongs to the structure over $X\times X$, which is why the base category has been kept visible.
\end{dcomment}


Ce foncteur s'identifie au foncteur de composition
\[
c \longmapsto c \circ \alpha
\]
où $\alpha$ est le foncteur d'inclusion
\[
\left\{
\begin{array}{lll}
\varepsilon \xrightarrow{\alpha} I & & \\[1ex]
\varepsilon = \{0,1\}, & \multicolumn{2}{l}{\text{catégorie discrète à deux éléments,}} \\[1ex]
\alpha \text{ donné par} & \alpha(0) = & \text{source de } I \\[1ex]
 & \alpha(1) = & \text{but de } I.
\end{array}
\right.
\]
Si $\omega'$ est l'orientation opposée de $I$, on a un diagramme commutatif
\[
\begin{tikzcd}
\Ch_{(I,\omega)}(X) \arrow[r, "\varphi", "\sim"'] \arrow[d, "{(\sigma,\beta)}"] & \Ch_{(I,\omega')}(X) \arrow[d, "{(\sigma',\beta')}"] \\
X \times X \arrow[r, "s"] & X \times X,
\end{tikzcd}
\]
où $s$ est le foncteur de symétrie
\[
s(x,y) = (y,x)
\]
(\footnote[6]{Le foncteur $\varphi$ est celui qui associe à tout chemin $c$, de type $\tau$, le \textit{chemin opposé $c^{\opp}$} ou $c^{-1}$, de type $\tau^{\opp}$. Si $\tau = \tau^{\opp}$, c'est donc un chemin \textit{de même type,} et $\varphi$ est un \textit{automorphisme} de $\Ch_\tau(X)$. Cet automorphisme (fonctoriel en $X$) est déduit d'un automorphisme $\theta$ de $I$, cf. ci-dessous.}). Donc l'isomorphisme $\varphi$ ne respecte pas les flèches structurales `source-but' $(\sigma,\beta)$.

Il y a cependant un embarras quand il existe un isomorphisme d'intervalles \textit{orientés}
\[
\theta : (I,\omega) \xrightarrow{\sim} (I,\omega'),
\]
i.e. quand $I$ est symétrique. Ce n'est pas le cas pour les $\bDelta_i$ $(i \geq 1)$ (ces catégories n'ont pas d'autres automorphismes que l'identité), mais c'est le cas pour les catégories $\Phi_1$, $\Psi_1$ de l'exemple $4^o$, plus généralement \msp{page 19} pour les catégories réduites de longueur \textit{paire.} Comme chaque fois qu'on a un isomorphisme $\theta$ d'intervalles \textit{orientés,} on trouve un isomorphisme entre les catégories de chemins de type $(I,\omega)$ et $(I,\omega)^{\opp} = (I,\omega')$, \textit{compatible} cette fois avec les foncteurs $(\sigma,\beta)$ :
\[
\begin{tikzcd}
\Ch_{(I,\omega)}(X) \arrow[rr, "\Psi_\theta", "\sim"'] \arrow[dr, "{(\sigma,\beta)}"'] & & \Ch_{(I,\omega')}(X) \arrow[dl, "{(\sigma',\beta')}"] \\
 & X \times X. &
\end{tikzcd}
\]

\printedpage{13}
% head: A. Grothendieck : Dérivateurs, Ch. VII, Édition des Universités de Montpellier II et Paris VII
% foot: Version 22 décembre 2001

\msp{page 20}

Il est bon de regarder un \textit{type} de chemins comme une donné non catégorique, mais combinatoire, une simple suite $\tau$ de $n$ signes $\tau_i$ $(1 \leq i \leq n)$ $+$ ou $-$ ($n =$ longueur des chemins envisagés). Cette donné combinatoire définit un intervalle orienté standard, $I_\tau$. Quand on a un intervalle orienté quelconque $(I,\omega)$, il définit un type $\tau$, et on l'identifie canoniquement au type $I_\tau$. Il vaut mieux écrire $\Ch_\tau(X)$, plutôt que $\Ch_{(I,\omega)}(X)$, pour bien indiquer que pour $X$ donné, c'est là une catégorie totalement `épinglée', sans automorphismes liés à la nature particulière de $\tau$ (à des automorphismes de $\tau$!). % [sic] une donné -> une donnée, missing final e; [sic] Cette donné -> Cette donnée, missing final e

On peut préciser ce point ainsi. Soit $\tau'$ un autre type d'intervalles, et proposons-nous de déterminer des `lois' qui permettent à tout chemin de type $\tau$ dans une catégorie arbitraire $X$, d'associer un chemin de type $\tau'$ de cette même catégorie, et ceci de façon `fonctorielle', i.e. compatible avec les applications horizontales
\[
\begin{tikzcd}
\mathrm{Ch}_\tau(X) \arrow[r, "\mathrm{Ch}_\tau(f)"] \arrow[d, "L_X"'] & \mathrm{Ch}_\tau(X') \arrow[d, "L_{X'}"] \\
\mathrm{Ch}_{\tau'}(X) \arrow[r, "\mathrm{Ch}_{\tau'}(f)"] & \mathrm{Ch}_{\tau'}(X')
\end{tikzcd}
\]
associées à un foncteur $f : X \longrightarrow X'$.

\msp{page 21}

\dgrblock{dgrb291}{Operations on paths are maps of their shapes}
\begin{dcomment}{exposition}{p20}{Operations on paths are maps of their shapes}
Grothendieck wants to do the combinatorics once, in a category small enough that every operation can be seen. A path in an arbitrary category carries many labels, but an operation valid for every category cannot depend on accidental properties of those labels. Its entire content is already present in the universal path. The passage from an oriented interval to its sign word makes this universal object unambiguous. It does not assert that a particular category of paths has no automorphisms; it removes the ambiguity coming from presenting the same unoriented interval in two different ways.

Write $P_I(X)=\Hom_{\Cat}(I,X)$ for the \emph{set} of paths of shape $I$, keeping $X^I=\uHom(I,X)$ for their category. The distinction between these two expressions matters in the displayed Yoneda argument. If
\[
 L_X:P_I(X)\longrightarrow P_J(X)
\]
is natural in all functors $X\to Y$, put $\varphi=L_I(\mathrm{id}_I):J\to I$. Applying naturality to the functor $c:I\to X$ gives
\[
 L_X(c)=c\circ L_I(\mathrm{id}_I)=c\varphi.
\]
Thus every such operation is precomposition with exactly one $\varphi$, and every $\varphi$ defines such an operation. This is the covariant Yoneda lemma, with the representing objects appearing in the reverse order; see \citet[Theorem~2.2.4 and Corollary~2.2.8]{Riehl2016}. No hypothesis about interval categories is needed for this argument.

There is then a canonical operation on morphisms of paths. If $v:c\Rightarrow d$, set $(L_Xv)_j=v_{\varphi(j)}$. For an arrow $a:j\to k$ of $J$, its naturality square is the naturality square of $v$ at $\varphi(a)$. Identities and composites are preserved component by component. Consequently the operation lifts to the functor
\[
 \varphi^*:X^I\longrightarrow X^J.
\]
This proves the extension asserted on manuscript page~23, without assuming a separate coherence theorem. Natural transformations between changes of shape have an equally concrete meaning: $\alpha:\varphi\Rightarrow\psi$ gives $c\alpha:c\varphi\Rightarrow c\psi$, naturally in $c$ and $X$.

The special advantage of intervals now enters. They are posets, so a functor into an interval is determined by its object function, and a natural transformation between two such functors exists precisely when the corresponding pointwise arrows exist. If it exists, it is unique. This uniqueness makes every required diagram of \emph{shape transformations} commute. It does not make the categories $X^I$ thin when $X$ is an arbitrary category: those categories still remember all morphisms of $X$.

For fixed endpoints the same universal-path argument is even more economical. Suppose $I$ and $J$ have specified sources and targets. The operation $c\mapsto c\varphi$ preserves both endpoints for every $X$ and $c$ exactly when $\varphi(0_J)=0_I$ and $\varphi(1_J)=1_I$. Necessity follows by taking $X=I$ and $c=\mathrm{id}_I$; sufficiency is direct evaluation. This is the reason the first category of all changes of shape will soon have to be restricted.
\end{dcomment}


Par exemple, prenant le type $\bDelta_2^+$ (i.e. $+1,+1$), correspondant aux chemins
\[
a_0 \xrightarrow{u_1} a_1 \xrightarrow{u_2} a_2,
\]
on peut associer à un tel chemin le chemin de type \texttt{[}$\bDelta_1^+$\texttt{]}
\[
a_0 \xrightarrow{u_2u_1} a_2,
\]
ou aussi le chemin de type $\bDelta_2^+$
\[
a_0 \xrightarrow{u_2u_1} a_2 \xrightarrow{\mathrm{id}} a_2,
\]
etc. Comme le foncteur
\[
\left\{
\begin{array}{cccl}
X & \longmapsto & \mathrm{Ch}_\tau(X) & \qquad (= \mathrm{Ob}\, \Ch_\tau(X)) \\[1ex]
\mathrm{Cat} & \longrightarrow & (\mathrm{Ens}) &
\end{array}
\right.
\]
est \textit{représenté} par l'objet $I_\tau$ de $X$ \texttt{[plutôt de} Cat\texttt{]}, le sorite des foncteurs représentables nous dit que de telles lois de formation de chemins de type $\tau$ à partir des chemins de type $\tau'$ \texttt{[plutôt} $\tau'$ \texttt{à partir des chemins de type} $\tau$\texttt{]} correspondent biunivoquement aux flèches dans Cat
\[
I_{\tau'} \xrightarrow{\varphi} I_\tau\ ,
\]

% ---- from Der-ChVII-456_p016-020.tex
\printedpage{14}
% head: A. Grothendieck : Dérivateurs, Ch. VII, Édition des Universités de Montpellier II et Paris VII
% foot: Version 22 décembre 2001

en prenant simplement la ‘loi’ correspondante à $\varphi$
\[
c \longmapsto c\circ\varphi \qquad \text{pour } c : I_\tau \longrightarrow X.
\]
De façon plus précise, l'application naturelle
\[
\underbrace{\Hom_{\Cat}(I_{\tau'},I_\tau)}_{=\mathrm{Ob}\,\uHom(I_{\tau'},I_\tau)} \longrightarrow \Hom(\mathrm{Ch}_\tau,\mathrm{Ch}_{\tau'}),
\]
\msp{page 22}
où le deuxième Hom est considéré comme un Hom d'objets dans $\Cat^{\vee}=\uHom(\Cat,\mathrm{Ens})$.
\par\Needspace{70mm}\noindent
Ainsi, dans les deux exemples plus haut, où $\tau=\bDelta_2^+$ et $\tau'=\bDelta_1^+$ ou $\bDelta_2^+$, les foncteurs en sens inverse sont
\[
\begin{tikzcd}
\bDelta_2^+\ \Big\{ & 0 \arrow[r] & 1 \arrow[r] & 2 \\
\bDelta_1^+\ \Big\{ & 0 \arrow[u, mapsto] \arrow[rr] & & 1 \arrow[u, mapsto]
\end{tikzcd}
\]
et
\[
\begin{tikzcd}
\bDelta_2^+\ \Big\{ & 0 \arrow[r] & 1 \arrow[r] & 2 \\
\bDelta_2^+\ \Big\{ & 0 \arrow[u, mapsto] \arrow[r] & 1 \arrow[ur, mapsto] \arrow[r] & 2\,. \arrow[u, mapsto]
\end{tikzcd}
\]
C'est là un principe général pour ramener toutes les ‘opérations’ $\pm$ canoniques sur des chemins dans des catégories arbitraires, et des énoncés les concernant, à des questions où $X$ disparaît et où on travaille dans les catégories intervalles elles-mêmes. Un des avantages techniques, c'est que ce sont des catégories fort petites, donnant lieu de plus à une intuition topologique sûre (qui risque de se volatiliser par moments en travaillant dans des catégories de chemins), et aussi que dans de telles catégories, donc aussi dans les catégories de foncteurs entre elles, il y a toujours au plus
\msp{page 23}
\textit{une seule flèche} allant d'un objet donné vers un autre. En d'autres termes, dans de telles catégories tous les diagrammes sont commutatifs.

Bien sûr, on voit aussi que si on a une loi de formation $L$ de chemins de type $\tau'$ à partir de chemins de type $\tau$ (avec la condition de fonctorialité par rapport à la catégorie $X$ où on prend les chemins), cette loi est nécessairement \textit{fonctorielle par rapport aux chemins :} si on a un homomorphisme
\[
c \xrightarrow{v} c' \qquad \text{dans } \uCh_\tau(X),
\]
\printedpage{15}
% head: A. Grothendieck : Dérivateurs, Ch. VII, Édition des Universités de Montpellier II et Paris VII
% foot: Version 22 décembre 2001
alors la loi $L$ ‘se prolonge’ à de tels $v$ pour donner
\[
L(c) \xrightarrow{L(v)} L(c'),
\]
avec les relations habituelles
\[
L(\mathrm{id}_C)=\mathrm{id}_{L(c)}, \qquad L(v'v)=L(v')L(v). % [sic] \mathrm{id}_C -> \mathrm{id}_c, capital C printed while the chemin is c (cf. L(c))
\]
Cela résulte de la description de $L$ comme une opération $L_\varphi$ en termes d'un foncteur bien déterminé
\[
I_{\tau'} \xrightarrow{\varphi} I_\tau,
\]
lequel à son tour donne pour toute $X$ un \textit{foncteur}
\[
\left\{
\begin{array}{rcl}
\uCh_\tau(X) & \longrightarrow & \uCh_{\tau'}(X)\\[1ex]
c & \longmapsto & c\circ\varphi
\end{array}
\right.
\]
plus précis que simplement un \textit{application} $\mathrm{Ch}_\tau(X)\longrightarrow\mathrm{Ch}_{\tau'}(X)$ entre les ensembles d'objets correspondants. % [sic] un application -> une application, gender agreement

\dgrblock{dgrb292}{Composition and insertion of an identity}
\begin{dcomment}{examples}{p23:entre les ensembles d'objets correspondants.}{Composition and insertion of an identity}
The two examples on manuscript pages~21--22 illustrate how a familiar operation becomes a finite calculation before any coefficient category is chosen. Take $I=[2]=(0\to1\to2)$ and $J=[1]$. The unique endpoint-preserving functor $\varphi:J\to I$ sends $0$ to $0$ and $1$ to $2$. A path $c=(a_0\xrightarrow{u}a_1\xrightarrow{v}a_2)$ therefore becomes the arrow $vu:a_0\to a_2$.

A morphism from $c$ to $d=(b_0\xrightarrow{u'}b_1\xrightarrow{v'}b_2)$ is a triple $(h_0,h_1,h_2)$ satisfying $h_1u=u'h_0$ and $h_2v=v'h_1$. The induced morphism between the composite arrows is the square with sides $h_0,h_2$. Its commutativity follows from the explicit computation
\[
 h_2vu=v'h_1u=v'u'h_0.
\]
The middle component has done its work and can be forgotten. The reversed variance is essential: an operation which reduces the length of a path is represented here by a functor from the shorter interval to the longer one.

For the second operation let $J=I=[2]$ and take the nondecreasing object function $\psi=(0,2,2)$. Then $\psi^*c=(a_0\xrightarrow{vu}a_2\xrightarrow{1}a_2)$, and $\psi^*h=(h_0,h_2,h_2)$. The first naturality square is the square just computed; the second is $h_2\,1=1\,h_2$. Since $\psi\psi=\psi$, the resulting operation on paths is strictly idempotent. This equality is checked on the three objects of $[2]$, rather than on all categories and all pairs of composable arrows.

Endpoint preservation alone has not made such maps rigid. On $[2]$ the three nondecreasing endpoint-preserving endofunctors have object functions $(0,0,2)$, $(0,1,2)$, and $(0,2,2)$. They respectively place the composite in the second edge, leave the path unchanged, and place the composite in the first edge. This also supplies a concrete reminder of the correction to Proposition~1(b): fixing the endpoints does not force an arbitrary endofunctor of a nonreduced interval to be the identity. The rigidity there belongs to isomorphisms, while the present theory deliberately needs the noninvertible maps.
\end{dcomment}


\msp{page 24}
On a le choix maintenant entre

\begin{enumerate}[label=\alph*)]
\item Une étude plus ou moins systématique des foncteurs entre intervalles. Je préfère attendre pour cela d'avoir besoin de ceci et de cela, et de ne faire que le strict minimum pour les constructions que j'ai en vue.
\item Une description de la compositions des chemins. % [sic] la compositions -> la composition, number agreement
\item Un formalisme des catégories de chemins \textit{de type non précisé.}
\end{enumerate}

Comme mon objectif est surtout de décrire de telles catégories $\uCh(X)$ (sans indice $\tau$!), avec des bonnes propriétés, je commencerai par c), regarderai ensuite b) comme un test crucial pour l'utilisabilité des notions obtenues, et enfin m'occuperai de a) quand je saurai de quoi j'aurai en besoin.

\printedpage{16}
% head: A. Grothendieck : Dérivateurs, Ch. VII, Édition des Universités de Montpellier II et Paris VII
% foot: Version 22 décembre 2001

\section*{4\quad Catégorie de chemins $\uCh(X)$}

Je veux ‘inclure’ toutes les catégories $\uCh_\tau(X)$ de chemins de tous les types qu'on veut, dans une même catégorie $\uCh(X)$. Rappelons
\[
\uCh_\tau(X)=\uHom(I_\tau,X),
\]
où $I_\tau$ est l'intervalle (orienté) standard de type $\tau$.

L'ensemble de tous les types d'intervalles
\msp{page 25}
peut être identifié au monoïde libre Ty à deux générateurs $\Delta^+$ et $\Delta^-$, dont les éléments sont les suites finies quelconques (eventuellement vides) d'objets alternativement égales à $\Delta^+$ et à $\Delta^-$. Si le premier élément est $\Delta^+$, on a un type de chemin à direction initiale directe, rétrograde dans le cas contraire. La décomposition de $I_\tau$, ou de $\mathrm{Fl}\,(I_\tau)$, en intervalles critiques correspond à l'écriture réduite d'un mot comme % [sic] eventuellement -> éventuellement, missing accent
\[
(\Delta^+)^{l_1}(\Delta^-)^{l_2}\ \cdots\ (\Delta^?)^{l_r} \qquad (\Delta^?=\Delta^+ \text{ si } r \text{ impair, } \Delta^?=\Delta^- \text{ si } r \text{ pair}),
\]
ou comme
\[
(\Delta^-)^{l_1}(\Delta^+)^{l_2}\ \cdots\ (\Delta^?)^{l_r} \qquad (\Delta^?=\Delta^- \text{ si } r \text{ impair, } \Delta^?=\Delta^+ \text{ si } r \text{ pair}).
\]
Les $l_i$ sont les longueurs des intervalles critiques successifs (ou des chaines structurales successives), leur nombre est la longueur réduite de l'intervalle.

On va considérer Ty comme l'ensemble des objets d'une catégorie $\uTy$ :
\[
\left\{
\begin{array}{l}
\mathrm{Ob}\,\uTy=\Ty=\ \text{monoïde libre à deux générateurs } \Delta^+,\ \Delta^- \text{ précédents.}\\[1ex]
\text{Si } \tau,\tau'\in\Ty,\ \Hom(\tau,\tau')=\Hom_{\Cat}(I_\tau,I_{\tau'}).
\end{array}
\right.
\]

(\footnote[7]{Définition un peu canulée, il faut se borner aux homomorphismes de $I_\tau$ dans $I_{\tau'}$ compatibles avec source et but - cf. p.~29.}). On notera qu'on a non seulement un ensemble $\Hom_{\Cat}(I_\tau,I_{\tau'})$, mais même une catégorie $\uHom(I_\tau,I_{\tau'})$, ce qui revient ici à dire que $\Hom(I_\tau,I_{\tau'})$ (i.e. $\Hom(\tau,\tau')$) peut être vu comme un ensemble \textit{ordonné.} Mais pour l'instant nous allons ignorer cette relation d'ordre.

\msp{page 26}
La composition des morphismes est la composition des foncteurs entre catégories $I_\tau$. Ainsi on a un foncteur \textit{pleinement fidèle} (\footnote[8]{faux, cf. annotation marginale page précédente - on ne prend pas \textit{tous} les foncteurs dans Cat entre $I_\tau$ et $I_{\tau'}$.}) canonique
\[
\tau \longmapsto I_\tau : \uTy \xrightarrow{\rho} (\Cat) \qquad \text{(\footnotemark[9]),}
\]
\footnotetext[9]{foncteur \textit{réalisation} d'un type (combinatoire) d'intervalles.}% footnote text printed at the foot of p. 17
\printedpage{17}
% head: A. Grothendieck : Dérivateurs, Ch. VII, Édition des Universités de Montpellier II et Paris VII
% foot: Version 22 décembre 2001
par quoi la catégorie $I_\tau$ \texttt{[plutôt} $\uTy$\texttt{]} pourrait s'identifier à une sous-catégorie pleine de (Cat) (celle des $I_\tau$). Mais je préfère ne pas identifier $\tau\in\Ty$ (donnée combinatoire) avec $I_\tau$ (donnée catégorique).

\dgrblock{dgrb293}{Three successive categories of types}
\begin{dcomment}{exposition}{p25}{Three successive categories of types}
Grothendieck is testing which changes of parameter should count as morphisms of paths. The changes of definition in these pages are purposeful, but a notation retained through all three stages can obscure which fullness assertion remains true. We distinguish the stages here, without altering the source's notation. All three categories have the finite sign words as objects, including the empty word $e$.

Let $T_{\mathrm{all}}$ have every functor $I_\tau\to I_\upsilon$ as a morphism; let $T_{\partial}$ retain just the functors preserving the two chosen endpoints; and let $T_{\uparrow}$ retain those which also preserve the external, linear order of vertices. Thus
\[
 T_{\uparrow}\subseteq T_{\partial}\subseteq T_{\mathrm{all}}.
\]
The category eventually denoted by $\uTy$ in these pages is $T_{\uparrow}$. On page~25 it is initially $T_{\mathrm{all}}$, while pages~29--31 pass through $T_{\partial}$. For $T_{\mathrm{all}}$, realization in $\Cat$ is fully faithful by definition. For $T_{\partial}$, realization is fully faithful into $\varepsilon\backslash\Cat$, where $\varepsilon$ is the discrete category on two objects; forgetting $\varepsilon$ loses fullness. For $T_{\uparrow}$ even the embedding into $\varepsilon\backslash\Cat$ is generally not full. The footnotes retracting fullness on page~26 and the assertion of fullness on page~31 are therefore compatible when read at their respective stages.

Here is a small witness to the last distinction. From the alternating interval $I_{+-+}$ to $I_+$ there is a functor with object function
\[
 (0,1,2,3)\longmapsto(0,1,0,1).
\]
Every forward transition goes to $0\to1$ and the backward transition, whose categorical direction is from its right vertex to its left vertex, also goes to $0\to1$. Both endpoints are preserved, but the vertex function is not nondecreasing. Thus it belongs to $T_{\partial}$ and not to $T_{\uparrow}$. The monotonicity condition excludes this backtracking.

The sign words themselves form a free monoid on $+$ and $-$. Arbitrary adjacent repetitions are allowed; the alternating feature is the alternation of signs between \emph{maximal constant runs}. A nonempty word has a unique expression $\varepsilon_1^{l_1}\cdots\varepsilon_r^{l_r}$ with $l_i>0$ and $\varepsilon_{i+1}=-\varepsilon_i$. Its length is $\sum l_i$, its useful length is $r$, and its critical intervals are these runs with their endpoint vertices restored. The empty word has length and useful length zero and no initial sign. This last exception should remain in force whenever the text speaks of the first generator.

In the undercategory description, $\varepsilon\to I_e$ identifies the two objects of $\varepsilon$. It is a structure map, not an embedding; the hooked arrows drawn on page~29 should not be read as asserting injectivity in this degenerate case. Keeping $e$ is indispensable: it will be the unit for concatenation.
\end{dcomment}


Soit $X$ une (petite?) catégorie. On a un contrafoncteur
\[
\left\{
\begin{array}{rcl}
S & \longmapsto & \uHom(S,X)\\[1ex]
(\Cat)^o & \longrightarrow & (\Cat)
\end{array}
\right.
\]
défini par $X$, d'où en composant avec $\rho$ un contrafoncteur
\[
\uTy^{o} \xrightarrow{\mathrm{CH}_X} \Cat.
\]
C'est le contrafoncteur
\[
\left\{
\begin{array}{l}
\tau \longmapsto \uHom(I_\tau,X) \overset{\text{déf}}{=} \uCh_\tau(X),\\[1ex]
\text{i.e. } \mathrm{CH}_X(\tau)=\uCh_\tau(X).
\end{array}
\right.
\]
Mais un contrafoncteur sur une petite catégorie $\uTy$ donne lieu à une \textit{catégorie fibrée} (et même scindée) sur celle-ci.

\msp{page 27}
Je vais désigner cette catégorie fibrée par $\uCh(X)$. (Provisoirement - si elle s'avère trop grosse, je la restreindrai \dots) Ainsi on a
\[
\begin{tikzcd}[row sep=large]
\uCh(X) \arrow[d, "\pi_X"] \\
\uTy\ .
\end{tikzcd}
\]
La fibre de $\uCh(X)$ en $\tau\in\mathrm{Ob}\,\uTy$ est la catégorie $\uCh_\tau(X)$ des chemins de type $\tau$ dans $X$. Cela nous dit donc quelles sont les flèches dans $\uCh(X)$ au dessus des flèches identiques dans $\uTy$. Mais que sont les flèches plus générales? On se donne une flèche dans $\uTy$
\[
\varphi : \tau' \longrightarrow \tau, \qquad \text{i.e. } I_{\tau'} \xrightarrow{\varphi} I_\tau,
\]
quelles sont les flèches de $\uCh(X)$ au dessus de $\varphi$? Par le sorite général des catégories fibrées, si
\[
\begin{array}{lclcl}
c\in\uCh(X)_\tau & = & \uCh_\tau(X) & \overset{\text{déf}}{=} & \uHom(I_\tau,X)\\[1ex]
c'\in\uCh(X)_{\tau'} & = & \uCh_{\tau'}(X) & \overset{\text{déf}}{=} & \uHom(I_{\tau'},X)
\end{array}
\]
on a
\[
\boxed{\Hom_\varphi(c',c)\simeq\Hom(c',\underbrace{\varphi^*(c)}_{=c\circ\varphi}).}
\]
Donc : les homomorphismes du
\printedpage{18}
% head: A. Grothendieck : Dérivateurs, Ch. VII, Édition des Universités de Montpellier II et Paris VII
% foot: Version 22 décembre 2001
\msp{page 28}
chemin $c'$ de type $\tau'$ dans $X$, dans le chemin $c$ de type $\tau$, ‘au dessus’ du morphisme $\varphi$ de $\tau'$ dans $\tau$, i.e. du foncteur
\[
\varphi : I_{\tau'} \longrightarrow I_\tau,
\]
est un homomorphisme $v$ de $c'$ dans $c\circ\varphi$, i.e.
\[
\begin{tikzcd}[row sep=large]
I_{\tau'} \arrow[r, "\varphi"] \arrow[dr, "c'"', ""{name=A}] & I_\tau \arrow[d, "c", ""'{name=B}] \\
& X.
\arrow[from=A, to=B, bend left=20, shorten <=6pt, shorten >=2pt, "\circled{v}"]
\end{tikzcd}
\]
La composition des flèches est évidente :
\[
\begin{tikzcd}[row sep=large]
I_{\tau''} \arrow[r, "\varphi'"] \arrow[drr, "c''"', ""{name=C}] & I_{\tau'} \arrow[r, "\varphi"] \arrow[dr, "c'"', ""{name=A}] & I_\tau \arrow[d, "c", ""'{name=B}] \\
& & X,
\arrow[from=A, to=B, bend left=20, shorten <=6pt, shorten >=2pt, "v"]
\arrow[from=C, to=A, shorten <=4pt, shorten >=4pt, "v'"]
\end{tikzcd}
\]
on a
\[
\begin{array}{rccl}
v : & c' & \longrightarrow & c\circ\varphi, \quad \text{d'où } c'\circ\varphi' \xrightarrow{v*\varphi'} \underbrace{(c\circ\varphi)\circ\varphi'}_{=c\circ(\varphi\varphi')}\\[3ex]
v' : & c'' & \longrightarrow & c'\circ\varphi'
\end{array}
\]
et
\[
v''=(v*\varphi')\circ v', \qquad \text{donc} \qquad
\begin{tikzcd}
c'' \arrow[r, "v'"] \arrow[rr, bend right=12, "\circled{v''}"'] & c'\circ\varphi' \arrow[r, "v*\varphi'"] & c\circ(\varphi\varphi').
\end{tikzcd}
\]
Mais a-t-on des foncteurs source et but
\[
\uCh(X) \overset{\sigma}{\underset{\beta}{\rightrightarrows}} X \qquad
\begin{array}{ccl}
c & \longmapsto & \text{source}(c)\\[1ex]
c & \longmapsto & \text{but}(c)
\end{array}
\qquad ?
\]
Il faudrait donc qu'un homomorphisme $v$ comme dessus
\[
v : c' \longrightarrow c\circ\varphi
\]
définisse
\[
\begin{array}{lrcl}
\sigma(c) : & \text{source}(c') & \longrightarrow & \text{source}(c)\\[1ex]
\beta(c) : & \text{but}(c') & \longrightarrow & \text{but}(c).
\end{array}
\]

\dgrblock{dgrb294}{The split fibration of all paths}
\begin{dcomment}{context}{p27}{The split fibration of all paths}
Grothendieck now uses the machinery of descent to assemble paths of different shapes. The construction is the one in \citet[Expos\'e~VI, \S\S8--9]{GrothendieckSGA1}: a strict contravariant functor to categories supplies a category with a chosen, strictly composable family of Cartesian lifts. The commentary on Chapter~VI, manuscript pages~26 and~53--56, discussed splittings and their universal constructions. Here the splitting needs no choices, because it is given by literal precomposition.

For any of the three shape categories $T$, put $F_X(\tau)=X^{I_\tau}$. An object of its Grothendieck construction is $(\tau,c)$, and a morphism
\[
 (\tau',c')\longrightarrow(\tau,c)
\]
is a pair $(\varphi,v)$ with $\varphi:\tau'\to\tau$ and $v:c'\Rightarrow c\varphi$. If $(\varphi',v'):(\tau'',c'')\to(\tau',c')$, their composite is
\[
 (\varphi,v)(\varphi',v')=
 \bigl(\varphi\varphi',(v\varphi')\circ v'\bigr).
\]
At $i\in I_{\tau''}$ the second component is $v_{\varphi'(i)}v'_i$. Threefold composition is associative because both parenthesizations have component $v_{\varphi'\varphi''(i)}v'_{\varphi''(i)}v''_i$ and the same shape map. The identity is $(1_\tau,1_c)$. These checks explain the word ``\emph{\'evidente}'' in the source without suppressing the variance.

The chosen lift of $\varphi:\tau'\to\tau$ at $(\tau,c)$ is
\[
 (\varphi,1_{c\varphi}):(\tau',c\varphi)\longrightarrow(\tau,c).
\]
To verify its Cartesian property, let a map from $(\omega,d)$ to $(\tau,c)$ have shape $\varphi\theta$ and transformation $w:d\Rightarrow c\varphi\theta$. There is exactly one factor over $\theta$ through the displayed lift, namely $(\theta,w)$. The chosen lifts compose on the nose, since $(c\varphi)\theta=c(\varphi\theta)$. More generally $(\varphi,v)$ is Cartesian exactly when $v$ is invertible: it factors as the vertical arrow $(1,v)$ followed by this chosen lift, and a vertical Cartesian arrow is an isomorphism.

For $T=T_{\partial}$ or $T_{\uparrow}$ the endpoint functors send $(\varphi,v)$ to $v_{0_{\tau'}}$ and $v_{1_{\tau'}}$. Their targets are $c(0_\tau)$ and $c(1_\tau)$ because $\varphi$ preserves endpoints. The formula for composition immediately gives functoriality. Likewise a functor $H:X\to Y$ acts by $(\tau,c)\mapsto(\tau,Hc)$ and $(\varphi,v)\mapsto(\varphi,Hv)$, preserving the chosen Cartesian lifts. Restricting the permitted shape maps is literally base change of this split fibration, which proves the fiber-product assertion on page~30.

Size causes no mathematical obstruction here. Fix universes $\mathcal U\in\mathcal V$. The set of finite words and their finite interval functors is $\mathcal U$-small. If $X$ is a $\mathcal V$-small category, each $X^{I_\tau}$ and the total construction are $\mathcal V$-small. One may therefore take $X=\Cat_{\mathcal U}$ after working in a sufficiently larger universe. This supplies the convention which the source deliberately leaves implicit.
\end{dcomment}


% ---- from Der-ChVII-456_p021-025.tex
\printedpage{19}
% head: A. Grothendieck : Dérivateurs, Ch. VII, Édition des Universités de Montpellier II et Paris VII
% foot: Version 22 décembre 2001

\msp{page 29}
Mais comme
\[
\left\{
\begin{array}{l}
\sigma(c)=c(0_\tau),\ \beta(c)=c(1_\tau)\\[1ex]
\text{où } 0_\tau,\ 1_\tau \text{ sont les objets source et but de l'intervalle orienté standard } I_\tau\ ,
\end{array}
\right.
\]
et de même pour $\tau'$, $v$ n'induira bien
\[
\underbrace{\sigma(c')}_{=c'(0_{\tau'})} \longrightarrow c(\varphi(0_{\tau'})) \overset{?}{=} \underbrace{\sigma(c)}_{c(0_\tau)},
\]
et itou pour $\beta(c')$ et $\beta(c)$, que si on impose
\[
\boxed{
\left\{
\begin{array}{lcl}
\varphi(0_{\tau'}) & = & 0_\tau\\
\varphi(1_{\tau'}) & = & 1_\tau,
\end{array}
\right.
}
\]
i.e. que le foncteur envisagé $\varphi$ soit \textit{compatible avec la source et le but.} C'est ce visiblement qu'il fallait imposer dans la définition de % [sic] C'est ce visiblement -> C'est visiblement ce, word order
\[
\mathrm{Hom}_{\ul{Ty}}(\tau',\tau)=\underline{\mathrm{Hom}}^{\circled{int}}(I_{\tau'},I_\tau)\quad\text{-}
\]
il faut prendre les foncteurs $\varphi$ respectant les structures d'intervalle, i.e. rendant commutatif le diagramme
\[
\begin{tikzcd}[column sep=small]
I_{\tau'} \arrow[rr] & & I_\tau \\
& \varepsilon=\{0,1\} \arrow[ul, hook', "i_{\tau'}"] \arrow[ur, hook, "i_\tau"'] &
\end{tikzcd}
\]
définissant l'orientation de $I_{\tau'}$ et de $I_\tau$ (\footnote[10]{En fait, il faut sans doute restreindre encore plus les morphismes en exigeant qu'il soient \textit{croissants} sur les objets - cf. page 31.}). Cette rectification faite, j'ai bon espoir que la catégorie $\ul{Ch}(X)$ qu'on vient de construire est `la bonne'. On pourrait désigner par $\ul{CH}(X)$ la catégorie plus grande, sur la catégorie $\ul{Typ}$ ayant mêmes objets que $\ul{Ty}$, mais plus de morphismes : à savoir % [sic] (in footnote 10) qu'il soient -> qu'ils soient, missing plural s
\msp{page 30}
\[
\begin{array}{lcl}
\mathrm{Hom}_{\ul{Typ}}(\tau,\tau') & = & \mathrm{Hom}_{\mathrm{Cat}}(I_\tau,I_{\tau'})\\[0.5ex]
\mathrm{Hom}_{\ul{Ty}}(\tau,\tau') & = & \mathrm{Hom}^{\mathrm{int}}_{\mathrm{Cat}}(I_\tau,I_{\tau'})\\[0.5ex]
 & & \text{(homomorphismes respectant la structure d'intervalle).}
\end{array}
\]
On aura
\[
\ul{Ch}(X)\simeq\ul{CH}(X)\times_{\ul{Typ}}\ul{Ty},
\]
\printedpage{20}
% head: A. Grothendieck : Dérivateurs, Ch. VII, Édition des Universités de Montpellier II et Paris VII
% foot: Version 22 décembre 2001
i.e. $\ul{Ch}(X)$ est la \textit{restriction} de la catégorie fibrée $\ul{CH}(X)$ sur $\ul{Typ}$, à la sous-catégorie (non-pleine) $\ul{Ty}$ de $\ul{Typ}$.

\dgrblock{dgrb295}{Changing shape need not have a direct image}
\begin{dcomment}{examples}{p28}{Changing shape need not have a direct image}
The projection to the category of types is a fibration because paths can always be restricted. Extending a path to a new shape asks a different question, and generally requires colimits in its coefficient category. This is the usual distinction between a fibration and a bifibration; compare \citet[Expos\'e~VI, \S10 and \S11(f)]{GrothendieckSGA1}.

Consider the interval $J=I_{-+}=(0\leftarrow1\to2)$ and its unique collapse $q:J\to e$. This collapse belongs even to $T_{\uparrow}$. Its inverse-image functor is the diagonal
\[
 q^*:X\longrightarrow X^J,
\]
sending an object to its constant span. If an opCartesian lift over $q$ existed at every span $c$, its target $q_!c$ would represent the functor
\[
 x\longmapsto\operatorname{Nat}(c,\mathrm{const}_x).
\]
Indeed every map from $c$ over $q$ to a constant path would factor uniquely through the opCartesian lift. This is precisely the universal property of the pushout of the span $c(0)\leftarrow c(1)\to c(2)$.

Take $X=J$ itself and $c=1_J$. There is no cocone on this diagram: neither $0$ nor $2$ maps to the other, and neither maps to $1$. Thus $q_!c$ does not exist, and the path projection is not an opfibration in this example. This does not reveal a gap in Grothendieck's assertion, which is only that the construction is split fibrational. It identifies precisely what would be added by a stronger assertion.

For $X=\mathrm{Set}$ the same calculation has a complete answer. Let $A\leftarrow C\to B$ be a span. Then $q_!c$ is the quotient of $A\amalg B$ by the equivalence relation generated by identifying the two images of each $z\in C$. A map from this quotient to $S$ is exactly a pair $A\to S$, $B\to S$ which agrees on $C$, so it is exactly a natural transformation $c\Rightarrow\mathrm{const}_S$. For a general change of shape $\varphi$ and a cocomplete $X$, left Kan extension supplies the corresponding left adjoint to $\varphi^*$. Such extra adjoints come from $X$; the split fibration itself is available without them.
\end{dcomment}


\dgrblock{dgrb296}{How much endpoint preservation is forced}
\begin{dcomment}{development}{p29}{How much endpoint preservation is forced}
The correction on page~29 gives exactly what is needed for endpoint maps obtained by taking components of a natural transformation. There is a slightly larger possibility if one allows those components to be followed by canonical arrows inside the target interval. Describing it clarifies why endpoint preservation is a structural choice, and why concatenation will justify the stricter choice.

\textbf{Endpoint evaluation criterion.}\ours{}
Let $T$ be a subcategory of $T_{\mathrm{all}}$ with all sign words as objects. Its path construction admits functors $\sigma_X,\beta_X:\int F_X\to X$, natural in $X$, restricting in each fiber to evaluation at the specified source and target, if and only if every $\varphi:\tau'\to\tau$ in $T$ admits arrows
\[
 a_\varphi:\varphi(0_{\tau'})\longrightarrow0_\tau,
 \qquad
 b_\varphi:\varphi(1_{\tau'})\longrightarrow1_\tau
 \quad\text{in }I_\tau.
\]
When they exist, these endpoint functors are unique. The permitted functors form a category $T_{\mathrm{ev}}$ containing $T_{\partial}$, and it is the largest such subcategory of $T_{\mathrm{all}}$.

\emph{Proof.} For necessity take $X=I_\tau$ and the universal path $c=1_{I_\tau}$. The Cartesian arrow
\[
 (\varphi,1_\varphi):(\tau',\varphi)\longrightarrow(\tau,1_{I_\tau})
\]
must be sent by $\sigma_X$ to an arrow $\varphi(0_{\tau'})\to0_\tau$, and by $\beta_X$ to an arrow $\varphi(1_{\tau'})\to1_\tau$. Each is unique when it exists, because an interval is thin. For composable maps $\tau''\xrightarrow{\psi}\tau'\xrightarrow{\varphi}\tau$, the arrows
\[
 a_\varphi\circ\varphi(a_\psi),\qquad
 b_\varphi\circ\varphi(b_\psi)
\]
have precisely the endpoints required for $\varphi\psi$. Identities satisfy the requirement as well. Hence the permitted functors form a category.

To construct the endpoint maps, define on $(\varphi,v):c'\to c$
\[
 \sigma_X(\varphi,v)=c(a_\varphi)\circ v_{0_{\tau'}},
 \qquad
 \beta_X(\varphi,v)=c(b_\varphi)\circ v_{1_{\tau'}}.
\]
For a second morphism $(\psi,w):c''\to c'$, naturality of $v$ at $a_\psi$ gives
\[
 v_{0_{\tau'}}c'(a_\psi)
   =c\varphi(a_\psi)v_{\psi(0_{\tau''})}.
\]
After composing with $c(a_\varphi)$ and $w_{0_{\tau''}}$, this says that the product of the two source arrows is
\[
 c(a_{\varphi\psi})
 v_{\psi(0_{\tau''})}w_{0_{\tau''}},
\]
which is the source arrow assigned to the composite. The same argument with $b_\psi$ proves the target identity. Identity morphisms and naturality in $X$ follow directly from the formulas.

Finally, every $(\varphi,v)$ factors as a vertical morphism $(1,v)$ followed by a chosen Cartesian arrow. The endpoint functors are prescribed on vertical arrows. Naturality in $X$ applies to $c:I_\tau\to X$ and forces their values on the Cartesian arrow to be $c(a_\varphi)$ and $c(b_\varphi)$. Thus the formulas are forced, proving uniqueness and maximality.\hfill$\square$

For example, the functor $e\to I_+$ selecting $0$ belongs to $T_{\mathrm{ev}}$: its source comparison is an identity and its target comparison is $0\to1$. It does not preserve both endpoints. The functor selecting $1$ does not belong to $T_{\mathrm{ev}}$, since there is no arrow $1\to0$. This already prevents endpoint evaluation on $T_{\mathrm{all}}$.

The enlargement does not solve the concatenation problem. Concatenate two copies of the admissible map $e\to I_+$ selecting $0$. The common vertex of $e*e=e$ would have to map both to vertex $0$ of the first target interval and to vertex $1$, the starting vertex of the second interval in $I_{++}$. The two prescribed maps therefore do not glue. Strict endpoint preservation makes precisely these images agree. The theorem is an elementary specialization of the general description of functors out of a split fibration in \citet[Expos\'e~VI, Proposition~12.1]{GrothendieckSGA1}. We checked that description and the Yoneda treatment cited above; we did not find this particular maximal class of interval maps stated there. The marker records this explicit deduction and proof, without a claim of priority.
\end{dcomment}


\textbf{NB} Les catégories fibrées $\ul{CH}(X)$, $\ul{Ch}(X)$ sur $\ul{Typ}$, $\ul{Ty}$ dépendent fonctoriellement de $X$, variant dans (Cat), comme il se doit. Comme de juste, on s'y intéresse surtout quand $X$ \textit{n'est pas} dans (Cat), i.e. quand $X$ est une `grosse' catégorie (telle que (Cat) elle-même). Je me dispense d'expliciter ici les questions pertinentes d'Univers.

Deux tests essentiels, pour s'assurer si la notion est bonne
\begin{enumerate}
\item[a)] La théorie des $\mathcal{M}\Sigma^{-1}$.
\item[b)] La composition des chemins doit bien marcher.
\end{enumerate}

\msp{page 31}
On a le diagramme de foncteurs canoniques
\[
\begin{tikzcd}
\ul{Ty} \arrow[r, hook] \arrow[d, hook] & \varepsilon\backslash\mathrm{Cat} \arrow[d] \\
\ul{Typ} \arrow[r, hook] & \mathrm{Cat},
\end{tikzcd}
\]
où les flèches horizontales sont pleinement fidèles.

\dgrblock{dgrb297}{The fraction test and hammock localization}
\begin{dcomment}{status}{p30}{The fraction test and hammock localization}
Grothendieck's two tests concern different levels of the construction. Concatenation should be an operation on actual paths and their morphisms; localization should arise after paths are subjected to specified identifications or homotopical comparisons. The category $\uCh(X)$ constructed here has not yet been asserted to satisfy the universal property of a localization.

For a category $\mathcal M$ and a specified class $\Sigma$ of arrows, ordinary localization is already constructed by \citet[Chapter~I, \S1.1 and Lemma~1.2]{GabrielZisman1967}. Its arrows are composable words of forward arrows and formal inverses of arrows in $\Sigma$, modulo the category relations and inverse relations. A functor out of the quotient is consequently equivalent to a functor from $\mathcal M$ sending $\Sigma$ to isomorphisms. This existence statement does not require the Ore hypotheses needed for a calculus with uniformly short fractions.

The homotopical comparison already available before these notes is quite specific. In \citet[Definition~2.1 and Propositions~2.2 and~3.1]{DwyerKan1980Calculating}, a hammock is a finite grid of zigzags, with backward horizontal arrows and vertical arrows required to be weak equivalences. Identity columns are removed and adjacent columns pointing in the same direction are composed. Their mapping simplicial sets are equivalent to the standard simplicial localization, and their sets of components recover ordinary localization:
\[
 \pi_0L^H(\mathcal M,\Sigma)(x,y)
   \cong\Hom_{\mathcal M[\Sigma^{-1}]}(x,y).
\]
Thus there is a proven construction in which changing the length and subdivision of zigzags is accompanied by the higher comparison data needed for mapping spaces.

This is a comparison of mathematical mechanisms, not an identification of the constructions. The present $\uCh(X)$ allows arbitrary vertical natural transformations, has not specified which backward arrows may be inverted, and uses a single ordinary category of types. The hammock conditions, the fixed endpoints, and the higher simplices must all be supplied before a localization theorem can be invoked. The source's demand that the fraction theory ``work'' is a design requirement at this point, not a research conjecture left open in these pages. Its exact implementation in Grothendieck's subsequent text should be checked on its own definitions.
\end{dcomment}


\printedpage{21}
% head: A. Grothendieck : Dérivateurs, Ch. VII, Édition des Universités de Montpellier II et Paris VII
% foot: Version 22 décembre 2001

\section*{5\quad Propriétés des morphismes d'intervalles}

Je me suis plus ou moins convaincu que pour se prêter à des calculs et à des raisonnements simples, il faut encore plus restreindre les morphismes dans la définition de la catégorie $\ul{Ty}$. Je vais exiger que les morphismes $\tau\longrightarrow\tau'$ soient des foncteurs
\[
w : I_\tau\longrightarrow I_{\tau'}
\]
qui sont compatibles avec la structure d'\textit{intervalles orientés} de $I_\tau$, $I_{\tau'}$, i.e. qu'ils transforment source en source, but en but, mais de plus qu'ils soient \textit{croissants} sur les objets (pour la relation d'ordre de progression des chemins). C'est une restriction qui géométriquement me paraît assez naturelle. (Elle est similaire à la restriction analogue, dans la définition pré-catégorique des morphismes entre $\bDelta_i$ \ldots\,) C'est cette hypothèse qui permet de faire une théorie de structure assez simple des morphismes d'intervalles. Mais pour l'instant je vais distinguer `morphismes' et `morphismes croissants'.

\dgrblock{dgrb298}{The same type category in later zigzag constructions}
\begin{dcomment}{later}{p31:Je vais exiger que les morphismes}{The same type category in later zigzag constructions}
Grothendieck's final restriction on changes of shape has a precise counterpart in the later literature on localization. The correspondence can be checked directly, before making any assertion about homotopy theory. Given a word $\tau$, mark the backward transitions of $I_\tau$ and their composites, together with the identities. This makes $I_\tau$ a relative category. An increasing endpoint-preserving functor automatically respects this marking: a nonconstant image of a backward transition is a composite of backward transitions. Conversely, forgetting the marking of an increasing endpoint-preserving relative functor gives exactly a morphism of $\uTy$. The objects, morphisms, and concatenation are therefore unchanged by this decoration.

Under $+\leftrightarrow A$ and $-\leftrightarrow W^{-1}$, this identifies $\uTy$ with the category $Z$ of \citet[Definition~3.5, Remark~3.6, and footnotes~8--10, pp.~341--342]{MazelGee2018Hammocks}. Their weak-equivalence labels restrict which diagrams in a coefficient category are allowed; they do not impose an additional restriction on these monotone shape maps. The empty word remains the monoidal unit. This is an identification of the finite indexing categories, rather than an inference from a resemblance between pictures.

The image factorization implicit in the manuscript also explains a useful structure. If $f:I_\tau\to I_\upsilon$ is increasing, the full subcategory on its image is an interval. Factor $f$ through its standard oriented model:
\[
 I_\tau\mathrel{\mathop{\longrightarrow}^{q}}I_\omega
 \mathrel{\mathop{\lhook\joinrel\longrightarrow}^{j}}I_\upsilon.
\]
The first map is surjective on vertices and the second injective. Their uniqueness follows from the unique increasing enumeration of the image. Nonidentity surjections lower the length, and nonidentity injections raise it. Thus length is a Reedy degree. The reduced interval gives the smallest image of an endpoint-preserving injection: Proposition~5 says that every such image contains all critical vertices. If the interval is not reduced, its reduced inclusion is consequently initial among the proper injections into it; for a reduced interval there are none. These statements follow from the full-subcategory argument already given, without a new hypothesis on coefficients.

This Reedy structure and this exact use of critical vertices occur in \citet[Example~4.12 and Lemma~4.16]{ArakawaCnossen2026}, where the category is denoted $\mathrm{Zig}_{\partial}$. Their work is an August~2026 preprint. We cite these known applications, rather than present the Reedy interpretation as an original result. What remains to be added to the present manuscript is the relative and homotopical diagram theory over this category of shapes.
\end{dcomment}


\msp{page 32}
Soit
\[
f : I_\tau\longrightarrow I_{\tau'}
\]
un morphisme d'intervalles orientés. Considérons les sous-catégories pleines, isomorphes à des $\Delta_i^+$ ou $\Delta_i^-$ comme intervalles orientés, définis par les intervalles critiques de $I_\tau$. Ils forment une suite $\Sigma_\tau$ de sous-catégories de $I_\tau$
\[
\Sigma_\tau=(\bDelta^{\varepsilon_1}_{l_1}(\tau),\bDelta^{\varepsilon_2}_{l_2}(\tau),\ldots,\bDelta^{\varepsilon_r}_{l_r}(\tau))
\]
dont $I_\tau$ est la somme amalgamée (on reviendra sur ce point plus \texttt{[bas]}), $r$ étant la longueur utile de $\tau$. On les appelle les \textit{sous-intervalles critiques} de $I_\tau$. Ainsi $f$ se décompose en une suite de morphismes
\[
f_i : \bDelta^{\varepsilon_i}_{l_i}\longrightarrow I_{\tau'},\qquad i\in[1,r],
\]
tels que
\begin{enumerate}
\item[a)] pour $i\in[1,r-1]$, le but de $f_i$ soit égal à la source de $f_{i+1}$.

De plus,
\item[b)] $f_1$ transforme source en source, et $f_r$ transforme but en but.
\end{enumerate}

Inversement, quand on se donne des $f_i$ satisfaisant ces deux conditions, ils définissent un morphisme d'intervalles $f$ unique. (Et si la condition b) n'est pas exigée, on trouve une description des morphismes de catégories $I_\tau\longrightarrow I_{\tau'}$.)

\msp{page 33}
On va d'abord regarder un foncteur quelconque
\[
f : \bDelta^{\varepsilon}_i\longrightarrow I_{\tau'}\qquad(\varepsilon\in\{\pm1\})\quad .
\]

\dgrblock{dgrb299}{A functor cannot cross a turn in one run}
\begin{dcomment}{exposition}{p32}{A functor cannot cross a turn in one run}
Grothendieck begins the structure theorem by cutting an interval where its direction changes. Each piece is a chain, and the image of a chain remains comparable even when its length changes. This elementary observation is what allows a map of long zigzags to be described by a finite amount of ordered data.

For a nonempty word $\tau$ let
\[
 0=s_0<s_1<\cdots<s_r=n
\]
be its critical vertices, and let $K_i=[s_{i-1},s_i]$ be its maximal constant-sign runs, regarded as full subcategories. The category $I_\tau$ is obtained by gluing the $K_i$ at their common endpoint vertices. To see the universal property directly, choose functors $f_i:K_i\to Y$ agreeing at shared endpoints. They determine the image of every generating transition of $I_\tau$. Each nonidentity arrow of $I_\tau$ is the unique composite of the transitions along a single run, so its image is forced to be the corresponding composite in $Y$. The only composition equations are associativity along those runs, and these hold in $Y$. Thus the $f_i$ extend uniquely to a functor. For $Y=I_{\upsilon}$, imposing the values of $f_1(s_0)$ and $f_r(s_r)$ additionally imposes endpoint preservation, exactly as on page~32.

Now let $K$ be a critical run of the source and suppose $f|K$ is nonconstant. The image vertices are pairwise comparable in $I_{\upsilon}$: any two source vertices are joined by an arrow in one direction, whose image joins their images. Let $a<b$ be the smallest and largest distinct image vertices in the external order. Their comparability means that all target transitions between $a$ and $b$ have one sign. They therefore lie in a unique maximal run $J'$, and so do all the image vertices between them. This proves both the existence and uniqueness in Proposition~7.

If the source run has sign $\varepsilon$ and $J'$ has sign $\eta$, the object function of $f|K$ is nondecreasing when $\varepsilon=\eta$ and nonincreasing when $\varepsilon=-\eta$. A nonconstant function which is nondecreasing in the external order cannot be of both sorts. Hence an increasing functor sends every nonconstant source run into a target run of the same sign. This proves the sign assertion as well; no endpoint hypothesis has been used.

The preceding paragraph of the source overstates the constant case. A vertex strictly inside a long target run belongs to one critical interval; an interior \emph{critical} vertex belongs to two; an endpoint of a positive-length interval belongs to one. For instance, the constant map with value $1$ in $I_{++}=[2]$ factors through its sole critical interval, although $1$ is not an endpoint. This is a printed slip, retained in the text, and it does not affect Proposition~7 because that proposition assumes nonconstancy. There is also a harmless zero-length exception to the footnote's comparable-subset lemma: $I_e$ has one vertex and no critical runs. A nonconstant map cannot have this target, so the proposition itself remains valid without an added exception.
\end{dcomment}


\printedpage{22}
% head: A. Grothendieck : Dérivateurs, Ch. VII, Édition des Universités de Montpellier II et Paris VII
% foot: Version 22 décembre 2001
Supposons pour fixer les idées que $\varepsilon=+1$, i.e. que le premier membre soit un intervalle critique direct \ldots{} On trouve que $f$ se factorise par un sous-intervalle critique de $I_{\tau'}$, soit $\bDelta^{\varepsilon'}_j$. Il est donc défini par une application croissante de $\bDelta^{\varepsilon}_i$ dans $\bDelta^{\varepsilon'}_j$ (si $\varepsilon$, $\varepsilon'$ \texttt{[sont]} de même signe), une application décroissante dans le sens inverse (\footnote[11]{\textbf{Lemme.} \textit{Si $I_{\tau'}$ est un intervalle, toute partie $J$ de $I_{\tau'}$ telle que pour deux éléments de $J$, il existe toujours une flèche de l'un dans l'autre, est contenue dans un intervalle critique (unique si} $\operatorname{card} J\geq 2$\textit{).}}). Mais si on exige que $f$ soit croissant, alors il faut que \textit{ou bien} $f$ soit constant (auquel cas, si sa valeur n'est pas une extrémité de $I_{\tau'}$, il y a exactement \textit{deux} intervalles critiques de $I_{\tau'}$ par lesquels se factorise $f$) \textit{ou bien} $f$ non constant et $\varepsilon=\varepsilon'$ (et en tous cas, si $f$ non constant, alors l'intervalle critique par lequel se factorise $f$ est unique). Pour résumer :

\textbf{Proposition 7.} \textit{Soient $I_\tau$, $I_{\tau'}$ deux intervalles orientés,}
\[
f : I_\tau\longrightarrow I_{\tau'}
\]
\textit{un foncteur. Alors pour tout intervalle critique $J\subseteq I_\tau$ tel que $f|J$ soit non constant, il existe un unique intervalle critique $J'\subseteq I_{\tau'}$ tel que $f|J$ se factorise par $J'$, i.e. tel que $f(J)\subseteq J'$. Si $f$ est croissant, alors $J$ et $J'$ sont de même signe,}
\msp{page 34}
\textit{i.e. tous deux directs, ou tous deux rétrogrades.}

Ainsi, si on regarde la suite $s_0,s_1,\ldots,s_n$ des sommets critiques de $\tau$, alors
\[
t_0=f(s_0),\ t_1=f(s_1),\ \ldots\ ,\ t_n=f(s_n)
\]
forment une suite de sommets de $I_{\tau'}$, telle que deux consécutifs soient reliés par une flèche de $I_{\tau'}$. Donc la sous-catégorie pleine de $I_{\tau'}$ définie par ces $t_i$ est connexe, c'est donc un sous-intervalle. Si $f$ est croissant, ce sous-intervalle est $[t_0,t_n]$. Si $f$ est un morphisme d'intervalles, il contient tous les sommets critiques de $I_{\tau'}$. Dans tous les cas, prenant $t_i$, $t_j$ $(i<j)$ et appliquant ce qui précède à $f|J$, où $J$ est le sous-intervalle de $I$ d'origine $s_i$, d'extrémité $s_j$, on voit que tout sommet critique de $I_{\tau'}$ compris entre $t_i$ et $t_j$ est parmi les $t_k$ $(i\leq k\leq j)$. Pour résumer :

\textbf{Corollaire.} \textit{Soient $s_i\leq s_j$ deux sommets critiques de $I_{\tau'}$. Alors tout sommet critique de $I_{\tau'}$ compris entre $f(s_i)=t_i$ et $f(s_j)=t_j$, est l'image d'un sommet critique $s_k$ compris entre $s_i$ et $s_j$ (i.e. $i\leq k\leq j$). En particulier, si $f$ est un morphisme d'intervalles,} % [sic] sommets critiques de I_{\tau'} (first occurrence) -> de I_\tau, the s_i are the critical vertices of \tau
\[
f(I_{\tau\,\red})\supseteq I_{\tau'\,\red}
\]
\textit{(où $I_{\tau\,\red}$ désigne l'ensemble des sommets critiques de $I_\tau$ ou aussi le sous-intervalle réduit qu'il définit).}

\dgrblock{dgrb300}{Critical vertices are crossed, not necessarily preserved}
\begin{dcomment}{exposition}{p34:sous-intervalle réduit qu'il définit}{Critical vertices are crossed, not necessarily preserved}
The next step records which turns in the target must be encountered by a map. This is weaker than preservation of all critical vertices, and the distinction is needed when the target is not reduced. It also prevents an ordinary interval of integers from being confused with a connected full subcategory of an interval category.

Let $s_0,\ldots,s_r$ be the critical vertices of the source and put $t_i=f(s_i)$. Consecutive $t_i$ are comparable, since they are images of the endpoints of a source run. Thus the full subcategory on the \emph{set} $\{t_0,\ldots,t_r\}$ is connected. Suppose $t_i$ and $t_j$ lie on opposite sides of an interior critical vertex $z$ of the target. If no $t_k$, $i\leq k\leq j$, were equal to $z$, some consecutive pair in this finite sequence would lie strictly on opposite sides of $z$. Such a pair cannot be comparable: the two transitions adjoining an interior critical vertex have opposite signs. This contradiction proves that $z=t_k$ for some $k$ in the indicated range. The argument permits backtracking and repetitions; it does not require monotonicity.

If $f$ preserves the two endpoints, take $i=0$ and $j=r$. Every target critical vertex is consequently the image of a source critical vertex. This is the content of the corollary, whose first occurrence of $I_{\tau'}$ should be read as $I_\tau$ when it describes the $s_i$. The stronger inclusion in the other direction is false. In the displayed example $I_{+-+}\to I_{++}$, the object function is
\[
 (0,1,2,3)\longmapsto(0,1,1,2).
\]
All four source vertices are critical. Their image contains the two target endpoints and also its noncritical vertex $1$; the middle source transition collapses. The example therefore gives exactly the one-sided containment asserted in the text.

There is a second printed overstatement just before the corollary. Even for increasing $f$, the full subcategory on the $t_i$ need not have \emph{every} integer vertex between $t_0$ and $t_r$. Take the identity of $I_{++}$. Its source critical vertices are $0,2$, so their image is $\{0,2\}$, a full subcategory isomorphic to $[1]$, whereas the ambient interval $[0,2]$ has three vertices. The correct statement is that the image subcategory has least and greatest vertices $t_0,t_r$, its order-convex hull is $[t_0,t_r]$, and it contains every ambient critical vertex in that hull. This is all that the argument uses.

One can now see why the image of $\Cr(\tau,f)\to\Cr(\upsilon)$ is an interval in the ordered set of target runs. Join successive numbers $t_i$ by real line segments, obtaining a continuous walk in $[0,m]$. Each nonconstant segment is contained in the unique target run attached to $K_i$. If the walk reaches two distinct target runs, it must traverse every intervening run, and each such traversal occurs on at least one nonconstant segment. Conversely, a target run not encountered by the walk cannot be attached to any $K_i$. The same argument shows that every target run meeting the segment between $f(0)$ and $f(n)$ in a positive-length segment is encountered. In particular, endpoint preservation makes the run map surjective. This completes the proofs of Corollary~2(a)--(c), including the parts for maps which backtrack.
\end{dcomment}


\msp{page 35}
On fera attention que même si $f$ est croissant, on \textit{n'a pas} en général $f(I_{\tau\,\red})\subseteq I_{\tau'\,\red}$, i.e. $f$ ne transforme pas forcément sommets critiques en sommets critiques, Exemple minimal :
\printedpage{23}
% head: A. Grothendieck : Dérivateurs, Ch. VII, Édition des Universités de Montpellier II et Paris VII
% foot: Version 22 décembre 2001

\noindent
\begin{minipage}[c]{0.6\textwidth}
\[
\begin{tikzcd}
\Phi_3^+\ : & s_0 \arrow[r] \arrow[d, mapsto] & s_1 \arrow[d, mapsto] & s_2 \arrow[l] \arrow[r] \arrow[dl, mapsto] & s_3 \arrow[dl, mapsto] \\
\bDelta_2^+\ : & {} \arrow[r] & {} \arrow[r] & {} &
\end{tikzcd}
\]
\end{minipage}\hfill
\begin{minipage}[c]{0.35\textwidth}
$\Phi_3^+$ est réduite, les deux sommets critiques transformés en le sommet non critique 1.
\end{minipage}

Désignons par $\mathrm{Cr}(\tau)$ l'ensembles des intervalles critiques de $I_\tau$, en correspondance biunivoque (respectant l'ordre total) avec l'ensemble des flèches de transition de $I_{\tau\,\red}$, % [sic] l'ensembles -> l'ensemble, spurious plural s
\[
\mathrm{Cr}(\tau)\simeq\mathrm{Fl}\,\mathrm{tr}\,(I_{\tau\,\red}).
\]
Soit, pour un foncteur
\[
f : I_\tau\longrightarrow I_{\tau'},
\]
$\mathrm{Cr}(\tau,f)$ l'ensemble des intervalles critiques de $I_\tau$ sur lesquels $f$ ne soit pas constant. On a alors une application
\[
\mathrm{Cr}(\tau,f)\longrightarrow\mathrm{Cr}(\tau') \tag*{$(*)$} % label (*) printed at the left margin
\]
induite par $f$ via la proposition 7, associant à tout $J\in\mathrm{Cr}(\tau,f)$ l'unique $J'\in\mathrm{Cr}(\tau')$ tel que $f(J)\subseteq J'$. On a de plus

\textbf{Corollaire 2.}
\begin{enumerate}
\item[\textit{a)}] \textit{L'image de $(*)$ est un segment de l'ensemble totalement ordonné} $\mathrm{Cr}(\tau')$\textit{, i.e. si deux intervalles critiques $J'$, $J''$ dans $I'_{\tau'}$ sont dans l'image, il en est de même de tous} % [sic] I'_{\tau'} -> I_{\tau'}, spurious prime
\msp{page 36}
\textit{ceux qui sont compris entre eux.}
\item[\textit{b)}] \textit{Cette image contient aussi l'intervalle critique de $I_{\tau'}$ (il y en a au plus un) qui contient $f(0_\tau)$ (resp. $f(1_\tau)$) et dont l'intersection avec le segment $[f(0_\tau),f(1_\tau)]$ est de cardinal $\geq 2$, ainsi que tous les intervalles critiques contenus dans ledit segment.}
\item[\textit{c)}] \textit{Si $f$ est un morphisme d'intervalles orientés, alors l'application induite}
\[
\mathrm{Cr}(f):\mathrm{Cr}(\tau,f)\longrightarrow\mathrm{Cr}(\tau')
\]
\textit{est surjective. De façon plus précise, soit $J'\subseteq$} $\mathrm{Cr}(\tau')$\textit{, et soit $J=f^{-1}(J')$. Alors} $\mathrm{Cr}(f)^{-1}(\{J'\})$ \textit{est formé des intervalles critiques dans} $\mathrm{Cr}(\tau,f)$ contenus dans $J$. \textit{On a de plus,} si $f$ est croissant, % [sic] J' \subseteq Cr(\tau') -> J' \in Cr(\tau'), J' is an element of Cr(\tau')
\[
\mathrm{long}(J')=\sum_{K\in\mathrm{Cr}(f)^{-1}(\{J'\})}\pi_{\mathrm{n}}(K,f),
\]
\textit{où $\pi_{\mathrm{n}}(K,f)$ est la `partie' (numérique) de $f|K$, i.e. la longueur de l'intervalle $[f(0_K),f(1_K)]$. (}\textbf{NB} \textit{Si $K$ est un intervalle critique de $I$ tel que $K\notin Cr(\tau,f)$, i.e. $f|K$ est constant, on posera $\pi_{\mathrm{n}}(K,f)=0$. Donc la formule précédente se lit aussi}
\[
\mathrm{long}(J')=\sum\limits_{\underbrace{K\in\mathrm{Cr}(I_\tau),\ K\subseteq J}_{\text{(i.e. }K\,\in\,\mathrm{Cr}(f^{-1}(J'))\text{)}}}\pi_{\mathrm{n}}(K,f).\ \textit{)}
\]
\end{enumerate}

% ---- from Der-ChVII-456_p026-030.tex
\printedpage{24}
% head: A. Grothendieck : Dérivateurs, Ch. VII, Édition des Universités de Montpellier II et Paris VII

\textbf{Corollaire 3.} \textit{Si $f$ est un morphisme} croissant \textit{d'intervalles orientés, alors pour tout $J' \in \Cr(\tau')$, on a}
\[
(0 \leq)\ \card \Cr(f)^{-1}(\{J\}) \longrightarrow \lng(J'),
\] % [sic] \longrightarrow -> \leq, the next sentence speaks of equality; and {J} -> {J'}, as in Cr(f)^{-1}({J'}) below
\msp{page 37}
\textit{l'égalité n'étant obtenue que si les parties de $f$ sur les intervalles critiques de $I_\tau$ qui sont au dessus de $J'$ sont toutes égales à $0$ ou $1$ (donc si et seulement si sur les $J \in \Cr(f)^{-1}(\{J'\})$, elles sont égales à $1$).}

\dgrblock{dgrb301}{The length budget over a target run}
\begin{dcomment}{exposition}{p36}{The length budget over a target run}
Monotonicity turns the preceding coverage statement into an additive formula. A map may pause, but it cannot spend the same portion of the target twice. This is the numerical mechanism behind Corollaries~3--5.

Let $f:I_\tau\to I_\upsilon$ be increasing and preserve endpoints. For a source critical run $K_i=[s_{i-1},s_i]$, write
\[
 d_i=f(s_i)-f(s_{i-1})\geq0.
\]
When $d_i>0$, Proposition~7 associates to $K_i$ the unique target run containing its image; this is $\Cr(f)(K_i)$. Fix a target run $J'=[a,b]$. The nonconstant segments $[f(s_{i-1}),f(s_i)]$ attached to $J'$ cover $[a,b]$, meet only at consecutive endpoints, and are encountered in their natural order. Coverage follows because $f(0)=0$, $f(n)=m$ and every intervening critical vertex is attained; disjointness of interiors follows from the monotonicity of the sequence $f(s_i)$. Their lengths therefore add to $b-a$:
\[
 b-a=\sum_{\substack{K_i\in\Cr(\tau,f)\\\Cr(f)(K_i)=J'}}d_i.
\]
This is the formula for $\pi_{\mathrm n}(K_i,f)$ in the source. It is valid even if $f$ skips noncritical target vertices, because length measures the whole segment spanned by the image, not the number of image vertices.

Every summand displayed is a positive integer. Consequently, if $k_{J'}=|\Cr(f)^{-1}(J')|$, then
\[
 1\leq k_{J'}\leq b-a,
\]
and equality on the right holds precisely when every positive $d_i$ in the sum is $1$. The printed arrow in Corollary~3 is a typesetting slip for $\leq$, and its unprimed $J$ in the singleton should be $J'$. The subsequent discussion of equality gives the intended reading independently of our proof.

For a first calculation take $I_{+-+}\to I_{++}$ with vertex function $(0,1,1,2)$. Its three source runs have lengths of image $1,0,1$; the single target run has length $2$. Thus $k_{J'}=2$ and the upper bound is attained. For the identity $I_{++}\to I_{++}$, the sole source run has image length $2$, so $k_{J'}=1<2$. Both maps are increasing and preserve endpoints. What is counted is the number of nonconstant \emph{runs}, not the number of nonidentity generating arrows.

The alternative summation on page~36 may harmlessly include every global source run contained in $f^{-1}(J')$: the added terms have zero length. Its parenthetical identification of these runs with all critical runs of the preimage is, however, too literal. A preimage can contain a truncated piece of a neighboring global run on which $f$ is constant. For example, take
\[
 f:I_{++--}\longrightarrow I_{+-},\qquad
 (0,1,2,3,4)\longmapsto(0,1,1,1,2).
\]
For $J'=[0,1]$, its preimage has vertices $0,1,2,3$ and critical runs $[0,2]$ and $[2,3]$. Only $[0,2]$ is a full global source run contained in the preimage; $[2,3]$ is a proper part of the global run $[2,4]$. The extra piece is constant with value $1$, hence contributes zero. The equality of sums survives, but equality of the indexing sets does not. One should either keep the global-run indexing or explicitly allow such zero-contribution boundary pieces.
\end{dcomment}


\textbf{Corollaire 4.} \textit{Supposons de plus $I'$} \texttt{[}$= I_{\tau'}$\texttt{]} \textit{réduit. Alors}
\[
\mathit{Cr}(f) : \mathit{Cr}(\tau, f) \longrightarrow \mathit{Cr}(\tau')\ (\simeq \mathrm{Fl}\,\mathrm{tr}(\tau'))
\]
\textit{est} bijective.

Donc dans ce cas, $\tau'$ réduit, on obtient une application en sens inverse
\[
\begin{tikzcd}[column sep=small]
\Tr(\tau') \simeq \Cr(\tau') \arrow[rr, "\Cr^*(f)"] \arrow[dr, "\Cr(f)^{-1}"'] & & \Cr(\tau) \\
 & \Cr(\tau, f), \arrow[ur, hook, "\text{inclusion}"'] &
\end{tikzcd}
\]
application qui transforme l'intervalle critique $J'$ $(= \{i - 1, i\})$ en l'\textit{unique} intervalle critique $J$ de $I_\tau$ tel que
\[
f(J) = J'.
\]

\msp{page 38}
Mais $J \simeq \bD_i^{\varepsilon}$, et on a donc un morphisme \textit{surjectif}
\[
f_0 : \bD_i^{\varepsilon} \longrightarrow J' \simeq \bD_1^{\varepsilon}.
\]
On voit tout de suite qu'il existe alors une \textit{unique} flèche de transition $u$ de $\bD_i^{\varepsilon}$, telle que $f(u) \neq \mathrm{id}$. Donc on trouve le corollaire suivant, plus précis que le corollaire 4 :

\textbf{Corollaire 5.} \textit{Soit $f : I_\tau \longrightarrow I_{\tau'}$ un morphisme croissant d'intervalles orientés, avec $I_{\tau'}$ réduit, et soit $\Tr(\tau, f)$ l'ensemble des flèches de transition $\underline{u}_i$ de $I_\tau$ telle que $f(\underline{u}_i) \neq \mathrm{id}$, i.e. telles que $\pi(u_i, f) = 1$. Alors $f$ induit une application} % [sic] telle -> telles, agreement with "flèches" (cf. "telles que" just after)
\[
\Tr(\tau, f) \longrightarrow \Tr(\tau')
\]
\textit{qui est bijective, et (bien sûr) croissante, donc un isomorphisme d'ensembles ordonnés.}

Ainsi l'application $\Tr(\tau') \longrightarrow \Cr(\tau, f)$ de tantôt se précise par une application bijective
\[
\Tr(f)^* : \Tr(\tau') \longrightarrow \Tr(\tau, f),
\]
reliée à $\Cr(f)^*$ plus haut par le triangle commutatif
\[
\begin{tikzcd}
 & \Tr(\tau, f) \arrow[dd, "\wr\ \alpha"] \\
\Tr(\tau') \arrow[ur, "\Tr(f)^*"] \arrow[dr, "\Cr(f)^*"'] & \\
 & \Cr(\tau, f),
\end{tikzcd}
\]
% foot: Version 22 décembre 2001
\printedpage{25}
% head: A. Grothendieck : Dérivateurs, Ch. VII, Édition des Universités de Montpellier II et Paris VII
\msp{page 39}
l'application $\alpha$ associant à toute $u_i \in \Tr(\tau, f)$ l'unique intervalle critique qui ‘contienne’ $u_i$ (i.e. tel que $u_i$ soit une flèche dudit intervalle, ou encore, fasse partie de la ‘chaîne critique’ associée \dots).

\dgrblock{dgrb302}{A reduced target remembers the surviving transitions}
\begin{dcomment}{examples}{p38}{A reduced target remembers the surviving transitions}
Grothendieck next specializes the length formula to a target with no redundant vertices. Every target run then has length one. There is exactly one nonconstant source critical run over it, and that source run can cross its two target vertices only once. This converts a numerical constraint into the combinatorial invariant used in Proposition~8.

More explicitly, on a source run $K=[s,t]$ mapping onto the two vertices $a,a+1$, the nondecreasing vertex function has the form
\[
 f(s)=\cdots=f(i-1)=a,
 \qquad
 f(i)=\cdots=f(t)=a+1
\]
for exactly one $i$ with $s<i\leq t$. Hence exactly one transition in $K$ has nonidentity image. The run and its target have the same sign, so this transition maps to the target's unique transition in its categorical direction, including when that direction is backward. Performing this construction for all target runs gives the bijection of Corollary~5. It is increasing because the source runs and their crossings occur in the same external order as the target runs. The triangle with $\alpha$ commutes because $\alpha$ forgets which transition was selected and remembers its containing source run.

Take $\tau=++--++$ and $\upsilon=+-+$. Choose transitions $2,4,6$ in the source. The associated vertex map is
\[
 (0,1,2,3,4,5,6)\longmapsto(0,0,1,1,2,2,3).
\]
The first source run maps across the first forward target edge at transition~2; the second maps across the backward target edge at transition~4; the third crosses the final forward edge at transition~6. The other three transitions become identities. If a path of type $\upsilon$ in $X$ is
\[
 x_0\xrightarrow{a}x_1\xleftarrow{b}x_2\xrightarrow{c}x_3,
\]
restriction along this map is the path
\[
 x_0\xrightarrow{1}x_0\xrightarrow{a}x_1
 \xleftarrow{1}x_1\xleftarrow{b}x_2
 \xrightarrow{1}x_2\xrightarrow{c}x_3.
\]
In each of the three source runs either of its two transitions could instead be the surviving one. These choices give $2^3=8$ different maps to this reduced target, all explicitly determined by their three crossing positions. The general theorem on the next pages will allow arbitrary sign words and express this determination as a signed-subsequence bijection.

When the target is $e$, there are no target runs and no nondegenerate transitions. The unique endpoint-preserving map $I_\tau\to I_e$ collapses every transition. Both sides of the bijection in Corollary~5 are empty, so the assertion includes this case. When the source is $e$ and the target has positive length, no endpoint-preserving map exists. These two degenerate cases explain why the later subset description needs the empty word as well as the alternating nonempty words.
\end{dcomment}


On a mieux encore :

\textbf{Proposition 8.} \textit{Soient $\tau, \tau' \in \mathrm{Ty}$, $\tau'$} réduit. \textit{Pour tout morphisme croissant d'intervalles orientés}
\[
f : I_\tau \longrightarrow I_{\tau'}
\]
\textit{(on va dire simplement que $f$ est une flèche dans $\uTy$, et écrire $f \in \Hom_{\uTy}(I_\tau, I'_\tau)$), associons lui la partie $\Tr(\tau, f)$ des flèches de transition de $I_\tau$ non dégénérées pour $f$. Ces flèches forment un ensemble totalement ordonné affecté de signes $+$ ou $-$ (caractère direct ou rétrograde des flèches), et $f$ induit une bijection d'ensembles totalement ordonnés} % [sic] I'_\tau -> I_{\tau'}, the target of f is I_{\tau'}
\[
\underbrace{E_f}_{\subseteq \Tr(\tau)} = \Tr(\tau, f) \simeq \Tr(\tau'),
\]
\textit{et même un isomorphisme d'ensembles totalement ordonnés ‘affectés de signes’. Ainsi, à toute $f$ on associe une partie $E_f$ de $\Tr(\tau)$, et telle que son ‘type’ soit égal à $\tau'$. Désignons par}
\[
\mathfrak{P}_{\tau'}(\Tr(\tau)) \subseteq \mathfrak{P}(\Tr(\tau))
\]
\textit{l'ensemble des parties $E$ de $\Tr(\tau)$ dont le type soit $\tau'$.}

\msp{page 40}
\textit{Alors l'application}
\[
f \longmapsto \Tr(\tau, f) : \Hom_{\uTy}(\tau, \tau') \longrightarrow \mathfrak{P}_{\tau'}(\Tr(\tau))
\]
\textit{est} bijective \textit{: se donner un morphisme $f : \tau \longrightarrow \tau'$ dans $\uTy$,} revient au même \textit{que de se donner une suite strictement croissante de flèches de transition}
\[
\underline{u}_{i_1},\ \underline{u}_{i_2},\ \dots\ ,\underline{u}_{i_{n'}} \qquad (n' = \lng \tau')
\]
\textit{dans $I_\tau$ $(i_1 < i_2 < \cdots < i_{n'})$, telle que la suite des signatures correspondante}
\[
(\varepsilon_{i_1}, \varepsilon_{i_2}, \dots, \varepsilon_{i_{n'}}) = \tau'
\]
\textit{soit égale à $\tau'$.}

\dgrblock{dgrb303}{Choosing the transitions that survive}
\begin{dcomment}{exposition}{p40:se donner une suite strictement croissante}{Choosing the transitions that survive}
The reduced target allows Grothendieck to replace a functor by a selection of edges. This is the first form of the normal form that will make concatenation manageable: each edge of the target has one identifiable ancestor in the source. The hypothesis that the target is reduced is doing exactly the work needed here.

Write the source signature as $\tau=(\varepsilon_1,\ldots,\varepsilon_n)$ and the reduced target as $\tau'=(\varepsilon'_1,\ldots,\varepsilon'_m)$. An increasing endpoint-preserving functor has a nondecreasing object map
\[
 f:[0,n]\longrightarrow[0,m],\qquad f(0)=0,\quad f(n)=m.
\]
The image of the $i$th transition is either an identity or a nonidentity arrow joining $f(i-1)$ to $f(i)$. In a reduced interval every nonidentity arrow is a transition: two successive transitions have opposite signs and cannot compose. Consequently
\[
 f(i)-f(i-1)\in\{0,1\}.
\]
Exactly $m$ differences are $1$. If their indices are $i_1<\cdots<i_m$, preservation of arrow direction gives $\varepsilon_{i_j}=\varepsilon'_j$. Conversely such a signed subsequence defines
\[
 f(k)=\#\{j\mid i_j\leq k\}.
\]
On a selected transition this map uses the corresponding target transition; on every other transition it uses an identity. These choices extend uniquely to a functor because the source is the free category on its oriented linear graph. They preserve the two endpoints and recover the selected subsequence. This proves Proposition~8, including its empty-target case: the empty subsequence gives the unique functor to the point.

For example, from the word $++-+$ to the reduced word $+-+$ there are exactly two morphisms. The negative transition must be the third source edge, the last positive transition must be the fourth, and the first positive transition may be either the first or the second. Their object maps are respectively
\[
 (0,1,1,2,3),\qquad (0,0,1,2,3).
\]
The two paths retain different occurrences of the first positive step. Thus the type records positions, not merely the numbers of positive and negative transitions.

For a single-edge target the same argument gives the stated representability of positive and negative transitions. A positive edge at position $i$ determines the map to $\Delta_+$ which is $0$ before that edge and $1$ after it; a negative edge determines the analogous map to $\Delta_-$. Naturality remains to be explained because an edge can be stretched over several target edges. The next construction supplies exactly that explanation. When the target is not reduced, the selection alone is insufficient: the unique morphism $\Delta_+\to\Delta_+^2$ stretches one edge across two, so the size of each surviving image must also be retained.
\end{dcomment}


\textbf{Corollaire 1.} \textit{Pour tout $\tau$, on a des bijections canoniques}
\[
\left\{
\begin{array}{lcl}
\Tr^+(\tau) & \simeq & \Hom_{\uTy}(\tau, \Delta^+) \\
\Tr^-(\tau) & \simeq & \Hom_{\uTy}(\tau, \Delta^-),
\end{array}
\right.
\]
fonctorielles en $\tau$ \textit{(variant dans} $\uTy$\textit{).}
% foot: Version 22 décembre 2001

\printedpage{26}
% head: A. Grothendieck : Dérivateurs, Ch. VII, Édition des Universités de Montpellier II et Paris VII

Pour donner un sens au corollaire, il faut définir
\[
\tau \longmapsto \Tr^+(\tau), \qquad \tau \longmapsto \Tr^-(\tau)
\]
comme \textit{foncteurs} contravariant en $\tau$, donc pour un morphisme
\[
f : I_{\tau_2} \longrightarrow I_{\tau_1}
\]
dans $\uTy$, définir des applications
\msp{page 41}
\[
\begin{array}{rcl}
\Tr^+(f)^* : \Tr^+(\tau_1) & \longrightarrow & \Tr^+(\tau_2) \\
\Tr^-(f)^* : \Tr^-(\tau_1) & \longrightarrow & \Tr^-(\tau_2),
\end{array}
\]
en d'autres termes, une application
\[
\Tr(f)^* : \Tr(\tau_1) \quad \longrightarrow \quad \Tr(\tau_2)
\]
‘compatible avec les signatures’. Dans le corollaire 5 plus haut, on a défini une telle application seulement pour $\tau'$ réduit \texttt{[}$\tau_1 = \tau'$\texttt{,} $\tau_2 = \tau$\texttt{]} (\footnote[12]{Ce qui suit ne fait pas intervenir les signatures - c'est un sorite sur les applications croissantes entre ensembles finis totalement ordonnés.}). Pour une flèche $u$ de $I_\tau$, appelons \textit{partie ensembliste} de $f$ sur $u$, où de $u$ pour $f$, la partie ensembliste de la flèche $f(u)$ dans $I_{\tau'}$, i.e. le segment $[f(i-1), f(i)] \subseteq I_{\tau'}$ (ici $u = \underline{u}_i$). Soit % [sic] où -> ou, alternative "de f sur u, ou de u pour f"
\[
\underline{u}_{i_1},\ \dots\ ,\underline{u}_{i_N} \qquad \text{ou } (v_1, v_2, \dots, v_n)
\] % [sic] v_n -> v_N, the sequence has N elements (cf. v_N below)
la suite croissante des éléments de $\Tr(\tau, f)$, on voit alors par récurrence sur $i$ que pour $i \in [1, N-1]$, on a
\[
\begin{array}{rcl}
f(\mathrm{but}(v_i)) & = & f(\mathrm{source}(v_{i+1}))\ \ (= t_i), \\
f(\mathrm{source}(v_1)) & = & 0_{\tau'}, \\
f(\mathrm{but}(v_N)) & = & 1_{\tau'}.
\end{array}
\]
Considérons la suite
\[
\{\overbrace{0_{\tau'}}^{= t_0}, t_1, t_2, \dots, t_{N-1}, \overbrace{1_{\tau'}}^{= t_N}\},
\]
elle est définie par le sous-ensemble correspondant de $I_{\tau'}$, qu'on appelera l'ensemble des \textit{sommets critiques de $I_{\tau'}$ pour $f$.} C'est une suite contenant $0_{\tau'}$, $1_{\tau'}$, et telle que deux éléments consécutifs sont reliés dans $I_{\tau'}$ par une flèche, ce qui équivaut à dire qu'elle contient les sommets critiques (intrinsèques) de $I_{\tau'}$. D'autre part, pour tout intervalle de transition $[i-1, i]$ de $I_{\tau'}$ $(i \in [1, n'])$, % [sic] appelera -> appellera, French spelling
% foot: Version 22 décembre 2001
\printedpage{27}
% head: A. Grothendieck : Dérivateurs, Ch. VII, Édition des Universités de Montpellier II et Paris VII
\msp{page 42}
il existe un et un seul intervalle de transition $[t_{j-1}, t_j]$ pour cette suite $(j \in [1, N])$, telle qu'on ait $[i-1, i] = [t_{j-1}, t_j]$. C'est dire qu'il existe une et une seule $v_j \in \Tr(\tau, f)$ telle que la partie de $f$ sur $v_j$ (savoir justement $[j-1, j]$) contienne (celle de) $\underline{u}_i$. Pour résumer :

\textbf{Corollaire 2.} \textit{Soit $f : \tau \longrightarrow \tau'$ dans $\uTy$. Alors pour toute $u \in \Tr(\tau') = \Tr(I_{\tau'})$, il existe une et une seule $v \in \Tr(\tau)$ telle que la partie de $f$ sur $v$ contienne $u$. De plus, les signatures de $u$ et de $v$ sont égales.}

C'est ainsi qu'on obtient donc une application
\[
\Tr(f)^* : \Tr(\tau') \quad \longrightarrow \quad \Tr(\tau) \qquad (\text{se factorisant par } \Tr(\tau, f))
\]
induisant
\[
\begin{array}{rcl}
\Tr^+(f)^* : \Tr^+(\tau') & \longrightarrow & \Tr^+(\tau) \\
\Tr^-(f)^* : \Tr^-(\tau') & \longrightarrow & \Tr^-(\tau).
\end{array}
\]
La transitivité pour un composé
\[
\tau \xrightarrow{f} \tau' \xrightarrow{g} \tau''
\]
ou
\[
I_\tau \xrightarrow{f} I_{\tau'} \xrightarrow{g} I_{\tau''}
\]
est immédiate. Si $u$, $u'$, $u''$ sont des flèches de transition dans $I_\tau$, $I_{\tau'}$, $I_{\tau''}$ respectivement, telles que
\[
\begin{array}{rcll}
\pi(g, u') & \supseteq & \pi(u''), & \qquad \text{i.e. } u' = \Tr(g)^*(u'') \\
\pi(f, u) & \supseteq & \pi(u'), & \qquad \text{i.e. } u = \Tr(f)^*(u'),
\end{array}
\]
alors on a
\[
\pi(gf, u) \ \supseteq\ \pi(u''), \qquad \text{i.e. } u = \Tr(gf)^*(u'').
\]

\msp{page 43}
Si les flèches en question ont comme partie intrinsèque $[i-1, i]$, $[i'-1, i']$, $[i''-1, i'']$ respectivement, l'hypothèse est
\[
\begin{array}{rclrcl}
g(i'-1) & \leq & i''-1, & g(i') & \geq & i'' \\
f(i-1) & \leq & i'-1, & f(i) & \geq & i'
\end{array}
\]
et comme $g$ est croissant, on a déduit
\[
\begin{array}{rcll}
gf(i-1) & \leq & g(i'-1) & (\leq i''-1) \\
gf(i) & \geq & g(i') & (\geq i''),
\end{array}
\]
\begin{flushright}
c.q.f.d.
\end{flushright}
% foot: Version 22 décembre 2001

\dgrblock{dgrb304}{A stretched edge and the containment convention}
\begin{dcomment}{examples}{p42}{A stretched edge and the containment convention}
Grothendieck is assigning every target transition to the unique source transition whose image crosses it. One small example fixes both the meaning of this assignment and two printed slips immediately before Corollary~2.

Let $f:\Delta_+\to\Delta_+^2$ have object map $f(0)=0$, $f(1)=2$. There is one surviving source transition, so $N=1$, $t_0=0$ and $t_1=2$. The two target transitions have underlying vertex segments $[0,1]$ and $[1,2]$. Both are contained in $[t_0,t_1]=[0,2]$, and neither equals it. Thus the printed relation
\[
 [i-1,i]=[t_{j-1},t_j]
\]
must be read mathematically as $[i-1,i]\subseteq[t_{j-1},t_j]$. The parenthesis identifying the part of $f$ on $v_j$ as $[j-1,j]$ has the same problem: that part is $[t_{j-1},t_j]$. Both readings occur in the edition; they are retained in the source and repaired here. Corollary~2 itself uses containment and is correct.

The intervals $[t_{j-1},t_j]$ may meet at a vertex, but their sets of transitions are disjoint. Since $t_0=0$, $t_N=m$ and the $t_j$ are strictly increasing, these transition sets partition $\{1,\ldots,m\}$. A transition cannot lie in two of them: its two endpoints would then have to belong to their intersection, which contains at most one vertex. This proves uniqueness without assuming that the target is reduced. Functoriality forces every edge in the block assigned to $v_j$ to have the sign of $v_j$.

The words ``source'' and ``but'' in the preceding display concerning the $v_j$ must here be read as the first and last vertices in the traversal order. On a retrograde transition they are the reverse of its categorical domain and codomain. Writing $f(i_j-1)$ and $f(i_j)$ avoids this ambiguity. If the target is the point there are no target transitions and the pullback map of transition sets is the unique map from the empty set; no first or last surviving transition should be invoked.
\end{dcomment}


\printedpage{28}
% head: A. Grothendieck : Dérivateurs, Ch. VII, Édition des Universités de Montpellier II et Paris VII

Ainsi, le corollaire 1 plus haut obtient un sens précis et il est démontré du même coup, vu que les correspondances du corollaire 1 sont données par
\[
\underbrace{f}_{\in \Hom_{\uTy}(\tau, \Delta^{\varepsilon})} \qquad \longmapsto \Tr(f)^*(\underline{u}),
\]
où $\underline{u}$ est l'unique flèche de transition de $\Delta^{\varepsilon}$.

\dgrblock{dgrb305}{Edges reverse the direction of a map}
\begin{dcomment}{exposition}{p43:Ainsi, le corollaire 1 plus haut}{Edges reverse the direction of a map}
The contravariant transition map contains more information than its notation initially suggests. It determines the whole functor, and gives a particularly economical modern description of the final category of types. This observation explains why the later calculations require only ordered lists and nonnegative integers.

Let $\mathcal O_{\pm}$ have as objects finite linearly ordered sets equipped with a sign at each element, including the empty set, and as arrows the nondecreasing maps which preserve signs. We use the skeletal version with underlying sets $[1,n]$. For $f:\tau\to\tau'$ in $\uTy$, where $\tau$ and $\tau'$ have lengths $n$ and $m$, define
\[
 q_f:[1,m]\longrightarrow[1,n],\qquad
 q_f(j)=\min\{i\mid j\leq f(i)\}.
\]
The endpoint conditions ensure that this minimum exists when $m>0$; its defining inequalities are
\[
 f(q_f(j)-1)<j\leq f(q_f(j)).
\]
This is Grothendieck's $\Tr(f)^*$. It is nondecreasing, and it preserves signs because the image of a source transition is a monochromatic string of target transitions of that same sign.

Conversely, suppose $q:[1,m]\to[1,n]$ is nondecreasing and sign-preserving. Put
\[
 F_q(i)=\#q^{-1}([1,i]),\qquad 0\leq i\leq n.
\]
The fibre $q^{-1}(i)$ is a consecutive block, possibly empty, whose target signs are all $\varepsilon_i$. Hence the consecutive vertices $F_q(i-1)$ and $F_q(i)$ are joined in the target by an arrow of the required direction, or coincide when the fibre is empty. Sending the $i$th source transition to this arrow defines a unique functor. Its object map is nondecreasing, begins at $0$, and ends at $m$. The inequalities above show $q_{F_q}=q$, while counting the target edges crossed by the first $i$ source edges gives $F_{q_f}(i)=f(i)$. These constructions are inverse, including the empty cases. There is no map from the point to a nonpoint interval in $\uTy$, just as there is no function from a nonempty edge set to the empty one.

If $g:\tau'\to\tau''$, a target edge of $\tau''$ has a unique ancestor in $\tau'$ and then a unique ancestor in $\tau$. Equivalently, applying the displayed inequalities twice gives
\[
 q_{gf}=q_fq_g.
\]
Identities induce identities. We therefore have an isomorphism of skeletal categories
\[
 \uTy\ \cong\ \mathcal O_{\pm}^{\mathrm{op}}.
\]
The reversal of variance is essential: a functor stretches a source edge into a target block, whereas $q_f$ sends each member of that block back to its source edge. The signs can be forgotten for the construction on ordered sets, as Grothendieck observes, but must be restored to say when the resulting vertex map is a functor between the oriented interval categories.

This identification also establishes the naturality promised in Corollary~1. If an edge of sign $\varepsilon$ represents $h:\tau'\to\Delta_\varepsilon$, the edge representing $hf$ is $q_f$ of the edge representing $h$. It is precisely ordinary precomposition in the representable presheaf $\Hom_{\uTy}(-,\Delta_\varepsilon)$. The underlying single-sign calculation is the familiar ordinal description of the simplicial category; see \citet[Appendix~D.1, pp.~713--715]{AguiarMahajan2010}. The signed formulation follows by retaining the labels throughout the inverse constructions above.
\end{dcomment}


On va utiliser ces sorites pour donner, pour tout $\tau \in \uTy$, une description combinatoire simple des catégories $\uTy/\tau$ et $\tau\backslash\uTy$, ou du moins de leurs ensembles d'objets.

\textbf{1$^o$ Description de} $\Ob\,(\tau\backslash\uTy)$\textbf{.}

À $f \in \Ob(\tau\backslash\uTy)$, i.e.
\[
f : \tau \longrightarrow \tau' \qquad \text{i.e. } f : I_\tau \longrightarrow I_{\tau'}
\]
dans $\uTy$, est associé $E = E_f = \Tr(\tau, f) \subseteq \Tr(\tau)$, l'ensemble des flèches de transition critiques (i.e. non dégénérées) pour $f$. À chaque $u \in E$, associons l'entier $\lng(\pi(f, u))$,
\msp{page 44}
la \textit{partie} de $f$ sur $u$. Alors $(\tau', f)$ est déterminé de façon unique par le couple
\[
(E, \pi), \qquad E \in \mathfrak{P}(\Tr(\tau)), \qquad \underbrace{\pi}_{\text{pondération}} : E \ \longrightarrow\ \underbrace{\mathbf{N}^*}_{\substack{= \text{ensemble des} \\ \text{entiers} \geq 1}} \quad .
\]
Ou si on préfère, par l'application $\pi = \pi_f$
\[
\Tr(\tau) \xrightarrow{\pi} \mathbf{N}, \qquad u \longmapsto \pi(u) = \pi(u, f) = \lng(\pi_{\tau'}(f(u)))
\]
(qui redonne $E$ comme
\[
E = \pi^{-1}(\mathbf{N}^*) \qquad ).
\]
Cette application n'est soumise à aucune restriction,
\[
\Ob\,(\tau\backslash\uTy) \ \xrightarrow[\sim]{\text{bij.}}\ \Hom_{\mathrm{Ens}}(\Tr(\tau), \mathbf{N}).
\]
On récupère $\tau'$ à partir de $\pi$ ainsi. On prend $E = \pi^{-1}(\mathbf{N}^*)$ (flèches de transition non dégénérées pour $\pi$), $E$ est totalement ordonné par l'ordre induit par celui de $\Tr(\tau)$, il est de plus signaturé \texttt{[sic]} par la signaturation \texttt{[sic]} de $\Tr(\tau)$, soit
\[
(\varepsilon_1, n_1),\ \dots\ ,\ (\varepsilon_r, n_r)
\]
la suite des signatures et pondérations, de la suite des éléments de $E$ $(\varepsilon_i \in \{\pm 1\}$, $n_i \in \mathbf{N}^*$ pour $i \in [1, r])$. Alors,
\[
\tau' = \Delta_{\varepsilon_1}^{n_1} \Delta_{\varepsilon_2}^{n_2} \cdots \Delta_{\varepsilon_r}^{n_r}, \qquad r = \card E,
\]
dans le monoïde libre $\mathrm{Ty}$, où on a désigné par $\Delta_+$, $\Delta_-$ (au lieu de $\Delta^+$, $\Delta^-$) les deux générateurs. L'application
\[
I_\tau \ \longrightarrow\ I_{\tau'}
\]
% foot: Version 22 décembre 2001

% ---- from Der-ChVII-456_p031-035.tex
\printedpage{29}
% head: A. Grothendieck : Dérivateurs, Ch. VII, Édition des Universités de Montpellier II et Paris VII
\msp{page 45}est décrite de façon unique en disant qu'elle est \textit{constante} sur chacune des `composantes connexes' (au sens de la numérotation des sommets de $I_\tau$) de l'ensemble $E' = E'(f) \subseteq I_\tau = [0,n]$ des sommets qui sont sources ou buts de flèches de transition dégénérées pour $f$, i.e. dans $\mathrm{Tr}(\tau)\backslash E$. Et sur le $\alpha$-ème élément de $E$ ($\alpha \in [1, r = \operatorname{card} E]$), d'indice $i_\alpha$ dans $\mathrm{Tr}(\tau)$, $f$ induit le foncteur
\[
\boldsymbol{\Delta}_{\varepsilon_\alpha} \simeq \underbrace{J_{i_\alpha}}_{\substack{\text{intervalle}\\ \text{de transition}\\ \text{d'indice } i_\alpha\\ \text{dans } I_\tau}} \longrightarrow I_{\tau'}
\]
donné par \textit{l'unique} \underline{Ty}-\textit{flèche} dans $I_\tau$ de $J_{i_\alpha}$ dans le $\alpha$-ème facteur, de type $\Delta^{n_\alpha}_{\varepsilon_\alpha}$, de $\tau'$, donné par $(*)$
\[
\Delta^{n_\alpha}_{\varepsilon_\alpha} \simeq \text{type de } \boldsymbol{\Delta}^{\varepsilon_\alpha}_{n_\alpha}\ \ \ldots
\]

\dgrblock{dgrb306}{Weight vectors and composition by block sums}
\begin{dcomment}{examples}{p44}{Weight vectors and composition by block sums}
The outgoing classification now has coordinates which are useful for actual calculations. For a fixed source $\tau=(\varepsilon_1,\ldots,\varepsilon_n)$, put
\[
 a_i=f(i)-f(i-1)=\#q_f^{-1}(i),\qquad
 A_i=\sum_{h=1}^{i}a_h,\quad A_0=0.
\]
Every $a_i$ is a nonnegative integer. The target is the word
\[
 \tau'=\varepsilon_1^{a_1}\varepsilon_2^{a_2}\cdots\varepsilon_n^{a_n},
 \qquad f(i)=A_i,
\]
where exponent zero means omission. These equations prove the claimed unrestricted parametrization of $\Ob(\tau\backslash\uTy)$: they construct an outgoing morphism from every weight vector and recover the vector from every outgoing morphism. When adjacent source signs agree, their target blocks may merge into one critical interval, but their individual lengths remain part of the morphism.

Take the source $+-++-$ and the weights $(2,0,1,3,0)$. Its target is $++++++$, and its vertex map is
\[
 (0,2,2,3,6,6).
\]
The pullback on target transitions is $(1,1,3,4,4,4)$. Both negative source edges have been collapsed. The two consecutive positive source edges at positions $3$ and $4$ survive separately, even though their target images lie in the same positive run.

To see composition, give the six target edges the further weights
\[
 (b_1,\ldots,b_6)=(1,0,2,1,0,3).
\]
The second map has vertex values $(0,1,1,3,4,4,7)$ and target $+^7$. The composite has weights $(1,0,2,4,0)$, since its $i$th weight is
\[
 c_i=\sum_{j=A_{i-1}+1}^{A_i}b_j.
\]
Thus its vertex map is $(0,1,1,3,7,7)$, which is also obtained by direct substitution into the two displayed object maps. An empty block contributes zero; a block of length greater than one contributes a sum. This is the numerical content of composition in Section~6.

The formula gives a useful direct verification of associativity. Given a third weight vector, first replacing every $b_j$ by the sum over its descendants and then summing over the $a_i$-block is the same finite sum as first merging the $b_j$-blocks and then summing over the merged block. The order and boundaries of the summands are unchanged. No cancellation, choice of representative, or equivalence relation enters this calculation.
\end{dcomment}


\subsection*{2$^o$ Détermination de $\mathrm{Ob}\,(\uTy/\tau')$.}

Si $(\tau, f) \in \mathrm{Ob}\,(\uTy/\tau')$, donc
\[
f : I_\tau \longrightarrow I_{\tau'},
\]
on lui associe tout d'abord la partie
\[
\mathrm{Cr}(\tau', f) = E'_f = E' \subseteq I_{\tau'}
\]
formée des sommets (dits `critiques pour $f$', cf. page 41) de $I_{\tau'}$ qui sont source ou but d'une flèche $f(u)$, où $u \in \mathrm{Tr}(\tau, f)$, i.e. $f(u)$ n'est pas dégénérée (\footnote[13]{dégager un lemme $f$ croissante inutile.}). $E'$ contient $f(I_{\tau\,\text{réd}})$ (i.e. les images par $f$ des sommets critiques de $I_\tau$),
\msp{page 46}mais peut être plus grand. (P.ex. si $\tau = \tau'$, $f = \mathrm{id}$, $E' = I_{\tau'}$, qui est donc strictement plus grand que $I_{\tau'\,\text{réd}}$ si $\tau'$ n'est pas réduit.) Pour tout $i \in E'$, considérons $f^{-1}(\{i\}) \subseteq I_\tau$. Comme $f$ \texttt{[est]} croissant, c'est un intervalle de l'ensemble totalement ordonné $I_\tau$, donc c'est un sous-intervalle de la catégorie intervalle $I_\tau$. Comme tel, il a une signature $\tau_i \in \Ty$. On a donc une application
\[
\eta_f \text{ ou } \eta : E' \longrightarrow \Ty.
\]
Je dis que $(\tau, f) \in \mathrm{Ob}\,(\uTy/\tau')$ est déterminé de façon unique par ce couple
\[
(E', \eta) \begin{cases} E' \in \mathfrak{P}(I_{\tau'}) = \mathfrak{P}([0,n']) \\[1ex] \eta : E' \longrightarrow \Ty\ . \end{cases}
\]
De plus, celui-ci est soumis à la seule restriction que
\[
E' \supseteq \mathrm{Cr}(\tau') = \text{ensemble des sommets critiques de } I_{\tau'}
\]
\printedpage{30}
% head: A. Grothendieck : Dérivateurs, Ch. VII, Édition des Universités de Montpellier II et Paris VII
(où ce qui revient au même, que $E'$ contient $0_{\tau'}$, $1_{\tau'}$ et que deux éléments consécutifs de $E'$ sont reliés par une flèche de $I_{\tau'}$) : % [sic] où -> ou, « ou ce qui revient au même »
\[
\left\{
\begin{array}{lcl}
\mathrm{Ob}\,(\uTy/\tau') & \simeq & \text{ensemble des couples } (E',\eta), \text{ avec} \\[1ex]
 & & E' \subseteq [0, \underbrace{\ n'\ }_{=\,\mathrm{long}\,\tau'}], \quad E' \supseteq \mathrm{Cr}(\tau'), \\[1ex]
 & & \eta : E' \longrightarrow \Ty\ .
\end{array}
\right.
\]
Soit
\[
0 = t_0 < t_1 < \ \cdots\ < t_r = n'
\]
la suite des éléments de $E'$, soit pour \ovalbox{$i \in [1,N]$}
\[
\left\|
\begin{array}{l}
l_i = t_i - t_{i-1} \\[1ex]
\varepsilon_i \quad \text{signature de } [t_{i-1}, t_i] \text{ dans } I_{\tau'}\ ,
\end{array}
\right.
\]
\msp{page 47}et pour \ovalbox{$i \in [0,N]$}\,, soit
\[
\|\ \ \tau_i = \eta(t_i).
\]
Alors le type $\tau$ est donné par
\[
\tau = \tau_0 \Delta_{\varepsilon_1} \tau_1 \Delta_{\varepsilon_2} \cdots \Delta_{\varepsilon_N} \tau_N, \tag*{$(*)$}
\]
\texttt{[de]} longueur $N + \sum_{i=0}^N \underbrace{\mathrm{long}(\tau_i)}_{\lambda_i \geq 0}$.

\dgrblock{dgrb307}{Fibres are components of the collapsed-edge graph}
\begin{dcomment}{exposition}{p45}{Fibres are components of the collapsed-edge graph}
Grothendieck wants a geometric reconstruction of the same weight vector. The intended reconstruction is correct once the graph of collapsed transitions is kept, but the printed description in terms of a set of vertices loses information that is needed to find its components.

Consider $\tau=+++$ with weights $(0,1,0)$. The corresponding map to $+$ is
\[
 f(0)=f(1)=0,\qquad f(2)=f(3)=1.
\]
The vertices incident to a collapsed transition form the whole set $\{0,1,2,3\}$. As a subset of the integers this set is connected. Consequently the assertion that $f$ is constant on its integer-connected components would make $f$ constant everywhere, contrary to the displayed functor. This is a printed mathematical slip, not a defective proof of the weight classification.

Instead, take the graph with all source vertices and only those edges $[i-1,i]$ for which $a_i=0$. Its components are consecutive strings of zero-weight edges, together with any isolated vertices. In the example they are $\{0,1\}$ and $\{2,3\}$. Since
\[
 f(j)-f(i)=\sum_{h=i+1}^{j}a_h\qquad(i\leq j),
\]
nonnegativity shows that two vertices have the same image if and only if every intervening edge has weight zero. The graph components are therefore exactly the nonempty vertex fibres of $f$. On each surviving transition the endpoint values are determined by the cumulative sum of the earlier weights. This proves both existence and uniqueness of the proposed geometric reconstruction after the correction.

The same observation produces a canonical factorization. First collapse the zero-weight edges, leaving the signed subword consisting of the active source edges; then stretch its successive edges by their positive weights. The first vertex map is surjective, the second strictly increasing, and both preserve the endpoints. Uniqueness follows either from the fibre components or from the unique image factorization of $q_f$ in the category of signed ordered sets. This factorization will let us distinguish a combinatorial contraction from a contraction which has an inverse on one side.
\end{dcomment}


\subsection*{3$^o$ Description de $\mathrm{Fl}(\uTy)$.}

Si
\[
f : \tau \longrightarrow \tau' \qquad \texttt{[est]} \ \text{dans } \uTy\ ,
\]
on lui associe comme tantôt un entier \underline{$N \geq 0$} ($N = \operatorname{card}\mathrm{Cr}(\tau', f){-}1$, $N = 0$ si et seulement si $\tau' = 1$, type de l'intervalle ponctuel), des applications
\[
\left\{
\begin{array}{lcl@{\qquad\qquad}ccl}
{[1,N]} & \longrightarrow & [+1,-1] \times \mathbf{N}^*, & i & \longmapsto & (\varepsilon_i, l_i) \\[1ex]
{[0,N]} & \longrightarrow & \Ty & i & \longmapsto & \tau_i
\end{array}
\right.
\]
(\footnote[14]{\textbf{NB} Si $N = 0$, $[1,N] = \emptyset$, et il n'y a pas de $\varepsilon_i$, $l_i$; l'application $\underbrace{[0,N]}_{=\,\{0\}} \longrightarrow \Ty$ se réduit à la donnée de $\tau_0 \in \Ty$.}). La flèche $f$ dans $\uTy$ est determinée par ces données. Le type $\tau$ est donné par la formule $(*)$ ci-dessus, le type $\tau'$ par % [sic] determinée -> déterminée, accent manquant
\[
\tau' = \Delta^{l_1}_{\varepsilon_1} \Delta^{l_2}_{\varepsilon_2}\ \cdots\ \Delta^{l_N}_{\varepsilon_N}, \qquad \text{longueur } n' = \textstyle\sum l_i, \tag*{$(*')$}
\]
\printedpage{31}
% head: A. Grothendieck : Dérivateurs, Ch. VII, Édition des Universités de Montpellier II et Paris VII
et l'application croissante de $I_\tau = [0, N + \sum \mathrm{long}(\tau_i)]$ dans $I_{\tau'} = [0, \sum l_i]$ est donnée ainsi :
\[
\begin{array}{l@{\qquad}cl}
\text{constante sur } [0, 0+\lambda_0] & \longmapsto & 0 \\[1ex]
\text{constante sur } [\lambda_0+1, \lambda_0+1+\lambda_1] & \longmapsto & l_1 \\[1ex]
\text{constante sur } [\lambda_0+\lambda_1+2, \lambda_0+\lambda_1+2+\lambda_2] & \longmapsto & l_2 \\[1ex]
\ldots & & \\[1ex]
\text{constante sur } [\lambda_0+\lambda_1+\cdots+\lambda_i+(i+1), \lambda_0+\cdots+\lambda_i+\lambda_{i+1}+(i+1)] & \longmapsto & l_i.
\end{array}
\]

\printedpage{32}
% head: A. Grothendieck : Dérivateurs, Ch. VII, Édition des Universités de Montpellier II et Paris VII
\msp{page 48}

\dgrblock{dgrb308}{Reading the morphism from its target}
\begin{dcomment}{exposition}{p47}{Reading the morphism from its target}
The incoming classification organizes a morphism by the places where its nonconstant pieces meet. It is the same normal form viewed from the target: choose the visited vertices, and at each visited vertex record the entire interval that has been collapsed there. This description is useful because the intervening noncollapsed transitions are then forced.

Let $f:I_\tau\to I_{\tau'}$ be increasing and endpoint-preserving, with target length $m$. The uniform definition is
\[
 E'=\operatorname{im}(f)\subseteq[0,m].
\]
When $m>0$, it agrees with the printed definition by endpoints of nondegenerate images. Indeed, take a vertex value $t$ assumed by $f$. If not every source vertex has value $t$, connectedness of the underlying source path supplies an edge from the nonempty set $f^{-1}(t)$ to its complement. That edge is nondegenerate and one of its image endpoints is $t$. Conversely every such endpoint is a value of $f$. This argument does not require monotonicity, and supplies the elementary lemma requested in footnote~13, provided the functor is nonconstant.

For $m=0$ all transition images are identities, so the printed endpoint definition would instead give $E'=\varnothing$. The later claim that $E'$ contains the endpoint, and the definition $N=\#E'-1$, then fail. Taking $E'=\operatorname{im}(f)=\{0\}$ repairs the nullary case and agrees with the explicit description of $N=0$ on manuscript page~47. This is an omitted boundary convention; it does not change the nonconstant construction.

Write the distinct elements of $E'$ as
\[
 0=t_0<t_1<\cdots<t_N=m.
\]
Every two consecutive values are joined by a target arrow: the source transition between their two adjacent nonempty fibres has exactly these image endpoints. Hence the segment from $t_{j-1}$ to $t_j$ is monochromatic, with sign $\varepsilon_j$. Equivalently, $E'$ contains the intrinsic critical vertices of the target, namely its endpoints and every place where the signs change. The notation $\mathrm{Cr}(\tau')$ has changed meaning at this point in the manuscript: here it denotes these vertices, whereas earlier it denoted critical intervals. Keeping these two sets distinct removes a potential confusion in the formula.

The fibre over $t_j$ is a nonempty consecutive interval of vertices; let its signature be $\tau_j$ and its length $\lambda_j$. Reading the source from left to right gives exactly
\[
 \tau=\tau_0\Delta_{\varepsilon_1}\tau_1
        \Delta_{\varepsilon_2}\cdots
        \Delta_{\varepsilon_N}\tau_N.
\]
Conversely, choose any subset $E'$ containing the critical vertices and any words $\tau_0,\ldots,\tau_N$. Form this source word. Collapse the $\tau_j$-block to $t_j$, and map the transition between consecutive blocks to the target arrow joining $t_{j-1}$ and $t_j$, in the direction prescribed by $\varepsilon_j$. This defines an increasing endpoint-preserving functor; its image is exactly $E'$, and its fibres have the prescribed words. Each arrow is uniquely determined by its endpoints in an interval category, so there is no additional gluing datum. The two constructions are inverse. For a point target they reduce to the arbitrary word $\tau_0$ and its unique map to the point.

For use in computations, set
\[
 l_j=t_j-t_{j-1},\qquad
 B_j=j+\sum_{h<j}\lambda_h.
\]
The correct full formula is
\[
 f(k)=t_j=\sum_{h=1}^{j}l_h
 \quad\text{for } B_j\leq k\leq B_j+\lambda_j,
 \qquad 0\leq j\leq N.
\]
The source's explicit table gives $l_2$ on the third fibre, where the value must be $l_1+l_2$; its final row similarly needs the corresponding cumulative sum. The formula above repairs all rows at once. For example, take target $++-$, visited vertices $E'=\{0,2,3\}$, and fibre words $\tau_0=-$, $\tau_1=+$, $\tau_2$ empty. Then the source is $-++-$ and the object map is $(0,0,2,2,3)$. The first positive transition crosses the two target positive edges, the next is collapsed, and the final negative transition crosses the final target edge. The two jumps have sizes $2$ and $1$, while the final vertex value is their sum $3$.
\end{dcomment}


\dgrblock{dgrb309}{Which contractions admit a section}
\begin{dcomment}{development}{p47}{Which contractions admit a section}
The normal form makes it possible to ask precisely when a deleted collection of transitions can be restored functorially. This is useful for the later path formalism: a surjection of underlying vertex sets need not be a split epimorphism in the category of signed types. The signs can obstruct every possible section.

\textbf{Section criterion for a contraction.}\ours{}
Let $\tau=(\varepsilon_1,\ldots,\varepsilon_n)$ and let $E=\{e_1<\cdots<e_k\}$ be a subset of its transition positions. Let
\[
 p_E:\tau\longrightarrow\tau|_E
\]
be the contraction with weights $1$ at the positions in $E$ and $0$ elsewhere. If $k=0$, it has a section exactly when $n=0$. If $k>0$, it has a section exactly under the following conditions. Every sign before $e_1$ equals $\varepsilon_{e_1}$, and every sign after $e_k$ equals $\varepsilon_{e_k}$. In each open gap $(e_j,e_{j+1})$, the sequence of signs is a possibly empty string of $\varepsilon_{e_j}$ followed by a possibly empty string of $\varepsilon_{e_{j+1}}$. When these conditions hold, the number of sections is
\[
 \prod_{\substack{1\leq j<k\\
                  \varepsilon_{e_j}=\varepsilon_{e_{j+1}}}}
       (e_{j+1}-e_j).
\]
An empty product is $1$.

\emph{Proof.} Under the signed-edge anti-isomorphism, $p_E$ corresponds to the inclusion $E\hookrightarrow[1,n]$. A section of $p_E$ therefore corresponds exactly to a nondecreasing sign-preserving retraction
\[
 r:[1,n]\longrightarrow E,\qquad r(e_j)=e_j.
\]
If $E$ is empty such a map exists exactly when its domain is empty. Assume henceforth $k>0$. Monotonicity forces every index before $e_1$ to map to $e_1$, and every index after $e_k$ to map to $e_k$. This proves the two exterior conditions, and these parts of $r$ are unique.

For $e_j<i<e_{j+1}$, monotonicity forces
$e_j\leq r(i)\leq e_{j+1}$. Since there is no element of $E$ strictly between them, $r(i)$ is one of these two endpoints. Moreover the indices mapped left must all precede the indices mapped right. There is therefore a single cut in the gap. Sign preservation is equivalent to the signs before the cut being $\varepsilon_{e_j}$ and those after it being $\varepsilon_{e_{j+1}}$. This proves necessity of the gap condition. If the endpoint signs differ, an admissible cut is unique. If they agree, every sign in the gap must have that common value and all $d+1$ cuts are admissible, where $d=e_{j+1}-e_j-1$.

Conversely, choose an admissible cut independently in every gap, use the forced exterior values, and fix every $e_j$. The resulting map is nondecreasing across each gap and its boundary, hence globally nondecreasing; it preserves signs by construction and retracts onto $E$. These choices describe every retraction exactly once. Multiplying their independent numbers gives the displayed product. The signed-edge correspondence turns each retraction back into a unique section of $p_E$, completing the proof.

For $\tau=+-$ and $E=\{1\}$ the trailing negative sign violates the exterior condition. Thus the surjective vertex map $+ -\to+$ has no section. For $\tau=+++$ and $E=\{1,3\}$ the gap has one omitted positive edge, and there are exactly two sections; their vertex maps $[0,2]\to[0,3]$ are $(0,1,3)$ and $(0,2,3)$. This distinguishes the collapse of an edge from a simplicial degeneracy with an available section.

We checked the ordinal description and comonoid interpretation in \citet[Appendix~D]{AguiarMahajan2010}, the zig-zag and Reedy structures of \citet[Definition~1.1 and Example~4.12]{ArakawaCnossen2026}, and searched for signed or coloured ordinal retractions and sections of Grothendieck's interval maps. We could not find this exact interval-form criterion and count in those sources. The result is an elementary consequence of the normal form, recorded here as a proved continuation of the calculation, without a claim of priority.
\end{dcomment}


\section*{6\quad Composition dans $\uTy$}

On a une loi de composition (partout définie) dans $\mathrm{Ob}\,(\uTy) = \Ty$, c'est celle du monoïde libre à deux générateurs $\Delta_+$, $\Delta_-$. Cette loi de composition s'interprète, au niveau catégorique (et non plus seulement combinatoire) comme une somme amalgamée
\[
I_{\tau\tau'} = I_\tau \amalg_{(e, \underline{1}_\tau, \underline{0}_{\tau'})} I_{\tau'},
\]
où $e$ est la catégorie unité, et où $\underline{1}_\tau$, $\underline{0}_{\tau'}$ représentent maintenant les morphismes $e \longrightarrow I_\tau$, $e \longrightarrow I_{\tau'}$ , ayant comme image l'élément but et l'élément source respectivement de l'intervalle orienté envisagé. On a donc un carré cocartésien dans (Cat)
\[
\begin{tikzcd}[column sep=small, row sep=small]
 & e \arrow[dl, hook', "\underline{1}_\tau"'] \arrow[dr, hook, "\underline{0}_{\tau'}"] & \\
I_\tau \arrow[dr, hook, "i_{\tau,\tau'}"'] & \text{\scriptsize cocart.} & I_{\tau'} \arrow[dl, hook', "j_{\tau,\tau'}"] \\
 & I_{\tau\tau'}, &
\end{tikzcd}
\]
où toutes les flèches sont des foncteurs mono et pleinement fidèles. Que ce diagramme (évidemment commutatif) soit bien cocartésien dans (Cat) revient, par définition, à dire que pour toute $X$ dans Cat, le diagramme transposé
\[
\begin{tikzcd}[column sep=small, row sep=small]
 & X & \\
\uCh_\tau(X) \arrow[ur, "\beta_\tau" near end, "\text{\scriptsize (but)}" near start] & & \uCh_{\tau'}(X) \arrow[ul, "\sigma_{\tau'}"' near end, "\text{\scriptsize (source)}"' near start] \\
 & \uCh_{\tau\tau'}(X), \arrow[ul] \arrow[ur] &
\end{tikzcd}
\]

\texttt{[est cartésien]}

\msp{page 49}ce qui revient à la tautologie suivante qu'un chemin $C$ dans $X$ de type $\tau\tau'$ `n'est autre chose' qu'un couple $(c,c')$ d'un chemin $c$ de type $\tau$ et d'un chemin $c'$ de type $\tau'$ `mis bout à bout' (dans cet ordre), i.e. tels que
\[
\mathrm{but}(c) = \mathrm{source}(c')
\]
(et on a alors
\[
\begin{array}{rcll}
\mathrm{source}(C) & = & \mathrm{source}(c) & \\[1ex]
\mathrm{but}(C) & = & \mathrm{but}(c') & \qquad\quad ).
\end{array}
\]

\printedpage{33}
% head: A. Grothendieck : Dérivateurs, Ch. VII, Édition des Universités de Montpellier II et Paris VII
Il se pose ici le sempiternel problème si on va noter ce chemin $c{\cdot}c'$, ce qui semble à première vue s'imposer, ou plutôt $c' \circ c$, regardant le chemin $c$ de $x = \mathrm{source}(c)$ vers $y = \mathrm{but}(c)$ comme une sorte de morphisme de $x$ dans $y$ (cf. plus bas le formalisme des catégories de fractions), $c'$ comme un morphisme de $y = \mathrm{source}(c')$ vers $z = \mathrm{but}(c')$, et en interprétant le composé $C$ comme le composé de ces deux morphismes $c$, $c'$ (dans cet ordre du point de vue des opérations à effectuer), qu'on note traditionellement $c' \circ c$ ou $c' \cdot c$, et non $c \circ c'$ ou $c \cdot c'$, hélas! C'est cette notation que nous allons suivre plus bas. % [sic] traditionellement -> traditionnellement, orthographe

Mais pour l'instant on ne va pas se soucier des chemins dans des catégories quelconques $X$, mais revenir au formalisme des $\tau$ et des $I_\tau$ eux-mêmes. L'interprétation catégorique de $\tau\tau'$ comme une somme amalgamée nous assure que cette opération de composition est fonctorielle :

\msp{page 50}
\[
\tau_1 \xrightarrow{\ f\ } \tau_2, \quad \tau'_1 \xrightarrow{\ f'\ } \tau'_2 \qquad\quad \text{donne } \tau_1\tau'_1 \xrightarrow{f*f'} \tau_2\tau'_2
\]
(on n'écrira pas $f \cdot f'$ pour la flèche notée $f * f'$, à cause de la confusion possible avec la composition des flèches dans $\uTy$). On a transitivité de façon évidente : si de plus $\tau_2 \xrightarrow{\ g\ } \tau_3$, $\tau'_2 \xrightarrow{\ g'\ } \tau'_3$, \texttt{[on a]}
\[
(gf) * (g'f') = (g * g') \circ (f * f').
\]
Donc on a un bifoncteur de composition
\[
\uTy \times \uTy \longrightarrow \uTy \qquad\quad (c,c') \longmapsto c \cdot c' \text{ ou } c * c'.
\]
(En fait, on a aussi un bifoncteur
\[
\uTyp \times \uTyp \longrightarrow \uTyp,
\]
mais on préfère oublier maintenant la catégorie \underline{Typ}, qui donne un formalisme nettement moins simple.)

\dgrblock{dgrb310}{Concatenation is a pushout with chosen endpoints}
\begin{dcomment}{context}{p49}{Concatenation is a pushout with chosen endpoints}
Grothendieck's reason for introducing types now becomes visible: the concatenation of words should encode the actual gluing of paths, including their transformations. An interval is a free category on an oriented linear graph. Gluing the last vertex of the first graph to the first vertex of the second produces the graph of the concatenated word, and the free-category construction preserves this pushout. Therefore
\[
 I_{\tau\tau'}=I_\tau\amalg_e I_{\tau'}
\]
is a pushout in $\Cat$ with the specified endpoint maps. Any directed path between two vertices of one summand stays in that summand: entering the other would require crossing the gluing vertex twice in an underlying linear graph. Thus the inclusions are fully faithful as well as injective on objects.

The corresponding assertion about path categories is also literal. Restriction gives an isomorphism
\[
 \operatorname{Fun}(I_{\tau\tau'},X)
 \cong
 \operatorname{Fun}(I_\tau,X)
 \times_X\operatorname{Fun}(I_{\tau'},X),
\]
where the fibre product uses evaluation at the two endpoints being identified. On objects, two functors agree at the common vertex exactly when their two paths concatenate. On morphisms, two natural transformations glue exactly when their components at that vertex are equal. Naturality along each generating edge then gives naturality on every composite. This verifies the categorical pullback, rather than only the corresponding assertion about sets of paths.

The convention on the order of composition now has a concrete meaning. The word $\tau\tau'$ is read first through $\tau$ and then through $\tau'$, whereas the corresponding categorical composition is written $c'\circ c$. This harmless reversal of written order should be fixed before discussing fractions. It is not an opposite-category operation and does not reverse any transition.

The endpoint data are indispensable. The parenthetical assertion on manuscript page~50 that concatenation also gives a bifunctor on $\uTyp$, as defined on page~30 using \emph{all} functors, is false. The one-arrow categories $I_+$ and $I_-$ are isomorphic in $\uTyp$. A bifunctor whose object operation is word concatenation would carry this isomorphism, together with $\mathrm{id}_{I_+}$, to an isomorphism between $I_{++}$ and $I_{-+}$. The former is a three-object chain; the latter has one initial object and two incomparable maximal objects. They are not isomorphic. Thus the claimed bifunctor cannot exist on that larger category. It does exist on endpoint-preserving functors, even without monotonicity, because their restrictions agree at the glued vertex. The final category $\uTy$ satisfies this requirement and remains closed under concatenation. This correction isolates precisely the structure that the path formalism needs.
\end{dcomment}


Cette composition est \textit{strictement} associative et \textit{strictement unitaire} (l'unité étant bien sûr l'élément unité 1 du monoïde libre Ty, correspondant à la catégorie ponctuelle $\{0\}$), et pas seulement `à isomorphisme près'. Et pour cause, puisque dans $\uTy$ les isomorphismes sont des identités. \textit{Donc} $\uTy$ \textit{est un monoïde dans la catégorie} (Cat) des petites catégories!

\msp{page 49$'$ (il y a un décalage dans la numérotation)}Cela implique que non seulement $\mathrm{Ob}\,\uTy = \Ty$, mais aussi $\mathrm{Fl}(\uTy)$, et plus généralement tous les ensembles de la forme $\mathrm{Hom}_{\mathrm{Cat}}(X, \uTy)$ sont des monoïdes, et ceci fonctoriellement en $X$, par exemple les composantes $\mathrm{Hom}_{\mathrm{Cat}}(\boldsymbol{\Delta}_i, \uTy)$ de l'ensemble semi-simplicial `nerf' de $\uTy$,
\[
\mathrm{Nerf}_i(\uTy) = \mathrm{Hom}_{\mathrm{Cat}}(\boldsymbol{\Delta}_i, \uTy),
\]
sont des monoïdes, et les applications semi-simpliciales entre elles sont des homomorphismes de monoïdes : le nerf de $\uTy$ est un \textit{monoïde semi-simplicial.}

% ---- from Der-ChVII-456_p036-040.tex
\printedpage{34}
% head: A. Grothendieck : Dérivateurs, Ch. VII, Édition des Universités de Montpellier II et Paris VII
On va décrire ces monoïdes $\mathrm{Nerf}_i(\uTy)$, et constater que ce sont tous des monoïdes \textit{libres,} tout comme $\mathrm{Nerf}_0(\uTy) = \mathrm{Ob}\,\uTy = \mathrm{Ty}$. Les applications semi-simpliciales entre ces monoïdes sont donc connus quand on les connaît sur les \textit{générateurs} de ces monoïdes libres (lesquels générateurs sont aussi les éléments indécomposables, donc définis intrinsèquement en termes de la structure multiplicative). % [sic] connus -> connues, agreement with « applications »

Pour $\mathrm{Nerf}_0(\uTy) = \mathrm{Ty}$, c'est une chose vue, les générateurs sont
\[
\Delta_+,\Delta_-,
\]
correspondant respectivement aux intervalles orientés
\[
\begin{array}{rccc}
\boldsymbol{\Delta}_+ : & 0 & \longrightarrow & 1\\[4pt]
\boldsymbol{\Delta}_- : & 0 & \longleftarrow & 1.
\end{array}
\]

\msp{page 50$'$}
Pour $\mathrm{Nerf}_1(\uTy) = \mathrm{Fl}(\uTy)$, i.e. pour les flèches, le résultat $3^o$) du paragraphe précédent (détermination de $\mathrm{Fl}(\uTy)$) s'interprète aussi en disant que $\mathrm{Fl}(\uTy)$ est le monoïde libre engendré par un système de générateurs infini
\[
\mathfrak{f}_{\varepsilon,n}\ , \qquad\quad (\varepsilon,n)\in\{+1,-1\}\times\mathbf{N},
\]
indexé par l'ensemble produit $\{+1,-1\}\times\mathbf{N}$. Ici $\mathfrak{f}_{\varepsilon,n}$ désigne l'unique flèche dans $\uTy$
\[
\mathfrak{f}_{\varepsilon,l} : \Delta_\varepsilon \longrightarrow \Delta^l_\varepsilon
\]
(où bien sûr l'exposant $l$ désigne une exponentation pour la loi de composition dans $\uTy$
\[
\Delta^l_\varepsilon = \underbrace{\Delta_\varepsilon\circ\Delta_\varepsilon\circ\ \cdots\ \circ\Delta_\varepsilon}_{l\text{ copies}},
\]
et non la puissance cartésienne dans Cat, où dans $\uTy$ (où elle n'existe que si $l\in\{0,1\}$)). % [sic] où dans -> ou dans, conjunction intended
Dans la description de $\mathrm{Fl}(\uTy)$ dans le paragraphe précédent, on a écrit le `terme général' du monoïde libre en groupant ensemble les groupes de foncteurs successifs de la forme $f_{+,0}$, $f_{-,0}$, i.e. corréspondant aux deux flèches de dégénérescence indécomposables (pour $*$) % [sic] corréspondant -> correspondant, accent
\[
\left\{
\begin{array}{lccc}
f_{+,0} : & \Delta_+ & \longrightarrow & e\\[4pt]
f_{-,0} : & \Delta_- & \longrightarrow & e
\end{array}
\right.
\]
(\footnote[15]{\textbf{NB} Le `mot' vide de $\mathrm{Fl}(\uTy)$ correspond à la flèche $\mathrm{id}_1$, où $1$ est l'élément unité de Ty, réalisant l'`intervalle' ponctuel.}). Le sous-monoïde libre de $\mathrm{Fl}(\uTy)$ engendré par ces deux éléments est le monoïde des flèches de dégénérescence (pure) quelconques
% foot: Version 22 décembre 2001
\printedpage{35}
% head: A. Grothendieck : Dérivateurs, Ch. VII, Édition des Universités de Montpellier II et Paris VII
\msp{page 51}
dans $\uTy$ :
\[
f_\tau : I_\tau \longrightarrow e \qquad\quad (\tau\in\mathrm{Ty}),
\]
et on a
\[
f_\tau * f_{\tau'} = f_{\tau\tau'},
\]
i.e. l'application
\[
\tau \longmapsto f_\tau :\quad \underbrace{\mathrm{Ob}\,(\uTy)}_{=\,\mathrm{Ty}\,=\,\mathrm{Nerf}_0(\uTy)} \quad\longrightarrow\quad \underbrace{\mathrm{Fl}(\uTy)}_{\mathrm{Nerf}_1(\uTy)}
\]
est un homomorphisme de monoïdes, transformant $\Delta_\varepsilon$ ($\varepsilon\in\{+1,-1\}$) dans $\mathfrak{f}_{\varepsilon,0}$. Ce n'est pas l'application de dégénérenscence de $\mathrm{Nerf}_0$ dans $\mathrm{Nerf}_1$, i.e. l'application $\tau\longmapsto\mathrm{id}_\tau$. Celle-ci est donnée par % [sic] dégénérenscence -> dégénérescence, spelling
\[
\begin{array}{ccc}
\mathrm{Ob}\,(\uTy) & \longrightarrow & \mathrm{Fl}(\uTy)\\[4pt]
\tau & \longmapsto & \mathrm{id}_\tau
\end{array}
\ :\
\left|\
\begin{array}{rcccl}
 & \Delta_\varepsilon & \longmapsto & \mathfrak{f}_{\varepsilon,1} & \qquad(\varepsilon\in\{+1,-1\})\\[4pt]
\text{i.e.} & \Delta_+ & \longmapsto & \mathfrak{f}_{+,1} & \\[4pt]
 & \Delta_- & \longmapsto & \mathfrak{f}_{-,1}\ . &
\end{array}
\right.
\]
Quant aux homomorphismes source et but
\[
\begin{tikzcd}[column sep=large]
\mathrm{Fl}(\uTy) \arrow[r, shift left, "\sigma\text{ (source)}"] \arrow[r, shift right, "\beta\text{ (but)}"'] & \mathrm{Ob}\,(\uTy),
\end{tikzcd}
\]
ils sont évidents sur les générateurs
\[
\left\{
\begin{array}{lcl}
\sigma(\mathfrak{f}_{\varepsilon,l}) & = & \Delta_\varepsilon\\[4pt]
\beta(\mathfrak{f}_{\varepsilon,l}) & = & \Delta^l_\varepsilon.
\end{array}
\right.
\]

\dgrblock{dgrb311}{The type category is generated by two comonoids}
\begin{dcomment}{later}{p50}{The type category is generated by two comonoids}
The strict monoidal structure places Grothendieck's calculation beside a classical universal property of the simplicial category. In the single-sign case, endpoint-preserving maps of finite intervals are opposite to monotone maps of finite ordered edge sets. The latter category, with ordinal sum, contains the universal monoid; its opposite contains the universal counital coassociative comonoid. An explicit account, with the weak-composition coordinates used here, is \citet[Appendix~D.1--D.2, especially Proposition~D.2, pp.~713--716]{AguiarMahajan2010}, which traces the universal property to Mac Lane's Proposition~VII.5.1. The two signs simply keep two independent copies of this structure in the same ordered word.

In $\uTy$ put
\[
 \delta_\varepsilon=\mathfrak f_{\varepsilon,2}:
      \Delta_\varepsilon\longrightarrow\Delta_\varepsilon^2,
 \qquad
 e_\varepsilon=\mathfrak f_{\varepsilon,0}:
      \Delta_\varepsilon\longrightarrow1,
 \qquad \varepsilon\in\{+,-\}.
\]
Both ways of applying $\delta_\varepsilon$ twice give the unique endpoint-preserving map from $\Delta_\varepsilon$ to $\Delta_\varepsilon^3$. Applying $e_\varepsilon$ to either output of $\delta_\varepsilon$ gives the identity, since its surviving weight is $1$. Thus each $\Delta_\varepsilon$ is a counital coassociative comonoid. There is no specified relation between the positive and negative comonoids.

More precisely, a strict monoidal functor from $\uTy$ to a strict monoidal category $(\mathcal C,\otimes,\mathbf1)$ is uniquely specified by a pair of counital coassociative comonoids $C_+$ and $C_-$ in $\mathcal C$. Here is the construction, which also proves the assertion. For either comonoid, let $\delta^{(a)}:C\to C^{\otimes a}$ denote the iterated comultiplication, with $\delta^{(0)}$ its counit and $\delta^{(1)}=\mathrm{id}$. Coassociativity and the counit laws give, for every list $b_1,\ldots,b_a$ of nonnegative integers,
\[
 \bigl(\delta^{(b_1)}\otimes\cdots\otimes\delta^{(b_a)}\bigr)
       \delta^{(a)}
 =\delta^{(b_1+\cdots+b_a)}.
\]
For completeness, the binary identity
$(\delta^{(p)}\otimes\delta^{(q)})\delta^{(2)}
 =\delta^{(p+q)}$ follows by induction on positive $p,q$ using coassociativity; the cases $p=0$ or $q=0$ are the counit identities, including $p=q=0$. Applying this binary identity successively to the first $a-1$ factors and the last factor proves the displayed formula by induction on $a$. The case $a=0$ says that the identity of the unit composed with the counit is the counit.

Send the word $\varepsilon_1\cdots\varepsilon_n$ to
$C_{\varepsilon_1}\otimes\cdots\otimes C_{\varepsilon_n}$ and send its outgoing morphism of weights $(a_1,\ldots,a_n)$ to
\[
 \delta_{\varepsilon_1}^{(a_1)}\otimes\cdots\otimes
 \delta_{\varepsilon_n}^{(a_n)}.
\]
The block-sum formula for composition and the preceding identity prove that this assignment is a functor. The all-one vector gives identities; the empty word gives the unit; concatenation of vectors gives the tensor product of maps. It is therefore strict monoidal. Conversely every outgoing morphism is the tensor product of its one-edge factors, and each such factor is an iterated comultiplication or counit. A strict monoidal functor must have the displayed values, establishing uniqueness.

This is a reinterpretation of the known ordinal/comonoid universal property, not a new homotopy theorem. It explains the two apparently different generators in the manuscript: $\mathfrak f_{\varepsilon,0}$ erases a transition, whereas $\mathfrak f_{\varepsilon,2}$ subdivides one. Their iterates produce every allowed reparametrization. No braiding is present: the order of the source transitions is part of the structure and cannot be permuted.
\end{dcomment}


Je vais décrire maintenant $\mathrm{Nerf}_2(\uTy)$, i.e. les systèmes
\[
\tau \xrightarrow{\ f\ } \tau' \xrightarrow{\ g\ } \tau''.
\]

\msp{page 52}
Se donner $f\in\mathrm{Fl}(\uTy)$ revient à se donner une suite ordonnée de générateurs $f_{\varepsilon,l}$, donc une suite
\[
\begin{aligned}
&\mathfrak{f}_{\varepsilon_1,l_1},\ \mathfrak{f}_{\varepsilon_2,l_2},\ \ \dots\ \ ,\mathfrak{f}_{\varepsilon_n,l_n}\\
&((\varepsilon_i,l_i))_{1\leq i\leq n}\in[\{\pm1\}\times\mathbf{N}]^n.
\end{aligned}
\]
On a alors
\[
\begin{aligned}
x &= \Delta_{\varepsilon_1}\Delta_{\varepsilon_2}\cdots\Delta_{\varepsilon_n} &\qquad\qquad& \mathrm{long}\,x = n\\
y &= \Delta^{l_1}_{\varepsilon_1}\Delta^{l_2}_{\varepsilon_2}\cdots\Delta^{l_n}_{\varepsilon_n} && \mathrm{long}\,y = \textstyle\sum_{i\in[1,n]} l_i.
\end{aligned}
\]
% foot: Version 22 décembre 2001

\printedpage{36}
% head: A. Grothendieck : Dérivateurs, Ch. VII, Édition des Universités de Montpellier II et Paris VII
Se donner $q : y\longrightarrow z$, i.e. un objet de $y\backslash\uTy$, revient à se donner une application de l'ensemble des $m=\sum l_i$ flèches de transition de $y$ dans $\mathbf{N}$. En fait, on a % [sic] q -> g, the map is called g below (« la donnée de g »)
\[
\mathrm{Tr}(y) = \mathrm{Tr}(\Delta^{l_1}_{\varepsilon_1})\amalg\mathrm{Tr}(\Delta^{l_2}_{\varepsilon_2})\amalg\ \cdots\ \amalg\mathrm{Tr}(\Delta^{l_n}_{\varepsilon_n})
\]
(`ensemble somme', i.e. réunion disjointe), donc
\[
\mathrm{Hom}_{\mathrm{Ens}}(\mathrm{Tr}(y)) \simeq \prod_{i\in[1,n]}\mathrm{Hom}_{\mathrm{Ens}}(\underbrace{\mathrm{Tr}(\Delta^{l_i}_{\varepsilon_i})}_{[1,l_i]},\mathbf{N}), % [sic] Hom_Ens(Tr(y)) -> Hom_Ens(Tr(y),N), second argument missing
\]
donc la donnée de $g$ équivaut à la donnée d'un système de $n$ applications
\[
\left\{
\begin{array}{l}
\pi_i : [1,l_i]\longrightarrow\mathbf{N} \qquad\quad (\text{pondération de }\Delta^{l_i}_{\varepsilon_i})\\[6pt]
\text{pour } i\in[1,n].
\end{array}
\right.
\]
Mais pour $i\in[1,n]$, considérons l'élément de $\mathrm{Nerf}_2$
\[
\Delta_{\varepsilon_1}\xrightarrow{\mathfrak{f}_{\varepsilon_i,l_i}}\Delta^{l_i}_{\varepsilon_i}\xrightarrow{g_{\varepsilon,\pi_i}}\underbrace{\Delta^{\pi_i(1)}_{\varepsilon_i}\Delta^{\pi_i(2)}_{\varepsilon_i}\cdots\Delta^{\pi_i(l_i)}_{\varepsilon_i}}_{\Delta_{\varepsilon_i}^{\Sigma_j\pi_i(j)\overset{\text{déf}}{=}\mathrm{long}(\pi_i)}}, % [sic] \Delta_{\varepsilon_1} -> \Delta_{\varepsilon_i}, and g_{\varepsilon,\pi_i} -> g_{\varepsilon_i,\pi_i}, index i intended
\tag*{(1)}
\]

\msp{page 53}
où $g_{\varepsilon_i,\pi_i}$ est la flèche définie en termes de $\pi_i$ au paragraphe précédent. Mais ce n'est autre aussi que
\[
g_{\varepsilon_i,\pi_i} = \mathfrak{f}_{\varepsilon_i,\pi_i(1)}*\mathfrak{f}_{\varepsilon_i,\pi_i(2)}*\cdots*\mathfrak{f}_{\varepsilon_i,\pi_i(l_i)}.
\tag*{(2)}
\]
L'élément de $\mathrm{Nerf}_2$ envisagé ne dépend que de
\[
(\varepsilon_i,l_i,\pi_i), \qquad\quad \text{où } (\varepsilon_i,l_i)\in\{\pm1\}\times\mathbf{N},\ \pi_i : [1,l_i]\longrightarrow\mathbf{N}.
\]
(\textbf{NB} Si $l_i=0$, $[1,l_i]=[1,0]=\emptyset$ par convention, et $\pi_i$ est l'unique application $\emptyset\longrightarrow\mathbf{N}$, donc la donnée $(\varepsilon_i,l_i,\pi_i)$ revient alors à $(\varepsilon_i,l_i)$.)

Si
\[
(\varepsilon,l,\pi)\ \text{ comme dessus,} \qquad\quad (\varepsilon,l)\in\{\pm1\}\times\mathbf{N},\ \pi : [1,l]\longrightarrow\mathbf{N},
\]
désignons par
\[
\mathfrak{f}_{\varepsilon,l,\pi}
\]
l'élément de $\mathrm{Nerf}_2$ donné par
\[
\Delta_\varepsilon\xrightarrow{\mathfrak{f}_{\varepsilon,l}}\Delta^l_\varepsilon\ \xrightarrow{\quad\mathfrak{f}_{\varepsilon,\pi(1)}\ \ *\ \cdots\ *\ \ \mathfrak{f}_{\varepsilon,\pi(l)}\quad}\Delta^{\pi(1)}_\varepsilon\Delta^{\pi(2)}_\varepsilon\cdots\Delta^{\pi(l)}_\varepsilon.
\]
On peut formuler la construction précédente en disant que tout élément de $\mathrm{Nerf}_2$ peut se décomposer de façon unique en un produit
\[
(f,g) = \mathfrak{f}_{\varepsilon_1,l_1,\pi_1}*\mathfrak{f}_{\varepsilon_2,l_2,\pi_2}*\cdots*\mathfrak{f}_{\varepsilon_n,l_n\pi_n}, % [sic] l_n\pi_n -> l_n,\pi_n, comma missing in index
\]
ce qui prouve que $\mathrm{Nerf}_2(\uTy)$ est le monoïde libre sur les générateurs précédents $\mathfrak{f}_{\varepsilon,l,\pi}$. (\textbf{NB} à la suite vide correspond
\[
1\xrightarrow{\mathrm{id}_1}1\xrightarrow{\mathrm{id}_1}1 \qquad\qquad .)
\]
% foot: Version 22 décembre 2001

\printedpage{37}
% head: A. Grothendieck : Dérivateurs, Ch. VII, Édition des Universités de Montpellier II et Paris VII
\msp{page 54}
Si l'objet de $\mathrm{Nerf}_2$ est
\[
x\xrightarrow{\ f\ }y\xrightarrow{\ g\ }z,
\]
la construction précédente nous dit, en termes de la donnée combinatoire, i.e. de la suite de données
\[
(\varepsilon_1,l_1,\pi_1),\ \dots\ ,\ (\varepsilon_n,l_n,\pi_n),
\]
quels sont $x$, $y$, $z$ : chacun d'eux est le produit dans $\mathrm{Nerf}_0=\mathrm{Ty}$ des $n$ facteurs $x_i$, $y_i$ ou $z_i$ respectivement, pour $i\in[1,n]$, donnés dans la définition des $\mathfrak{f}_{\varepsilon_i,l_i,\pi_i}$, i.e. de
\[
\Delta_{\varepsilon_i},\ \Delta^{l_i}_{\varepsilon_i},\ \text{ et }\Delta^{\pi_i(1)}_{\varepsilon_i}\Delta^{\pi_i(2)}_{\varepsilon_i}\cdots\Delta^{\pi_i(l_i)}_{\varepsilon_i}=\Delta^{\mathrm{long}(\pi_i)}_{\varepsilon_i}
\]
(où $\mathrm{long}(\pi_i)=\sum_{j\in[1,l_i]}\pi_i(j)$). Donc
\[
\left\{
\begin{array}{lcl}
x & = & \prod^{\mathrm{Ty}}_{i\in[1,n]}\Delta_{\varepsilon_i}\\[4pt]
y & = & \prod^{\mathrm{Ty}}_{i\in[1,n]}\Delta^{l_i}_{\varepsilon_i}\\[4pt]
z & = & \prod^{\mathrm{Ty}}_{i\in[1,n]}\Delta^{\mathrm{long}(\pi_i)}_{\varepsilon_i}
\end{array}
\right.
\]
(\footnote[16]{attention, les produits doivent être pris en respectant l'ordre des facteurs donné par la notation indicielle en $i$.}), et on trouve de même $f$, $g$ comme des produits (au sens $*$) de $n$ facteurs $f_i$ resp. $g_i$ ($i\in[1,n]$) explicités dans les formules (1), (2) (p. 52, 53) :
\[
\begin{array}{lclccl}
f_i & = & \mathfrak{f}_{\varepsilon_i,l_i} : & \Delta_{\varepsilon_i} & \longrightarrow & \Delta^{l_i}_{\varepsilon_i}\\[4pt]
g_i & = & g_{\varepsilon_i,l_i} : & \Delta^{l_i}_{\varepsilon_i} & \longrightarrow & \Delta^{\mathrm{long}(\pi_i)}_{\varepsilon_i}, % [sic] g_{\varepsilon_i,l_i} -> g_{\varepsilon_i,\pi_i}, notation of (1), (2)
\end{array}
\]
$g_i$ lui-même étant représenté comme un $*$-produit de $l_i$ facteurs indécomposables, de sorte que $f$ et $g$ apparaissent sous forme
\msp{page 55}
de $*$-produits de $n$, resp. de $\sum_{i\in[1,n]}l_i$ facteurs de la forme $\mathfrak{f}_{\varepsilon,\nu}$. Pour terminer de tenir en mains l'objet semi-simplicial tronqué $(\mathrm{Nerf}_i(\uTy))_{0\leq i\leq 2}$, il reste à décrire les deux applications de dégénérescence $\mathrm{Nerf}_1\longrightarrow\mathrm{Nerf}_2$ (cas où $f$ ou $g$ est une flèche identique) et surtout l'application de \textit{composition} des flèches
\[
\mathrm{Nerf}_2(\uTy)\xrightarrow{\ \gamma\ }\mathrm{Nerf}_1(\uTy)=\mathrm{Fl}(\uTy).
\]
Il suffit de donner $\gamma$ sur les générateurs, et on a tautologiquement, par définition des $\mathfrak{f}_{\varepsilon,l,\pi}$
\[
\gamma(\mathfrak{f}_{\varepsilon,l,\pi})=\mathfrak{f}_{\varepsilon,\mathrm{long}(\pi)}.
\]
(\textbf{NB} Il n'existe au plus qu'un seul morphisme de $\Delta_\varepsilon$ (de $\Delta_+$ ou de $\Delta_-$) dans un intervalle $I_\tau$. Il existe si et seulement si $\tau$ est de la forme $\Delta^l_\varepsilon$, $l\in\mathbf{N}$.)
% foot: Version 22 décembre 2001

\dgrblock{dgrb312}{Free monoids in degrees one and two}
\begin{dcomment}{exposition}{p53}{Free monoids in degrees one and two}
Grothendieck uses the nerve to keep composition together with concatenation. A string of composable reparametrizations can still be cut along the original source edges. The resulting factors are independent, and this is the reason the monoid in each nerve degree is free.

In degree one, a morphism with source word $\varepsilon_1\cdots\varepsilon_n$ and weights $(a_1,\ldots,a_n)$ has the unique decomposition
\[
 f=\mathfrak f_{\varepsilon_1,a_1}*\cdots*
       \mathfrak f_{\varepsilon_n,a_n}.
\]
The factors are indexed by source edges, even when their weights vanish. A generator $\mathfrak f_{\varepsilon,0}$ is therefore not the unit of this free monoid: its source has one edge. The monoid unit is the empty factorization, namely the identity of the point. Nor is $\mathfrak f_{\varepsilon,1}$ the monoid unit; it is the identity of a one-edge object and is indecomposable under concatenation. This separates the categorical identity operation from the unit of the tensor product.

Now consider $\tau\xrightarrow f\tau'\xrightarrow g\tau''$. The $i$th original edge produces an ordered block of $a_i$ edges after $f$. Let $\pi_i:[1,a_i]\to\mathbb N$ record the weights that $g$ assigns to that block. All descendants have the original sign $\varepsilon_i$. Consequently the pair has the unique factorization
\[
 (f,g)=\prod_{i=1}^{n,*}\mathfrak f_{\varepsilon_i,a_i,\pi_i}.
\]
Existence comes from cutting every intermediate word into its descendant blocks; uniqueness comes from the fixed order of the original edges and the unique weight vector of each restriction. The nullary case $a_i=0$ has one empty function $\pi_i$, so it contributes a legitimate generator whose two targets are the point. For $n=0$ the whole pair is the identity string at the point and is the monoid unit. This proves the claimed freeness in degree two without discarding extinct blocks.

Writing $N_r=N_r(\uTy)$, the three faces of a degree-two generator are
\[
\begin{aligned}
 d_2(\mathfrak f_{\varepsilon,a,\pi})
    &=\mathfrak f_{\varepsilon,a},\\
 d_1(\mathfrak f_{\varepsilon,a,\pi})
    &=\mathfrak f_{\varepsilon,\sum_{j=1}^{a}\pi(j)},\\
 d_0(\mathfrak f_{\varepsilon,a,\pi})
    &=\mathfrak f_{\varepsilon,\pi(1)}*\cdots*
        \mathfrak f_{\varepsilon,\pi(a)}.
\end{aligned}
\]
The last line is the empty product when $a=0$. The two degeneracies of a degree-one generator insert the identity before or after its map:
\[
 s_0(\mathfrak f_{\varepsilon,a})
       =\mathfrak f_{\varepsilon,1,(a)},\qquad
 s_1(\mathfrak f_{\varepsilon,a})
       =\mathfrak f_{\varepsilon,a,1_a}.
\]
Here $(a)$ is the one-entry weight vector and $1_a$ the all-one vector of length $a$. These formulas recover precisely the composition and degeneracies described on manuscript pages~54--56, while displaying the other two faces as well.

The word ``semi-simplicial'' in this edition includes degeneracies; in current terminology the nerve is a simplicial monoid. The faces and degeneracies are monoid homomorphisms because concatenation is a bifunctor and is strictly associative and unital. This statement concerns the whole degreewise free monoid. Its distinguished basis need not itself be a simplicial set: $d_0$ can send one generator to a product of several generators, or to the unit.
\end{dcomment}


\printedpage{38}
% head: A. Grothendieck : Dérivateurs, Ch. VII, Édition des Universités de Montpellier II et Paris VII
Le cas $g=\mathrm{id}$ correspond au cas où
\[
\pi_i(j)=1 \qquad\quad \forall j\in[1,l_i],
\]
les $\mathfrak{f}_{\varepsilon_i,l_i,\pi_i}$ sont alors détérminés en fonction des $\mathfrak{f}_{\varepsilon_i,l_i}$, et l'application de dégénérescence est % [sic] détérminés -> déterminés, accent
\[
\mathfrak{f}_{\varepsilon,l}\longmapsto\mathfrak{f}_{\varepsilon,l,1_l} \qquad\quad (\text{où } 1_l \text{ est la fonction constante de valeur 1 sur } [1,l]).
\]
Le $f=\mathrm{id}$ est celui où les $l_i$ sont égaux à 1. Alors chaque $\pi_i : [1,1]\longrightarrow\mathbf{N}$ peut être vu comme un élément de $\mathbf{N}$, la donnée de $(\varepsilon_i,l_i,\pi_i)$ revient à celle des $(\varepsilon_i,n_i)\in\{\pm1\}\times\mathbf{N}$, i.e. la donnée % [sic] Le f = id -> Le cas f = id, word « cas » missing
\msp{page 56}
d'une suite de $n$ générateurs parmi les $\mathfrak{f}_{\varepsilon,l}$. L'application de dégénérescence $\mathrm{Nerf}_1\longrightarrow\mathrm{Nerf}_2$ est ici donnée sur les générateurs par
\[
\mathfrak{f}_{\varepsilon,l}\longmapsto\mathfrak{f}_{\varepsilon,1,l}.
\]
Heureusement, on n'a plus à prouver l'associativité de la composition.

\textbf{NB} La simplicité de ces descriptions renforce mon sentiment que la catégorie $\uTy$ définie à présent, avec les \textit{morphismes d'intervalles orientés croissants} $I_\tau\longrightarrow I_{\tau'}$ comme flèches de $\uTy$, est bien `la bonne'. On va déjà voir si elle donne une description sympa des catégories de fractions $X\Sigma^{-1}$, sans tortillements comme précédemment.

Mais on a promis de déterminer $\mathrm{Nerf}_r(\uTy)$, pour tout $r\in\mathbf{N}$. Le générateur général de $\mathrm{Nerf}_r$ est de la forme (pour $\varepsilon\in\{\pm1\}$ fixé)
\[
\underbrace{\Delta_\varepsilon}_{=\,\Delta^{l_0}_\varepsilon\text{avec } l_0=1}\xrightarrow{\ f_1\ }\Delta^{l_1}_\varepsilon\xrightarrow{\ f_2\ }\Delta^{l_2}_\varepsilon\longrightarrow\ \cdots\ \xrightarrow{\ f_r\ }\Delta^{l_r}_\varepsilon,
\]
où les $f_i : \Delta^{l_{i-1}}_\varepsilon\longrightarrow\Delta^{l_i}_\varepsilon$, $i\in[1,r]$, sont donnés par des applications
\[
\pi_{i-1} : [1,l_{i-1}]\longrightarrow\mathbf{N} \qquad\quad \text{telles que }\ \underbrace{\mathrm{long}(\pi_{i-1})}_{\overset{\text{déf}}{=}\ \Sigma_{j\in[1,l_i]}\pi_i(j)}\ = l_i. % [sic] \Sigma_{j\in[1,l_i]}\pi_i(j) -> \Sigma_{j\in[1,l_{i-1}]}\pi_{i-1}(j), indices shifted with respect to long(\pi_{i-1})
\]
(NB $\pi_0 : [1,1]=\{1\}\longrightarrow\mathbf{N}$ est identifié à $\mathrm{long}(\pi_0)=\pi_0(1)=l_1$.)

\msp{page 57}
Pour ne pas faire dépendre la définition du domaine des $\pi_i$ ($i\in[0,r-1]$) des entiers $l_i$, on peut aussi considérer les $\pi_i$ comme des applications
\[
\mathbf{N}^*\xrightarrow{\ \pi\ }\mathbf{N} \qquad\quad (\mathbf{N}^*=\{i\in\mathbf{N}\mid i\geq 1\})
\]
à supports finis (i.e. tels que $\pi(i)=0$ pour $i$ grand), et exiger de façon précise
\[
\begin{aligned}
&\mathrm{supp}\,\pi_0\ \subseteq\ \{1\}\\
&\mathrm{supp}\,\pi_1\ \subseteq\ [1,\mathrm{long}(\pi_0)]\\
&\qquad\ \cdots\\
&\boxed{\mathrm{supp}\,\pi_i\ \subseteq\ [1,\mathrm{long}(\pi_{i-1})]}, \qquad\quad \text{pour }\ \boxed{i\in[0,r-1]}\quad .
\end{aligned}
\tag*{$(**)$}
\]
% foot: Version 22 décembre 2001


\dgrblock{dgrb313}{All nerve degrees as ordered forests}
\begin{dcomment}{exposition}{p57}{All nerve degrees as ordered forests}
The final assertion of manuscript page~57 is that the same decomposition works for arbitrarily long strings of reparametrizations. It is useful to complete the argument now, since it explains both the support conditions at the bottom of the page and the one face map that behaves differently from the others.

Fix $r\geq0$. Call an $r$-step string \emph{rooted} when its initial type has one transition. Every intermediate type is then a power of the same sign, and the string has the form
\[
 \Delta_\varepsilon^{l_0}\longrightarrow
 \Delta_\varepsilon^{l_1}\longrightarrow\cdots\longrightarrow
 \Delta_\varepsilon^{l_r},\qquad l_0=1.
\]
For $0\leq a<r$ its $a$th weight vector is
\[
 \pi_a:[1,l_a]\longrightarrow\mathbb N,
 \qquad l_{a+1}=\sum_{j=1}^{l_a}\pi_a(j).
\]
One can regard this as an ordered rooted tree with $r$ levels of edges: a transition at level $a$ has $\pi_a(j)$ ordered descendants at level $a+1$. A vertex may have no descendants. Once a whole level is empty, all subsequent levels are empty because there is no morphism from the point to a nonpoint type.

For a general string, let its initial word have transitions $1,\ldots,n$. At the first stage the descendants of each original transition form a consecutive monochromatic block. At the next stage the descendants of all transitions in one such block again form a consecutive block: concatenate their individual image blocks, allowing zero lengths. Induction gives such a block at every level. Restricting all maps to the descendants of a fixed original edge produces a rooted string, and concatenating the $n$ rooted strings reconstructs the original string. Each boundary is determined by cumulative weights, so this factorization is unique. A rooted string cannot be a product of two nonunit strings, since the lengths of their initial words would be positive integers summing to one. Conversely a string whose initial word has at least two edges splits after its first edge. Thus rooted strings are exactly the indecomposables, and
\[
 N_r(\uTy)\cong\text{the free monoid on rooted $r$-step strings}.
\]
For $r=0$ there are exactly two generators, the two signs. For an empty initial word there is only the all-identity string at the point, which is the monoid unit. These boundary cases establish the assertion for every $r$.

The simplicial operators can now be read directly on these forests. An inner face removes one level and replaces the two successive weight vectors by the sums over consecutive descendant blocks. The last face forgets the last level. A degeneracy repeats a level, inserting weight $1$ at each of its transitions. These operations preserve a single root. The first face removes the original root level: a rooted string with $l_1$ first-generation transitions becomes the concatenation of their $l_1$ rooted descendant strings. This is the higher-dimensional version of the formula for $d_0$ above. Since these are exactly restriction and composition of actual functors $[r]\to\uTy$, they satisfy all simplicial identities; the unique forest decomposition identifies their formulas with those categorical operators.

For $r>0$, to use fixed-domain functions as on manuscript page~57, extend each $\pi_a$ by zero to $\mathbb N^*$. The exact conditions are
\[
 \operatorname{supp}(\pi_0)\subseteq\{1\},\qquad
 \operatorname{supp}(\pi_a)\subseteq
       [1,\operatorname{long}(\pi_{a-1})]
       \quad(1\leq a<r).
\]
The source's final boxed range starts at $a=0$, although $\pi_{-1}$ is undefined; it must begin at $1$. The preceding display defining $\operatorname{long}(\pi_{i-1})$ likewise needs the indices $j\in[1,l_{i-1}]$ and $\pi_{i-1}(j)$. These printed indexing slips do not affect the argument once the successive domains are kept in view.

One should not infer from this result that the simplicial monoid is freely generated by a simplicial set of indecomposables: the first-face calculation rules that out for this particular basis. Nor does the algebraic freeness make its classifying space complicated. The point is terminal in $\uTy$, since every weight vector admits the unique further collapse to zero. Appending that terminal object gives the usual contraction of the nerve. The free presentations describe how all reparametrizations fit together, while this terminal object describes their ordinary homotopy type.
\end{dcomment}


\dgrblock{dgrb314}{What the shape calculus does for localization}
\begin{dcomment}{status}{p57}{What the shape calculus does for localization}
The free monoids just constructed organize subdivision and concatenation. To obtain mapping spaces, a second ingredient is needed: comparisons of labelled paths must carry the specified weak-equivalence information. For a relative category $(C,W)$ and fixed $x,y$, let $Z_\tau(C,W)_{x,y}$ be the category of paths of type $\tau$ from $x$ to $y$, with backward arrows in $W$, and natural transformations whose components lie in $W$ and whose endpoint components are identities. Precomposition makes these categories contravariant in $\uTy$. Their nerves, or their groupoid completions, retain comparisons that the set of path components alone would discard.

The analogous construction for relative infinity-categories is developed by \citet[Definitions~3.9 and~3.17]{MazelGee2018Hammocks}. Under a homotopical three-arrow calculus, his Theorem~4.4 identifies the geometric realizations of its hammock mapping objects with the groupoid completions of three-arrow zigzag categories. This is an additional hypothesis, not a consequence of the finite normal form in the manuscript. Independently, \citet[Theorem~3.8]{MazelGee2019Rezk} proves that completing the relative Rezk nerve presents localization. That universal property supplies the criterion with which a proposed path model must be compared.

For the general mapping-space comparison, \citet[Theorem~1.9 and Corollary~1.11]{ArakawaCnossen2026} prove in their preprint that
\[
 \operatorname{Map}_{C[W^{-1}]}(x,y)
 \simeq \operatorname*{colim}_{\tau\in\uTy^{\mathrm{op}}}
          \bigl|Z_\tau(C,W)_{x,y}\bigr|,
\]
where the colimit is taken in spaces and $|{-}|$ is groupoid completion. The displayed formula uses the ordinary relative categories just defined; the passage from their strict endpoint fibres to the invariant formula is supplied by Convention~4.2, Definition~4.3, Theorem~4.7, and Proposition~4.8(b) of the same paper. The general infinity-categorical statement uses the corresponding infinity-categories of relative diagrams and their endpoint fibres. Their Corollary~5.7 supplies the localization identification for Mazel-Gee's hammock construction. These references distinguish the established Rezk-nerve theorem from the more recent explicit comparison.

The comparison makes the scope of the present pages clear. They provide the shapes, their functorial changes, and strict concatenation. They do not yet choose $W$ or impose the conditions on vertical arrows used above. No new open conjecture is required to finish their finite combinatorics. The subsequent localization construction still has to be checked with its own morphisms and universal property; the contractibility of an interval, by itself, does not perform that check.
\end{dcomment}
\begin{otherlanguage}{english}
\dgroverview{VII-a}{Chapter VII \textperiodcentered\ \textup{VII-a} \textperiodcentered\ Overview}
\selectlanguage{english}
\noindent\textbf{Chapter VII, manuscript pages 1--57: the shapes of paths.}

The opening of Chapter VII is a restart. Grothendieck has found the earlier path formalism too cumbersome to serve as a foundation for localization, and he now separates a path from the finite combinatorial object on which it is parametrized. The distinction gives these pages their unity. First one understands intervals independently of any coefficient category; next one regards every operation on paths as a change of interval; finally one assembles all types and makes their concatenation strictly functorial. The later theory of fractions is a test of the construction, not yet a theorem established in these pages. The commentary on localization and path completion in Chapter~VI, especially at manuscript pages~64--70, supplies the background; the present chapter changes the finite apparatus on which such arguments are to rest.

An interval of length $n$ is the free category on an oriented path with vertices $0,\ldots,n$. Its signature is a word $\tau=(\varepsilon_1,\ldots,\varepsilon_n)$ in two signs. There are two orders to remember. The progression order is the linear order of the vertices; the categorical order records the existence of arrows. A path proceeds from vertex~$0$ to vertex~$n$ even when some of its arrows point backwards. Reversing this progression order changes the signature to $(-\varepsilon_n,\ldots,-\varepsilon_1)$, whereas taking the opposite category changes it to $(-\varepsilon_1,\ldots,-\varepsilon_n)$. Their distinction is already needed to understand why a span is symmetric as an interval but is not isomorphic to its opposite cospan.

The maximal constant runs of the signature are the structural chains. Their endpoints are the critical vertices; retaining just these vertices gives the reduced interval. These notions organize all full subcategories of an interval. A nonempty full subcategory has interval components, and a subset with extreme vertices $i<j$ is connected precisely when it retains every critical vertex between them. The proof is elementary and decisive: two consecutive retained vertices can be joined by an arrow exactly when no change of sign has been omitted between them. Thus reduction remembers the turns while suppressing subdivisions inside a composable run. It does not mean that arbitrary diagrams on the interval may forget the intermediate objects.

Several literal statements near the beginning need repair. Proposition~1(b) asserts that every endpoint-fixing endofunctor is the identity. On $[2]$, the function $(0,1,2)\mapsto(0,0,2)$ is a counterexample. The proof on the following pages establishes rigidity of endpoint-fixing \emph{isomorphisms}, which is sufficient for the classification of orientations. In fact, the stronger endofunctor assertion holds exactly for reduced intervals. Proposition~3's deletion criterion also needs the vertices incomparable with the deleted vertex to be separated from its incoming and outgoing links. Finally, an interval-shaped subcategory need not be full. All these readings occur in the edition and remain in the source text. The accompanying blocks give counterexamples and prove the repaired assertions, so that subsequent correct arguments do not inherit an unnecessary uncertainty.

Reduction has a useful additional interpretation. For the full inclusion $j:I_{\mathrm{red}}\hookrightarrow I$, send an interior vertex of each run to that run's categorical target or categorical source, and fix all critical vertices. The resulting maps satisfy
\[
 p_{\mathrm{target}}\dashv j\dashv p_{\mathrm{source}}.
\]
The commentary proves these adjunctions by comparing the relevant hom-sets. They give categorical deformation retractions and explain exactly why reduction is harmless for groupoid-valued paths. For a general category $X$, restriction $X^I\to X^{I_{\mathrm{red}}}$ need not be an equivalence. This is the first place where the geometry of the shape and the information in a labelled diagram must be kept distinct.

Once the shape is fixed, a path of type $\tau$ in $X$ is a functor $I_\tau\to X$, and its morphisms are natural transformations. The category is $X^{I_\tau}$; its set of objects is the represented functor $\Hom_{\Cat}(I_\tau,X)$. Yoneda then says that every operation on paths natural in the coefficient category is precomposition by a unique functor between their shapes, in the reverse direction. The proof evaluates the operation on the identity path in $I_\tau$. This explains both composition of successive arrows and insertion of identities, and supplies their action on morphisms without an additional coherence choice. The manuscript's distinction between the set and the category of paths must therefore survive the typography.

The category of types changes twice while Grothendieck tests it. Initially it permits all functors between intervals; then it retains functors preserving the two endpoints; finally it also requires their vertex functions to be nondecreasing in the progression order. Only the first version is fully faithful in $\Cat$. The second is fully faithful in the category of categories under the discrete two-point category. The final version, denoted $\uTy$, is generally not full there either. These are successive definitions, not contradictory assertions about a fixed category. The final restriction preserves directed progress through a path and excludes changes of parameter which backtrack.

For any chosen version $T$, the strict contravariant functor $\tau\mapsto X^{I_\tau}$ gives a split Grothendieck fibration over $T$. A morphism from $(\tau',c')$ to $(\tau,c)$ is a pair
\[
 \varphi:\tau'\to\tau,
 \qquad v:c'\Rightarrow c\varphi.
\]
Composition is whiskering followed by composition of transformations. This is the machinery of \citet[Expos\'e~VI, \S\S8--10]{GrothendieckSGA1}, specialized to restriction along finite shapes. Endpoint-preserving shape maps make evaluation at the two ends into functors on the total category. The commentary also determines the maximal enlargement for which natural endpoint evaluations can still exist, allowing comparison arrows at the ends. That enlargement illustrates why the subsequent requirement of strict gluing is stronger. No colimits in $X$ are needed for the fibration; extending paths in the opposite direction would require Kan extensions, and already a collapsed span asks for a pushout.

The main calculation of the second half has a particularly short modern form. Give the transition set $\Tr(\tau)=\{1,\ldots,n\}$ its order and its signs. If $f:I_\tau\to I_\upsilon$ is a morphism of the final type category, every target transition belongs to the image interval of exactly one source transition. This defines a monotone, sign-preserving map
\[
 f^*:\Tr(\upsilon)\longrightarrow\Tr(\tau).
\]
Conversely, given such a map $h$, recover the vertex function by
\[
 f(i)=\#\{j:h(j)\leq i\}.
\]
These assignments are inverse and reverse composition. Thus the type category is opposite to the category of finite signed linear orders with monotone sign-preserving maps. Its many normal forms are descriptions of this one correspondence.

For a reduced target, $f^*$ is injective, so a map is specified by a signed subsequence of the source word. In general, a map out of a fixed source is specified by nonnegative integer weights: an edge of sign $\varepsilon$ and weight $a$ becomes a run of $a$ edges of sign $\varepsilon$. Weight zero collapses that edge. A map into a fixed target is instead described by its image vertices and by the signatures of the fibers over them. The latter description must use cumulative sums for the target coordinates, and collapsed fibers are components of the graph of collapsed edges, not components of the union of their endpoint vertices with every ambient adjacency restored. The point target is treated separately. These corrections preserve the two classification results.

Concatenation glues the final vertex of one interval to the initial vertex of the next. On standard words and endpoint-preserving maps it is strictly associative and unital, and makes $\uTy$ a strict monoidal category. It cannot extend with this object operation to the category allowing \emph{all} interval functors: there $+$ and $-$ are isomorphic, whereas $++$ and $-+$ are not. The restriction on morphisms is therefore essential to the source's own construction. Through the signed-edge description, $\uTy$ is the free strict monoidal category on two counital coassociative comonoids. The one-generator background is the opposite of the augmented simplex category; see \citet[Appendix~D, Proposition~D.2]{AguiarMahajan2010}. The two-generator version follows by ordered concatenation, without a symmetry interchanging the letters.

The nerve calculation is now transparent. A chain of $r$ composable type maps splits uniquely according to the edges of its initial word. Each factor begins at a single signed edge; its later stages are successive ordered expansions of that edge, allowing complete disappearance. Consequently every $N_r(\uTy)$ is a free monoid. The face deleting the initial object can send a generator to several generators, or to the empty word, so this assertion does not identify the nerve with the free simplicial monoid on its degreewise generators. The commentary gives every face and degeneracy by the corresponding operation on ordered forests and resolves the final support-indexing slip.

The later comparison can be made at the level of the indexing category itself. Marking the negative transitions identifies $\uTy$ with the zigzag type category in \citet[Definition~3.5]{MazelGee2018Hammocks}. The closing status block distinguishes this finite identification from the additional construction of localization mapping spaces, including the recent comparison in \citet[Theorem~1.9 and Corollary~1.11]{ArakawaCnossen2026}. The development blocks prove explicit consequences left unworked here: the reduction adjunctions and sharp rigidity criterion, the endpoint evaluation criterion, and a criterion with a counting formula for sections of contractions. Their markers record our proofs and limited literature searches, not a claim of priority. The unresolved task at manuscript page~57 is the promised construction of fractions with its precise morphisms; the finite combinatorics considered in these pages has been settled in the commentary.
\end{otherlanguage}
\dgrchunkend
